% Generated mechanically from frozen input inputs/p04/ja/properties.tex.
% Do not edit this generated file; edit the bound locale source instead.

% OK, start here.
%

\setcounter{chapter}{27}
\chapter{スキームの性質}



\phantomsection
\label{properties-section-phantom}


\section{はじめに}
\label{properties-section-introduction}

\noindent
本章では、スキームのいくつかの絶対的性質を導入する。
基本的な参考文献は \cite{EGA} である。




\section{構成可能集合}
\label{properties-section-constructible}

\noindent
構成可能集合および局所構成可能集合は、位相の章の節
\ref{topology-section-constructible} で導入されている。
スキームの局所構成可能部分集合は次のように特徴づけられる。

\begin{lemma}
\label{properties-lemma-locally-constructible}
$X$ をスキームとする。部分集合 $E$ が $X$ において局所構成可能であるための
必要十分条件は、$X$ において次が成り立つことである：$E \cap U$ が $U$ において
構成可能である。ただし $U$ は $X$ の任意のアフィン開集合を動く。
\end{lemma}

\begin{proof}
$E$ が局所構成可能であると仮定する。このとき、開被覆
$X = \bigcup U_i$ が存在して、$E \cap U_i$ は $U_i$ において各 $i$ ごとに
構成可能である。$V \subset X$ を任意のアフィン開集合とする。
有限アフィン開被覆 $V = V_1 \cup \ldots \cup V_m$ を、各 $j$ について
$V_j \subset U_i$ がある $i = i(j)$ に対して成り立つように取ることができる。
位相の章の補題 \ref{topology-lemma-open-immersion-constructible-inverse-image}
により、各 $E \cap V_j$ は $V_j$ において構成可能である。包含射
$V_j \to V$ は準コンパクトなので（スキームの章の補題
\ref{schemes-lemma-quasi-compact-affine} を参照）、位相の章の補題
\ref{topology-lemma-collate-constructible} により、$E \cap V$ は
$V$ において構成可能である。逆の含意は直ちに従う。
\end{proof}

\begin{lemma}
\label{properties-lemma-generic-point-in-constructible}
$X$ をスキームとし、$E \subset X$ を局所構成可能部分集合とする。
$\xi \in X$ を $X$ のある既約成分の生成点とする。
\begin{enumerate}
\item $\xi \in E$ ならば、$\xi$ のある開近傍は $E$ に含まれる。
\item $\xi \not \in E$ ならば、$\xi$ のある開近傍は $E$ と交わらない。
\end{enumerate}
\end{lemma}

\begin{proof}
局所構成可能部分集合の補集合も局所構成可能なので、(2) を示せば十分である。
$X$ はアフィンであり、したがって $E$ は構成可能であると仮定してよい
（補題 \ref{properties-lemma-locally-constructible}）。この場合 $X$ はスペクトル空間である
（代数の章の補題 \ref{algebra-lemma-spec-spectral}）。すると
$\xi \not \in E$ から $\xi \not \in \overline{E}$ が従う。これは位相の章の補題
\ref{topology-lemma-constructible-stable-specialization-closed} と、$X$ には
$\xi$ と異なり $\xi$ に特殊化する点が存在しないという事実による。
\end{proof}

\begin{lemma}
\label{properties-lemma-quasi-separated-quasi-compact-open-retrocompact}
$X$ を準分離スキームとする。$X$ の任意の二つの準コンパクト開集合の共通部分は
$X$ の準コンパクト開集合である。$X$ の任意の準コンパクト開集合は $X$ において
逆コンパクトである。
\end{lemma}

\begin{proof}
$U$ と $V$ が準コンパクト開集合ならば
$U \cap V = \Delta^{-1}(U \times V)$ であり、ここで
$\Delta : X \to X \times X$ は対角射である。$X$ は準分離なので
$\Delta$ は準コンパクトである。したがって $U \cap V$ は準コンパクトである。
実際、$U \times V$ は準コンパクトである（詳細は省略する。$U \times V$ が有限個の
アフィンの合併であることを示すには、スキームの章の補題
\ref{schemes-lemma-affine-covering-fibre-product} を用いよ）。
残りの主張は最初の主張と位相の章の補題
\ref{topology-lemma-topology-quasi-separated-scheme} から従う。
\end{proof}

\begin{lemma}
\label{properties-lemma-quasi-compact-quasi-separated-spectral}
$X$ を準コンパクトかつ準分離なスキームとする。このとき $X$ の台位相空間は
スペクトル空間である。
\end{lemma}

\begin{proof}
位相の章の定義 \ref{topology-definition-spectral-space} により、$X$ がソバーであり、
準コンパクトであり、準コンパクト開集合からなる基底を持ち、かつ任意の二つの
準コンパクト開集合の共通部分が準コンパクトであることを確認すればよい。
これはスキームの章の補題 \ref{schemes-lemma-scheme-sober} と
\ref{schemes-lemma-basis-affine-opens}、および上の補題
\ref{properties-lemma-quasi-separated-quasi-compact-open-retrocompact} から従う。
\end{proof}

\begin{lemma}
\label{properties-lemma-constructible-quasi-compact-quasi-separated}
$X$ を準コンパクトかつ準分離なスキームとする。$X$ の任意の局所構成可能部分集合は
構成可能である。
\end{lemma}

\begin{proof}
$X$ は準コンパクトなので、有限アフィン開被覆
$X = V_1 \cup \ldots \cup V_m$ を選べる。$X$ は準分離なので、補題
\ref{properties-lemma-quasi-separated-quasi-compact-open-retrocompact} により各 $V_i$ は
$X$ において逆コンパクトである。したがって位相の章の補題
\ref{topology-lemma-collate-constructible} により、$E \subset X$ が $X$ において
構成可能であるための必要十分条件は、各 $E \cap V_j$ が $V_j$ において
構成可能であることである。ゆえに補題 \ref{properties-lemma-locally-constructible} により
結論を得る。
\end{proof}

\begin{lemma}
\label{properties-lemma-retrocompact}
$X$ をスキームとする。部分集合 $E$ が $X$ の部分集合として $X$ において
逆コンパクトであるための必要十分条件は、$E \cap U$ が準コンパクトとなることである。
ただし $U$ は $X$ の任意のアフィン開集合を動く。
\end{lemma}

\begin{proof}
$X$ の任意の準コンパクト開集合が有限個のアフィン開集合の合併であることから直ちに従う。
\end{proof}

\begin{lemma}
\label{properties-lemma-stratification-locally-finite-constructible}
各部分が逆コンパクトである分割 $X = \coprod_{i \in I} X_i$（スキーム $X$ のもの）が
局所有限であるための必要十分条件は、各部分が
局所構成可能であることである。
\end{lemma}

\begin{proof}
逆コンパクト、分割、および局所有限の定義については、位相の章の定義
\ref{topology-definition-quasi-compact},
\ref{topology-definition-paritition}, および
\ref{topology-definition-locally-finite} を参照せよ。

\medskip\noindent
分割が局所有限であり、$U \subset X$ がアフィン開集合ならば、
$U = \coprod_{i \in I} U \cap X_i$ は有限分割である（より正確には、有限個を除く
すべての部分が空である）。したがって $U \cap X_i$ は準コンパクトであり、その補集合は
有限個の逆コンパクトな部分の合併として $U$ において逆コンパクトである。ゆえに
位相の章の補題 \ref{topology-lemma-locally-closed-constructible-image} により、
$U \cap X_i$ は構成可能である。補題 \ref{properties-lemma-locally-constructible} により、
$X_i$ は局所構成可能である。

\medskip\noindent
各部分が局所構成可能であると仮定する。このとき任意のアフィン開集合
$U \subset X$ に対して、構成可能部分集合による被覆
$U = \coprod X_i \cap U$ を得る。構成可能位相は準コンパクトなので（位相の章の補題
\ref{topology-lemma-constructible-hausdorff-quasi-compact} を参照）、この被覆は
有限細分を持つ。すなわち、分割は局所有限である。
\end{proof}

\section{整スキーム、既約スキーム、および被約スキーム}
\label{properties-section-integral}

\begin{definition}
\label{properties-definition-integral}
$X$ をスキームとする。$X$ が {\it 整}であるとは、それが非空であり、任意の
非空アフィン開集合 $\Spec(R) = U \subset X$ に対して環 $R$ が整域であることをいう。
\end{definition}

\begin{lemma}
\label{properties-lemma-characterize-reduced}
$X$ をスキームとする。次の条件は同値である。
\begin{enumerate}
\item スキーム $X$ は被約である。スキームの章の定義
\ref{schemes-definition-reduced} を参照せよ。
\item アフィン開被覆 $X = \bigcup U_i$ が存在し、各
$\Gamma(U_i, \mathcal{O}_X)$ は被約である。
\item 任意のアフィン開集合 $U \subset X$ に対して環
$\mathcal{O}_X(U)$ は被約である。
\item 任意の開集合 $U \subset X$ に対して環 $\mathcal{O}_X(U)$ は被約である。
\end{enumerate}
\end{lemma}

\begin{proof}
スキームの章の補題 \ref{schemes-lemma-reduced} および
\ref{schemes-lemma-affine-reduced} を参照せよ。
\end{proof}

\begin{lemma}
\label{properties-lemma-characterize-irreducible}
$X$ をスキームとする。次の条件は同値である。
\begin{enumerate}
\item スキーム $X$ は既約である。
\item アフィン開被覆 $X = \bigcup_{i \in I} U_i$ が存在し、$I$ は空でなく、
$U_i$ はすべての $i \in I$ に対して既約であり、
$U_i \cap U_j \not = \emptyset$ がすべての $i, j \in I$ に対して成り立つ。
\item スキーム $X$ は非空であり、任意の非空アフィン開集合
$U \subset X$ は既約である。
\end{enumerate}
\end{lemma}

\begin{proof}
(1) を仮定する。スキームの章の補題 \ref{schemes-lemma-scheme-sober} により、
$X$ は一意な生成点 $\eta$ を持つ。すると
$X = \overline{\{\eta\}}$ である。したがって $\eta$ は任意の非空アフィン開集合
$U \subset X$ の元である。これは $\eta \in U$ が稠密であり、ゆえに $U$ が
既約であることを意味する。また、任意の二つの非空アフィン開集合が交わることも
意味する。したがって (1) は (2) と (3) の両方を含意する。

\medskip\noindent
(2) を仮定する。$X = Z_1 \cup Z_2$ が二つの閉部分集合の合併であるとする。
すべての $i$ について、$U_i \subset Z_1$ または $U_i \subset Z_2$ のいずれかである。
ある $i \in I$ を選び、$U_i \subset Z_1$ と仮定する（必要なら $Z_1$, $Z_2$ の番号を
入れ替える）。任意の $j \in I$ に対して、開部分集合 $U_i \cap U_j$ は $U_j$ において
稠密であり、閉部分集合 $Z_1 \cap U_j$ に含まれる。したがって
$U_j \subset Z_1$ でもある。ゆえに望むとおり $X = Z_1$ である。

\medskip\noindent
(3) を仮定する。アフィン開被覆 $X = \bigcup_{i \in I} U_i$ を選ぶ。
各 $U_i$ は非空であると仮定してよい。$X$ は非空なので $I$ は空でない。
仮定により各 $U_i$ は既約である。$U_i \cap U_j = \emptyset$ となる組
$i, j \in I$ があるとする。このとき開集合
$U_i \amalg  U_j = U_i \cup U_j$ はアフィンである。スキームの章の補題
\ref{schemes-lemma-disjoint-union-affines} を参照せよ。したがって仮定によりこれは
既約であるが、矛盾である。(3) が (2) を含意することが分かった。補題が証明された。
\end{proof}

\begin{lemma}
\label{properties-lemma-characterize-integral}
スキーム $X$ が整であるための必要十分条件は、被約かつ既約であることである。
\end{lemma}

\begin{proof}
$X$ が既約ならば、任意のアフィン開集合 $\Spec(R) = U \subset X$ は既約である。
$X$ が被約ならば、上の補題 \ref{properties-lemma-characterize-reduced} により $R$ は被約である。
したがって $R$ は被約であり、$(0)$ は素イデアルである。すなわち $R$ は整域である。

\medskip\noindent
$X$ が整ならば、任意の非空アフィン開集合 $\Spec(R) = U \subset X$ に対して環 $R$ は
被約であり、したがって補題 \ref{properties-lemma-characterize-reduced} により $X$ は被約である。
さらに、任意の非空アフィン開集合は既約である。したがって $X$ は既約である。
補題 \ref{properties-lemma-characterize-irreducible} を参照せよ。
\end{proof}

\noindent
例の章の節 \ref{examples-section-connected-locally-integral-not-integral} では、
すべての局所環が整域であるにもかかわらず整ではない連結アフィンスキームを構成する。




\section{環の性質によって定義されるスキームの種類}
\label{properties-section-properties-rings}

\noindent
本節では、環のどのような性質を用いてスキームの局所的性質を定義できるかを調べる。

\begin{definition}
\label{properties-definition-property-local}
$P$ を環の性質とする。次の条件が成り立つとき、$P$ は {\it 局所的}であるという：
\begin{enumerate}
\item 任意の環 $R$ と任意の $f \in R$ に対して
$P(R) \Rightarrow P(R_f)$ が成り立つ。
\item 任意の環 $R$ と $f_i \in R$ であって、
$(f_1, \ldots, f_n) = R$ を満たすものに対して
$\forall i, P(R_{f_i}) \Rightarrow P(R)$ が成り立つ。
\end{enumerate}
\end{definition}

\begin{definition}
\label{properties-definition-locally-P}
$P$ を環の性質とする。$X$ をスキームとする。$X$ が {\it 局所的に $P$} であるとは、
任意の $x \in X$ に対して、アフィン開近傍 $U$（$x$ のもので、$X$ におけるもの）で、
$\mathcal{O}_X(U)$ が性質 $P$ を持つものが存在することをいう。
\end{definition}

\noindent
これは、その性質が局所的である場合にのみ有用な概念である。$P$ が局所的な性質で
あっても、別の箇所で定義を明示しない限り、この定義を自動的に用いてスキームを
「局所的に $P$」とは呼ばない。

\begin{lemma}
\label{properties-lemma-locally-P}
$X$ をスキームとする。$P$ を環の局所的性質とする。次の条件は同値である：
\begin{enumerate}
\item スキーム $X$ は局所的に $P$ である。
\item 任意のアフィン開集合 $U \subset X$ に対して性質
$P(\mathcal{O}_X(U))$ が成り立つ。
\item アフィン開被覆 $X = \bigcup U_i$ が存在し、各
$\mathcal{O}_X(U_i)$ は $P$ を満たす。
\item 開被覆 $X = \bigcup X_j$ が存在し、各開部分スキーム $X_j$ は
局所的に $P$ である。
\end{enumerate}
さらに、$X$ が局所的に $P$ ならば、すべての開部分スキームも局所的に $P$ である。
\end{lemma}

\begin{proof}
もちろん (1) $\Leftrightarrow$ (3) であり、(2) $\Rightarrow$ (1) である。
(3) $\Rightarrow$ (2) ならば補題の最後の主張が成り立ち、(4) も (1) と同値である
ことが容易に従う。したがって (3) $\Rightarrow$ (2) を示す。

\medskip\noindent
$X = \bigcup U_i$ をアフィン開被覆とし、$U_i = \Spec(R_i)$ と書く。
$P(R_i)$ を仮定する。$\Spec(R) = U \subset X$ を任意のアフィン開集合とする。
スキームの章の補題 \ref{schemes-lemma-good-subcover} により、$U = \Spec(R)$ の
標準開集合 $D(f_j)$ による標準被覆で、各環 $R_{f_j}$ が環 $R_i$ のいずれかの
主局所化となるものが存在する。定義 \ref{properties-definition-property-local} (1) により
$P(R_{f_j})$ を得る。すると定義 \ref{properties-definition-property-local} (2) により
$P(R)$ を得る。
\end{proof}

\noindent
応用例を一つ挙げる。

\begin{lemma}
\label{properties-lemma-reduced-is-locally-reduced}
$X$ をスキームとする。このとき、$X$ が被約であるための必要十分条件は、定義
\ref{properties-definition-locally-P} の意味で $X$ が「局所的に被約」であることである。
\end{lemma}

\begin{proof}
補題 \ref{properties-lemma-characterize-reduced} から明らかである。
\end{proof}

\begin{lemma}
\label{properties-lemma-properties-local}
環 $R$ の次の性質は局所的である。
\begin{enumerate}
\item （Cohen--Macaulay。）
環 $R$ は Noether 環かつ CM である。代数の章の定義
\ref{algebra-definition-ring-CM} を参照せよ。
\item （正則。）
環 $R$ は Noether 環かつ正則である。代数の章の定義
\ref{algebra-definition-regular} を参照せよ。
\item （絶対 Noether。）
環 $R$ は $Z$ 上有限型である。
\item 必要に応じてここにさらに追加する。\footnote{ただし、すでに別の箇所で
個別に扱った性質はここには掲げない。}
\end{enumerate}
\end{lemma}

\begin{proof}
省略する。
\end{proof}




\section{Noether スキーム}
\label{properties-section-noetherian}

\noindent
環 $R$ が {\it Noether} であるとは、イデアルの昇鎖条件を満たすことだった。
同値な言い方をすれば、$R$ のすべてのイデアルは有限生成である。

\begin{definition}
\label{properties-definition-noetherian}
$X$ をスキームとする。
\begin{enumerate}
\item $X$ が {\it 局所 Noether} であるとは、任意の $x \in X$ がアフィン開近傍
$\Spec(R) = U \subset X$ であって、環 $R$ が Noether であるものを持つことをいう。
\item $X$ が局所 Noether かつ準コンパクトであるとき、$X$ は
{\it Noether} であるという。
\end{enumerate}
\end{definition}

\noindent
局所 Noether スキームを特徴づける標準的な結果は次のとおりである。

\begin{lemma}
\label{properties-lemma-locally-Noetherian}
$X$ をスキームとする。次の条件は同値である：
\begin{enumerate}
\item スキーム $X$ は局所 Noether である。
\item 任意のアフィン開集合 $U \subset X$ に対して環 $\mathcal{O}_X(U)$ は
Noether である。
\item アフィン開被覆 $X = \bigcup U_i$ が存在し、各 $\mathcal{O}_X(U_i)$ は
Noether である。
\item 開被覆 $X = \bigcup X_j$ が存在し、各開部分スキーム $X_j$ は
局所 Noether である。
\end{enumerate}
さらに、$X$ が局所 Noether ならば、すべての開部分スキームは局所 Noether である。
\end{lemma}

\begin{proof}
これを示すには、Noether であることが環の局所的性質であることを示せば十分である。
補題 \ref{properties-lemma-locally-P} を参照せよ。Noether 環の任意の局所化は Noether である。
代数の章の補題 \ref{algebra-lemma-Noetherian-permanence} を参照せよ。
代数の章の補題 \ref{algebra-lemma-cover} により、定義
\ref{properties-definition-property-local} の第二の条件が分かる。
\end{proof}

\begin{lemma}
\label{properties-lemma-immersion-into-noetherian}
任意の埋め込み $Z \to X$ は、$X$ が局所 Noether ならば準コンパクトである。
\end{lemma}

\begin{proof}
閉埋め込みは明らかに準コンパクトである。準コンパクト射の合成は準コンパクトである。
位相の章の補題 \ref{topology-lemma-composition-quasi-compact} を参照せよ。
したがって、局所 Noether スキームへの開埋め込みが準コンパクトであることを示せば十分である。
スキームの章の補題 \ref{schemes-lemma-quasi-compact-affine} を用いて、$X$ がアフィンで
ある場合に帰着する。Noether 環のスペクトルの任意の開部分集合は準コンパクトである
（例えば、代数の章の補題 \ref{algebra-lemma-Noetherian-topology} と位相の章の補題
\ref{topology-lemma-Noetherian} および
\ref{topology-lemma-Noetherian-quasi-compact} を組み合わせよ）。
\end{proof}

\begin{lemma}
\label{properties-lemma-locally-Noetherian-quasi-separated}
局所 Noether スキームは準分離である。
\end{lemma}

\begin{proof}
スキームの章の補題 \ref{schemes-lemma-characterize-quasi-separated} により、二つの
アフィン開集合の共通部分 $U \cap V$（$X$ のもの）が準コンパクトであることを示せばよい。
これは、例えば開埋め込み $U \cap V \to U$ を考え、上の補題
\ref{properties-lemma-immersion-into-noetherian} を用いれば従う。（しかし実際には、これは単に
Noether 環のスペクトルの任意の開集合が準コンパクトであることによる。）
\end{proof}

\begin{lemma}
\label{properties-lemma-Noetherian-topology}
（局所）Noether スキームの台位相空間は（局所）Noether である。位相の章の定義
\ref{topology-definition-noetherian} を参照せよ。
\end{lemma}

\begin{proof}
これは、Noether スキームが Noether 環のスペクトルの有限合併であること、代数の章の補題
\ref{algebra-lemma-Noetherian-topology}、および位相の章の補題
\ref{topology-lemma-finite-union-Noetherian} による。
\end{proof}

\begin{lemma}
\label{properties-lemma-locally-closed-in-Noetherian}
（局所）Noether スキームの任意の局所閉部分スキームは（局所）Noether である。
\end{lemma}

\begin{proof}
省略する。ヒント：Noether 環の任意の商および任意の局所化は Noether である。
Noether の場合には、Noether 空間の任意の部分集合が（誘導位相に関して）Noether 空間で
あることも用いる。
\end{proof}

\begin{lemma}
\label{properties-lemma-Noetherian-irreducible-components}
Noether スキームの既約成分は有限個である。
\end{lemma}

\begin{proof}
Noether スキームの台位相空間は Noether であり（補題
\ref{properties-lemma-Noetherian-topology}）、Noether 位相空間の既約成分が有限個しかないので
結論を得る（位相の章の補題 \ref{topology-lemma-Noetherian}）。
\end{proof}

\begin{lemma}
\label{properties-lemma-morphism-Noetherian-schemes-quasi-compact}
任意のスキームの射 $f : X \to Y$ は、$X$ が Noether ならば準コンパクトである。
\end{lemma}

\begin{proof}
補題 \ref{properties-lemma-Noetherian-topology} と、Noether 位相空間の任意の部分集合が
準コンパクトであることを用いよ（位相の章の補題
\ref{topology-lemma-Noetherian} および
\ref{topology-lemma-Noetherian-quasi-compact} を参照）。
\end{proof}

\noindent
次は面白い補題である。任意の局所 Noether スキームには閉点が豊富にある
（少なくとも任意の閉部分集合に一つある）ことを述べている。

\begin{lemma}
\label{properties-lemma-locally-Noetherian-closed-point}
任意の非空局所 Noether スキームは閉点を持つ。局所 Noether スキームの任意の
非空閉部分集合は閉点を持つ。同値な言い方をすれば、局所 Noether スキームの任意の点は
閉点に特殊化する。
\end{lemma}

\begin{proof}
第二の主張は第一の主張から従う（スキームの章の補題
\ref{schemes-lemma-reduced-closed-subscheme} と補題
\ref{properties-lemma-locally-closed-in-Noetherian} を用いよ）。任意の非空アフィン開集合
$U \subset X$ を考える。$x \in U$ を閉点とする。$x$ が $X$ の閉点ならば終わりである。
そうでなければ、$X_0 \subset X$ を $\overline{\{x\}}$ 上の被約誘導閉部分スキーム構造と
する。スキームの章の補題 \ref{schemes-lemma-closed-subspace-scheme} により
$U_0 = U \cap X_0$ は $X_0$ のアフィン開集合であり、$U_0 = \{x\}$ である。
$y \in X_0$, $y \not = x$ を $x$ の特殊化とする。局所環
$R = \mathcal{O}_{X_0, y}$ を考える。補題
\ref{properties-lemma-locally-closed-in-Noetherian} により $X_0$ は Noether なので、これは
Noether 局所環である。$V \subset \Spec(R)$ を、$U_0$ の $\Spec(R)$ における逆像、
すなわち標準射 $\Spec(R) \to X_0$ による逆像とする（スキームの章の節
\ref{schemes-section-points} を参照）。構成により $V$ は一元集合であり、その一意な点は
$x$ に対応する（スキームの章の補題 \ref{schemes-lemma-specialize-points} を用いよ）。
代数の章の補題 \ref{algebra-lemma-Noetherian-local-domain-dim-2-infinite-opens}
により $\dim(R) = 1$ である。言い換えれば、$y$ は $x$ の直後の特殊化である
（位相の章の定義 \ref{topology-definition-dimension-function} を参照）。さらに言い換えれば、
$y \not = x$ であって $x \leadsto y$ を満たす任意の点は、$x$ の直後の特殊化である。
明らかにこれらの各点は閉点であり、これが望む結論である。
\end{proof}




\begin{lemma}
\label{properties-lemma-locally-Noetherian-specialization-dvr}
$X$ を局所 Noether スキームとする。$x' \leadsto x$ を $X$ の点の特殊化とする。このとき
\begin{enumerate}
\item 離散付値環 $R$ と射 $f : \Spec(R) \to X$ が存在し、生成点 $\eta$（$\Spec(R)$
のもの）は $x'$ に写り、閉点は $x$ に写る。
\item $x \not = x'$ であるとし、有限生成体拡大 $K/\kappa(x')$ が与えられれば、
$\kappa(\eta)/\kappa(x')$（$f$ により誘導される拡大）が与えられた拡大と同型になるように
選ぶことができる。
\end{enumerate}
\end{lemma}

\begin{proof}
まず $x' \leadsto x$ が $X$ における特殊化で、$x \not = x'$ であると仮定し、
$K/\kappa(x')$ を有限生成体拡大とする。スキームの章の補題
\ref{schemes-lemma-specialize-points} と、スキームの章の補題
\ref{schemes-lemma-characterize-points} に続く議論により、環準同型
$\mathcal{O}_{X, x} \to \kappa(x') \to K$ を得る。$x \not = x'$ なので、
$\mathcal{O}_{X, x}$ の $K$ における像は体ではない（詳細は省略する）。
$R \subset K$ を、分数体が $K$ であり、$\mathcal{O}_{X, x} \to K$ の像を支配する
任意の離散付値環とする。代数の章の補題 \ref{algebra-lemma-exists-dvr} を参照せよ。
環準同型 $\mathcal{O}_{X, x} \to R$ は射 $f : \Spec(R) \to X$ を誘導する。
スキームの章の補題 \ref{schemes-lemma-morphism-from-spec-local-ring} を参照せよ。
構成により、この射は望む性質をすべて持つ。$x = x'$ ならば
$R = \kappa(x)[t]_{(t)}$ と置けばよい。
\end{proof}

\begin{lemma}
\label{properties-lemma-thin-infinite-sequence}
$S$ を Noether スキームとする。$T \subset S$ を無限部分集合とする。このとき、
無限部分集合 $T' \subset T$ であって、その点 $T'$ の間に非自明な特殊化が存在しない
ものが存在する。
\end{lemma}

\begin{proof}
$T_0 \subset T$ を、$t \in T$ であって $T$ の別の点に特殊化しないもの全体の集合とする。
$T_0$ が無限ならば $T' = T_0$ とすればよい。したがって $T_0$ は有限であると
仮定してよい。帰納的に、$i > 0$ に対して集合 $T_i \subset T$ を考える。その元
$t \in T$ は次を満たす：
\begin{enumerate}
\item $t \not \in T_{i - 1} \cup T_{i - 2} \cup \ldots \cup T_0$ である。
\item 非自明な特殊化 $t \leadsto t'$ であって $t' \in T_{i - 1}$ を満たすものが存在する。
\item 任意の非自明な特殊化 $t \leadsto t'$ であって $t' \in T$ を満たすものに対して、
$t' \in T_{i - 1} \cup T_{i - 2} \cup \ldots \cup T_0$ である。
\end{enumerate}
再び、$T_i$ が無限ならば $T' = T_i$ とすればよい。$d$ を局所環
$\mathcal{O}_{S, t}$（$t \in T_0$）の次元の最大値とする。$d$ は整数である。実際、
$T_0$ は有限であり、局所環の次元は代数の章の命題
\ref{algebra-proposition-dimension} により有限である。このとき
$T_i = \emptyset$ である（$i > d$）。実際、
$t \in T_i$ ならば、非自明な特殊化の列
$t = t_i \leadsto t_{i - 1} \leadsto \ldots \leadsto t_0$ であって $t_0 \in T_0$ を
満たすものを取れる。点
$t = t_i, t_{i - 1}, \ldots, t_0$ は $\Spec(\mathcal{O}_{S, t_0})$ に属するので
（スキームの章の補題 \ref{schemes-lemma-specialize-points}）、$i \leq d$ が分かる。
したがって $\bigcup T_i = T_d \cup \ldots \cup T_0$ は $T$ の有限部分集合である。

\medskip\noindent
$t \in T$ が $\bigcup T_i$ に属さないと仮定する。このとき
$t \leadsto t'$ という特殊化であって、$t' \in T$ かつ $t' \not \in \bigcup T_i$ を
満たすものが存在しなければ
ならない。（実際、$t$ のすべての特殊化が有限集合 $T_d \cup \ldots \cup T_0$ に属するなら、
$i$ であって、特殊化 $t \leadsto t'$ と $t' \in T_i$ が成り立つようなものの最大値があり、
構成により
$t \in T_{i + 1}$ となる。）したがって $T \setminus \bigcup T_i$ の点の間に、
非自明な特殊化の無限列
$$
t \leadsto t' \leadsto t'' \leadsto \ldots
$$
を得る。$S$ の台位相空間は補題
\ref{properties-lemma-locally-Noetherian-quasi-separated} により Noether なので、これは不可能である。
\end{proof}

\begin{lemma}
\label{properties-lemma-maximal-points}
$S$ を Noether スキームとする。$T \subset S$ を部分集合とする。$T_0 \subset T$ を、
$t \in T$ であって、非自明な特殊化 $t' \leadsto t$ で $t' \in T'$ を満たすものが
存在しないもの全体の集合とする。このとき、(a) $T_0$ の点の間に特殊化はなく、
(b) $T$ のすべての点は $T_0$ の点の特殊化であり、(c) $T$ と $T_0$ の閉包は等しい。
\end{lemma}

\begin{proof}
$\dim(\mathcal{O}_{S, s}) < \infty$（任意の $s \in S$）であることを思い出そう。
代数の章の命題 \ref{algebra-proposition-dimension} を参照せよ。$t \in T$ とする。
$t' \leadsto t$ ならば、次元論により
$\dim(\mathcal{O}_{S, t'}) \leq \dim(\mathcal{O}_{S, t})$ であり、等号が成り立つための
必要十分条件は $t' = t$ である。したがって、$t' \leadsto t$ を
$\dim(\mathcal{O}_{T, t'})$ が最小となるように選べば $t' \in T_0$ である。
言い換えれば、すべての $t \in T$ は
$T_0$ の元の特殊化である。
\end{proof}

\begin{lemma}
\label{properties-lemma-countable-dense-subset}
$S$ を Noether スキームとする。$T \subset S$ を無限稠密部分集合とする。このとき、
可算部分集合 $E \subset T$ であって $S$ において稠密なものが存在する。
\end{lemma}

\begin{proof}
$T'$ を、$s \in S$ であって、$\overline{\{s\}} \cap T$ が、その閉包が
$\overline{\{s\}}$ となる可算部分集合を含むもの全体の集合とする。
有限集合は可算なので $T \subset T'$ である。
$s \in T'$ に対して、そのような可算部分集合
$E_s \subset \overline{\{s\}} \cap T$ を選ぶ。
$E' = \{s_1, s_2, s_3, \ldots\} \subset T'$ を可算部分集合とする。このとき、$E'$ の
$S$ における閉包は、可算部分集合 $\bigcup_n E_{s_n}$（$T$ のもの）の閉包である。
したがって、$Z$ が
$E'$ の閉包の既約成分ならば、$Z$ の生成点は $T'$ に属する。

\medskip\noindent
$T'_0 \subset T'$ を、補題 \ref{properties-lemma-maximal-points} と同様に、$t \in T'$ であって、
非自明な特殊化 $t' \leadsto t$ で $t' \in T'$ を満たすものが存在しないもの全体の
部分集合とする。同補題の結果を以後断りなく用いる。
$T'_0$ が無限ならば、可算部分集合 $E' \subset T'_0$ を選ぶ。第一段落の議論により、
$E'$ の閉包の既約成分の生成点は $T'$ に属する。しかし、これらの点の一つが
$E' \subset T'_0$ の相異なる無限個の元に特殊化するので、これは矛盾である。
したがって $T'_0$ は有限であり、$T'_0 = \{s_1, \ldots, s_m\}$ と書く。このとき、
$S$（$T$ の閉包）は $\{s_1, \ldots, s_m\}$ の閉包に含まれ、これはさらに望むとおり
可算部分集合 $E_{s_1} \cup \ldots \cup E_{s_m} \subset T$ の閉包に含まれる。
\end{proof}




\section{Jacobson スキーム}
\label{properties-section-jacobson}

\noindent
すべての閉部分集合において閉点が稠密である空間を {\it Jacobson} ということを
思い出そう。位相の章の節 \ref{topology-section-space-jacobson} を参照せよ。

\begin{definition}
\label{properties-definition-jacobson}
スキーム $S$ の台位相空間が Jacobson であるとき、このスキームは {\it Jacobson} であるという。
\end{definition}

\noindent
環 $R$ が Jacobson であるとは、$R$ のすべての根基イデアルが極大イデアルの共通部分で
あることを思い出そう。代数の章の定義 \ref{algebra-definition-ring-jacobson} を参照せよ。

\begin{lemma}
\label{properties-lemma-affine-jacobson}
アフィンスキーム $\Spec(R)$ が Jacobson であるための必要十分条件は、環 $R$ が
Jacobson であることである。
\end{lemma}

\begin{proof}
これは代数の章の補題 \ref{algebra-lemma-jacobson} である。
\end{proof}

\noindent
Jacobson スキームを特徴づける標準的な結果は次のとおりである。直観的には
Jacobson $\Leftrightarrow$ 局所的に Jacobson であることを主張している。

\begin{lemma}
\label{properties-lemma-locally-jacobson}
$X$ をスキームとする。次の条件は同値である：
\begin{enumerate}
\item スキーム $X$ は Jacobson である。
\item スキーム $X$ は定義 \ref{properties-definition-locally-P} の意味で「局所的に Jacobson」である。
\item 任意のアフィン開集合 $U \subset X$ に対して環 $\mathcal{O}_X(U)$ は Jacobson である。
\item アフィン開被覆 $X = \bigcup U_i$ が存在し、各 $\mathcal{O}_X(U_i)$ は
Jacobson である。
\item 開被覆 $X = \bigcup X_j$ が存在し、各開部分スキーム $X_j$ は Jacobson である。
\end{enumerate}
さらに、$X$ が Jacobson ならば、すべての開部分スキームは Jacobson である。
\end{lemma}

\begin{proof}
補題の最後の主張は位相の章の補題 \ref{topology-lemma-jacobson-inherited} により成り立つ。
(5) と (1) の同値性は位相の章の補題 \ref{topology-lemma-jacobson-local} である。
したがって補題 \ref{properties-lemma-affine-jacobson} を用いると、
(1) $\Leftrightarrow$ (2) が分かる。補題の証明を終えるには、「Jacobson」が環の
局所的性質であることを示せば十分である。補題 \ref{properties-lemma-locally-P} を参照せよ。
Jacobson 環を一つの元で局所化した環は Jacobson である。代数の章の補題
\ref{algebra-lemma-Jacobson-invert-element} を参照せよ。$R$ を環とし、
$f_1, \ldots, f_n \in R$ が単位イデアルを生成し、各 $R_{f_i}$ が Jacobson であるとする。
すると $\Spec(R) = \bigcup D(f_i)$ はすべて Jacobson である開部分集合の合併だから、
再び位相の章の補題 \ref{topology-lemma-jacobson-local} により $\Spec(R)$ は
Jacobson である。これで定義 \ref{properties-definition-property-local} の第二の条件が証明された。
\end{proof}

\noindent
代数幾何で通常用いられる多くのスキームは Jacobson である。射の章の補題
\ref{morphisms-lemma-ubiquity-Jacobson-schemes} を参照せよ。ここでは次の興味深い場合を
挙げる。

\begin{lemma}
\label{properties-lemma-complement-closed-point-Jacobson}
Noether Jacobson スキームの例。
\begin{enumerate}
\item $(R, \mathfrak m)$ が Noether 局所環ならば、穿孔スペクトル
$\Spec(R) \setminus \{\mathfrak m\}$ は Jacobson スキームである。
\item $R$ が Jacobson 根基 $\text{rad}(R)$ を持つ Noether 環ならば、
$\Spec(R) \setminus V(\text{rad}(R))$ は Jacobson スキームである。
\item $(R, I)$ が Zariski 対であり（代数の続編の定義
\ref{more-algebra-definition-zariski-pair}）、$R$ が Noether ならば、
$\Spec(R) \setminus V(I)$ は Jacobson スキームである。
\end{enumerate}
\end{lemma}

\begin{proof}
(3) の証明。$\Spec(R) - V(I)$ は、アフィン開集合 $\Spec(R_f)$（$f \in I$）に
よる被覆を持つことに注意する。環 $R_f$ は代数の続編の補題
\ref{more-algebra-lemma-noetherian-zariski-jacobson-complement} により Jacobson である。
したがって補題 \ref{properties-lemma-locally-jacobson} により
$\Spec(R) \setminus V(I)$ は Jacobson である。(1) と (2) は (3) の特別な場合である。

\medskip\noindent
(1) の直接証明。$\Spec(R)$ は Noether スキームなので、$S$ は Noether スキームである
（補題 \ref{properties-lemma-locally-closed-in-Noetherian}）。したがって $S$ はソバーな Noether
位相空間である（スキームの章の補題 \ref{schemes-lemma-scheme-sober} を用いよ）。
矛盾を導くため、$S$ は Jacobson でないと仮定する。位相の章の補題
\ref{topology-lemma-non-jacobson-Noetherian-characterize} により、非閉点 $\xi \in S$ であって
$\{\xi\}$ が局所閉となるものが存在する。これは素イデアル $\mathfrak p \subset R$ であって、
(1) 素イデアル $\mathfrak q$, $\mathfrak p \subset \mathfrak q \subset \mathfrak m$ が存在し、
両方の包含が真であり、(2) $\{\mathfrak p\}$ が $\Spec(R/\mathfrak p)$ において開である
ものに対応する。これは代数の章の補題
\ref{algebra-lemma-Noetherian-local-domain-dim-2-infinite-opens} により不可能である。
\end{proof}




\section{正規スキーム}
\label{properties-section-normal}

\noindent
環 $R$ のすべての局所環が正規整域であるとき、この環は正規であるということを思い出そう。
代数の章の定義 \ref{algebra-definition-ring-normal} を参照せよ。正規整域とは、その分数体に
おいて整閉である整域をいう。代数の章の定義 \ref{algebra-definition-domain-normal} を参照せよ。
したがって、正規スキームを次のように定義することは自然である。

\begin{definition}
\label{properties-definition-normal}
スキーム $X$ が {\it 正規} であるための必要十分条件は、すべての $x \in X$ に対して
局所環 $\mathcal{O}_{X, x}$ が正規整域であることである。
\end{definition}

\noindent
これは EGA で用いられている定義のようである。\cite[0, 4.1.4]{EGA} を参照せよ。
$X = \Spec(A)$ とし、$A$ が被約であると仮定する。このとき $X$ が正規であることは、
$A$ がその全商環において整閉であることと同値ではない。しかし $A$ が Noether ならば
同値である（代数の章の補題 \ref{algebra-lemma-characterize-reduced-ring-normal} を参照）。

\begin{lemma}
\label{properties-lemma-locally-normal}
$X$ をスキームとする。次の条件は同値である：
\begin{enumerate}
\item スキーム $X$ は正規である。
\item 任意のアフィン開集合 $U \subset X$ に対して環 $\mathcal{O}_X(U)$ は正規である。
\item アフィン開被覆 $X = \bigcup U_i$ が存在し、各 $\mathcal{O}_X(U_i)$ は正規である。
\item 開被覆 $X = \bigcup X_j$ が存在し、各開部分スキーム $X_j$ は正規である。
\end{enumerate}
さらに、$X$ が正規ならば、すべての開部分スキームは正規である。
\end{lemma}

\begin{proof}
定義から明らかである。
\end{proof}

\begin{lemma}
\label{properties-lemma-normal-reduced}
正規スキームは被約である。
\end{lemma}

\begin{proof}
定義から直ちに従う。
\end{proof}

\begin{lemma}
\label{properties-lemma-integral-normal}
$X$ を整スキームとする。このとき、$X$ が正規であるための必要十分条件は、任意の
非空アフィン開集合 $U \subset X$ に対して環 $\mathcal{O}_X(U)$ が正規整域であることである。
\end{lemma}

\begin{proof}
これは代数の章の補題 \ref{algebra-lemma-normality-is-local} から従う。
\end{proof}

\begin{lemma}
\label{properties-lemma-normal-locally-finite-nr-irreducibles}
$X$ を、任意の準コンパクト開集合が有限個の既約成分を持つようなスキームとする。
次の条件は同値である：
\begin{enumerate}
\item $X$ は正規である。
\item $X$ は正規整スキームの非交和である。
\end{enumerate}
\end{lemma}

\begin{proof}
(2) が (1) を含意することは定義から直ちに分かる。$X$ を、任意の準コンパクト開集合が
有限個の既約成分を持つ正規スキームとする。$X$ がアフィンならば、代数の章の補題
\ref{algebra-lemma-characterize-reduced-ring-normal} により $X$ は (2) を満たす。
一般の $X$ に対して、$X = \bigcup X_i$ をアフィン開被覆とする。各 $X_i$ も有限個しか
既約成分を持たず、各 $X_i$ に対して補題が成り立つことに注意する。$T \subset X$ を
既約成分とする。アフィンの場合により、各共通部分 $T \cap X_i$ は $X_i$ において開であり、
正規整スキームである。したがって $T \subset X$ は開であり、正規整スキームである。
これにより、$X$ はその既約成分の非交和であり、それらが正規整スキームであることが
証明された。
\end{proof}

\begin{lemma}
\label{properties-lemma-normal-Noetherian}
$X$ を Noether スキームとする。次の条件は同値である：
\begin{enumerate}
\item $X$ は正規である。
\item $X$ は正規整スキームの有限非交和である。
\end{enumerate}
\end{lemma}

\begin{proof}
これは補題 \ref{properties-lemma-normal-locally-finite-nr-irreducibles} の特別な場合である。
実際、Noether スキームの台位相空間は Noether である（補題
\ref{properties-lemma-Noetherian-topology} および位相の章の補題
\ref{topology-lemma-Noetherian}）。
\end{proof}

\begin{lemma}
\label{properties-lemma-normal-locally-Noetherian}
$X$ を局所 Noether スキームとする。次の条件は同値である：
\begin{enumerate}
\item $X$ は正規である。
\item $X$ は整正規スキームの非交和である。
\end{enumerate}
\end{lemma}

\begin{proof}
省略する。ヒント：補題 \ref{properties-lemma-normal-Noetherian} から従う純粋に位相的な主張である。
\end{proof}

\begin{remark}
\label{properties-remark-normal-connected-irreducible}
$X$ を正規スキームとする。$X$ が局所 Noether ならば、補題
\ref{properties-lemma-normal-locally-Noetherian} により、$X$ が整であるための必要十分条件は
$X$ が連結であることである。しかし、連結アフィンスキーム $X$ であって、
$\mathcal{O}_{X, x}$ がすべての $x \in X$ に対して整域であるにもかかわらず、$X$ が
既約でないものが存在する。例の章の節
\ref{examples-section-connected-locally-integral-not-integral} を参照せよ。
この例は正規スキームでさえある（証明は省略する）ので、注意されたい。
\end{remark}

\begin{lemma}
\label{properties-lemma-normal-integral-sections}
\begin{slogan}
正規スキーム上の関数環は正規である。
\end{slogan}
$X$ を整正規スキームとする。このとき $\Gamma(X, \mathcal{O}_X)$ は正規整域である。
\end{lemma}

\begin{proof}
$R = \Gamma(X, \mathcal{O}_X)$ と置く。$R$ が整域であることは明らかである。
$f = a/b$ を、その分数体の元であって $R$ 上整であるものとする。すなわち
$f^d + \sum_{i = 0, \ldots, d - 1} a_i f^i = 0$ であり、ここで
$a_i \in R$ であるとする。
$U \subset X$ を非空アフィン開集合とする。$b \in R$ は零でなく、$X$ は整なので、
$b|_U \in \mathcal{O}_X(U)$ も零ではない。したがって $a/b$ は
$\mathcal{O}_X(U)$ の分数体の元であり、$\mathcal{O}_X(U)$ 上整である
（同じ多項式 $f^d + \sum_{i = 0, \ldots, d - 1} a_i|_U f^i = 0$ を $U$ 上で
用いることができるからである）。$\mathcal{O}_X(U)$ は正規整域なので（補題
\ref{properties-lemma-integral-normal}）、$f_U = (a|_U)/(b|_U) \in \mathcal{O}_X(U)$ である。
$f_U|_V = f_V$ は、$V \subset U \subset X$ が非空アフィン開集合ならば常に成り立つことは
明らかである。したがって局所切断 $f_U$ は貼り合わさって元
$g \in R = \Gamma(X, \mathcal{O}_X)$ を与える。このとき $bg$ と $a$ は
$\mathcal{O}_X(U)$ の同じ元に制限される（上のようなすべての $U$ に対して）。
したがって $bg = a$、言い換えれば $g$ は $f$ に写る（$R$ の分数体において）。
\end{proof}




\section{Cohen--Macaulay スキーム}
\label{properties-section-Cohen-Macaulay}

\noindent
代数の章の定義 \ref{algebra-definition-local-ring-CM} を参照して思い出そう。
局所 Noether 環 $(R, \mathfrak m)$ が Cohen--Macaulay であるとは
$\text{depth}_{\mathfrak m}(R) = \dim(R)$ であることをいう。また Noether 環 $R$ が
Cohen--Macaulay であるとは、すべての局所環 $R_{\mathfrak p}$（$R$ のもの）が
Cohen--Macaulay であることをいう。代数の章の定義 \ref{algebra-definition-ring-CM} を
参照せよ。

\begin{definition}
\label{properties-definition-Cohen-Macaulay}
$X$ をスキームとする。$X$ が {\it Cohen--Macaulay} であるとは、任意の $x \in X$ に
対してアフィン開近傍 $U \subset X$（$x$ のもの）であって、環 $\mathcal{O}_X(U)$ が
Noether かつ Cohen--Macaulay であるものが存在することをいう。
\end{definition}

\begin{lemma}
\label{properties-lemma-characterize-Cohen-Macaulay}
$X$ をスキームとする。次の条件は同値である：
\begin{enumerate}
\item $X$ は Cohen--Macaulay である。
\item $X$ は局所 Noether であり、そのすべての局所環は Cohen--Macaulay である。
\item $X$ は局所 Noether であり、任意の閉点 $x \in X$ に対して局所環
$\mathcal{O}_{X, x}$ は Cohen--Macaulay である。
\end{enumerate}
\end{lemma}

\begin{proof}
代数の章の補題 \ref{algebra-lemma-localize-CM} は、Cohen--Macaulay 局所環の局所化が
Cohen--Macaulay であることを述べている。これを補題 \ref{properties-lemma-locally-Noetherian}、
局所 Noether スキーム上の閉点の存在（補題
\ref{properties-lemma-locally-Noetherian-closed-point}）、および定義と組み合わせれば補題が従う。
\end{proof}

\begin{lemma}
\label{properties-lemma-locally-Cohen-Macaulay}
$X$ をスキームとする。次の条件は同値である：
\begin{enumerate}
\item スキーム $X$ は Cohen--Macaulay である。
\item 任意のアフィン開集合 $U \subset X$ に対して環 $\mathcal{O}_X(U)$ は
Noether かつ Cohen--Macaulay である。
\item アフィン開被覆 $X = \bigcup U_i$ が存在し、各 $\mathcal{O}_X(U_i)$ は
Noether かつ Cohen--Macaulay である。
\item 開被覆 $X = \bigcup X_j$ が存在し、各開部分スキーム $X_j$ は
Cohen--Macaulay である。
\end{enumerate}
さらに、$X$ が Cohen--Macaulay ならば、すべての開部分スキームは Cohen--Macaulay である。
\end{lemma}

\begin{proof}
補題 \ref{properties-lemma-locally-Noetherian} と \ref{properties-lemma-characterize-Cohen-Macaulay} を
組み合わせよ。
\end{proof}

\noindent
Cohen--Macaulay スキームと深さについての詳細は、スキームのコホモロジーの章の節
\ref{coherent-section-depth} にある。







\section{正則スキーム}
\label{properties-section-regular}

\noindent
代数の章の定義 \ref{algebra-definition-regular-local} を参照して思い出そう。
局所 Noether 環 $(R, \mathfrak m)$ が {\it 正則}であるとは、$\mathfrak m$ が
$\dim(R)$ 個の元で生成されることをいう。また Noether 環 $R$ が {\it 正則}であるとは、
すべての局所環 $R_{\mathfrak p}$（$R$ のもの）が正則であることをいう。代数の章の定義
\ref{algebra-definition-regular} を参照せよ。

\begin{definition}
\label{properties-definition-regular}
$X$ をスキームとする。$X$ が {\it 正則}、または {\it 非特異}であるとは、任意の
$x \in X$ に対してアフィン開近傍 $U \subset X$（$x$ のもの）であって、環
$\mathcal{O}_X(U)$ が Noether かつ正則であるものが存在することをいう。
\end{definition}

\begin{lemma}
\label{properties-lemma-characterize-regular}
$X$ をスキームとする。次の条件は同値である：
\begin{enumerate}
\item $X$ は正則である。
\item $X$ は局所 Noether であり、そのすべての局所環は正則である。
\item $X$ は局所 Noether であり、任意の閉点 $x \in X$ に対して局所環
$\mathcal{O}_{X, x}$ は正則である。
\end{enumerate}
\end{lemma}

\begin{proof}
代数の章の定義 \ref{algebra-definition-regular} に先立つ議論により、正則局所環の局所化が
正則であることが分かっている。これを補題 \ref{properties-lemma-locally-Noetherian}、局所 Noether
スキーム上の閉点の存在（補題 \ref{properties-lemma-locally-Noetherian-closed-point}）、および定義と
組み合わせれば補題が従う。
\end{proof}

\begin{lemma}
\label{properties-lemma-locally-regular}
$X$ をスキームとする。次の条件は同値である：
\begin{enumerate}
\item スキーム $X$ は正則である。
\item 任意のアフィン開集合 $U \subset X$ に対して環 $\mathcal{O}_X(U)$ は
Noether かつ正則である。
\item アフィン開被覆 $X = \bigcup U_i$ が存在し、各 $\mathcal{O}_X(U_i)$ は
Noether かつ正則である。
\item 開被覆 $X = \bigcup X_j$ が存在し、各開部分スキーム $X_j$ は正則である。
\end{enumerate}
さらに、$X$ が正則ならば、すべての開部分スキームは正則である。
\end{lemma}

\begin{proof}
補題 \ref{properties-lemma-locally-Noetherian} と \ref{properties-lemma-characterize-regular} を組み合わせよ。
\end{proof}

\begin{lemma}
\label{properties-lemma-regular-normal}
正則スキームは正規である。
\end{lemma}

\begin{proof}
代数の章の補題 \ref{algebra-lemma-regular-normal} を参照せよ。
\end{proof}



























\section{次元}
\label{properties-section-dimension}

\noindent
スキームの次元とは、その台位相空間の次元にほかならない。

\begin{definition}
\label{properties-definition-dimension}
$X$ をスキームとする。
\begin{enumerate}
\item $X$ の {\it 次元}とは、位相空間としての $X$ の次元にほかならない。位相の章の定義
\ref{topology-definition-Krull} を参照せよ。
\item $x \in X$ に対し、$\dim_x(X)$ で、位相の章の定義
\ref{topology-definition-Krull} における $X$ の台位相空間の $x$ での次元を表す。
$\dim_x(X)$ を {\it $X$ の $x$ における次元}という。
\end{enumerate}
\end{definition}

\noindent
スキームは sober な台位相空間を持つので（スキームの章の補題
\ref{schemes-lemma-scheme-sober}）、$X$ の次元は、長さ $n$ の鎖
$$
T_0 \subset T_1 \subset \ldots \subset T_n
$$
であって $X$ の既約閉部分集合からなるものの上限として、または長さ $n$ の特殊化の鎖
$$
\xi_n \leadsto \xi_{n - 1} \leadsto \ldots \leadsto \xi_0
$$
であって $X$ の点からなるものの上限として計算できる。

\begin{lemma}
\label{properties-lemma-dimension}
$X$ をスキームとする。次の量は等しい：
\begin{enumerate}
\item $X$ の次元。
\item $X$ の局所環の次元の上限。
\item $\dim_x(X)$ の、$x \in X$ にわたる上限。
\end{enumerate}
\end{lemma}

\begin{proof}
$X$ の点の特殊化の鎖
$$
\xi_n \leadsto \xi_{n - 1} \leadsto \ldots \leadsto \xi_0
$$
が与えられたとき、すべての点 $\xi_i$ は、スキームの章の補題
\ref{schemes-lemma-specialize-points} により、$X$ の $\xi_0$ における局所環の素イデアルに
対応することに注意せよ。したがって、$X$ の次元はその局所環の次元の上限である。
特に、$\dim_x(X) \geq \dim(\mathcal{O}_{X, x})$ である。実際、$\dim_x(X)$ は $x$ の
開近傍の次元の最小値だからである。よって
$\sup_{x \in X} \dim_x(X) \geq \dim(X)$ である。一方、
$\sup_{x \in X} \dim_x(X) \leq \dim(X)$ であることは、$\dim(U) \leq \dim(X)$ が
任意の $X$ の開部分集合について成り立つので明らかである。
\end{proof}

\begin{lemma}
\label{properties-lemma-codimension-local-ring}
$X$ をスキームとする。$Y \subset X$ を既約閉部分集合とし、$\xi \in Y$ をその生成点とする。
このとき
$$
\text{codim}(Y, X) = \dim(\mathcal{O}_{X, \xi})
$$
である。ここで余次元は位相の章の定義
\ref{topology-definition-codimension} による。
\end{lemma}

\begin{proof}
位相の章の補題 \ref{topology-lemma-codimension-at-generic-point} により、$X$ を $\xi$ の
アフィン開近傍で置き換えてよい。この場合、主張は代数の章の補題
\ref{algebra-lemma-irreducible-components-containing-x} から容易に従う。
\end{proof}

\begin{lemma}
\label{properties-lemma-generic-point}
$X$ をスキームとし、$x \in X$ とする。$x$ が $X$ の既約成分の生成点であるための必要十分条件は
$\dim(\mathcal{O}_{X, x}) = 0$ である。
\end{lemma}

\begin{proof}
例えば補題 \ref{properties-lemma-codimension-local-ring} から従う。
\end{proof}

\begin{lemma}
\label{properties-lemma-locally-Noetherian-dimension-0}
次元 $0$ の局所 Noether スキームは、Artin 局所環のスペクトルの非交和である。
\end{lemma}

\begin{proof}
次元 $0$ の Noether 環は Artin 局所環の有限積である。代数の章の命題
\ref{algebra-proposition-dimension-zero-ring} を参照せよ。したがって、局所
Noether スキーム $X$ の次元が $0$ ならば、そのアフィン開集合は離散な台位相空間を持つ。
これは $X$ の位相が
離散であることを意味する。以上から補題は容易に従う。
\end{proof}

\begin{lemma}
\label{properties-lemma-dimension-zero}
\begin{reference}
2016 年 6 月 4 日付 Ofer Gabber からの電子メール
\end{reference}
$X$ を次元零のスキームとする。次の条件は同値である：
\begin{enumerate}
\item $X$ は準分離である。
\item $X$ は分離的である。
\item $X$ はハウスドルフである。
\item すべてのアフィン開集合は閉である。
\end{enumerate}
この場合、$X$ の連結成分は点であり、$X$ のすべての準コンパクト開集合はアフィンである。
特に、$X$ が準コンパクトならば、$X$ はアフィンである。
\end{lemma}

\begin{proof}
$X$ の次元は零なので、任意のアフィン開集合 $U \subset X$ に対して空間 $U$ は
プロ有限であり、以下で自由に用いるほかの多くの性質も満たす。代数の章の補題
\ref{algebra-lemma-ring-with-only-minimal-primes} を参照せよ。アフィン開被覆
$X = \bigcup U_i$ を選ぶ。

\medskip\noindent
(4) が成り立つとする。このとき $U_i \cap U_j$ は $U_i$ の閉部分集合であるから準コンパクトであり、
従ってスキームの章の補題 \ref{schemes-lemma-characterize-quasi-separated} により $X$ は
準分離である。よって (1) が成り立つ。

\medskip\noindent
(1) が成り立つとする。このとき $U_i \cap U_j$ は $U_i$ の準コンパクト開集合なので
$U_i$ において閉である。従って $U_i \cap U_j \to U_i$ は像が閉である開埋め込みであり、
ゆえに閉埋め込みである。特に $U_i \cap U_j$ はアフィンであり、
$\mathcal{O}(U_i) \to \mathcal{O}_X(U_i \cap U_j)$ は全射である。
従ってスキームの章の補題 \ref{schemes-lemma-characterize-separated} により $X$ は分離的である。
よって (2) が成り立つ。

\medskip\noindent
(2) を仮定し、$x, y \in X$ とする。$x \in U_i$ としよう。$y \in U_i$ でもあるなら、
$x$ と $y$ の互いに交わらない開近傍を取ることができる。なぜなら $U_i$ はハウスドルフだからである。
$y \not \in U_i$ かつ $y \in U_j$ としよう。このとき $y \not \in U_i \cap U_j$ であり、
この集合は $U_j$ のアフィン開集合なので $U_j$ において閉である。従って
$y$ の、$U_i$ と交わらない開近傍を取ることができ、$X$ はハウスドルフであると結論する。
よって (3) が成り立つ。

\medskip\noindent
(3) を仮定する。$U \subset X$ をアフィン開集合とする。位相の章の補題
\ref{topology-lemma-quasi-compact-in-Hausdorff} により、$U$ は $X$ において閉である。
これで (4) が成り立つことが示された。

\medskip\noindent
$X$ が同値な条件 (1) -- (4) を満たすと仮定する。補題の最後の主張を証明する。
$x, y \in X$ かつ $x \not = y$ としよう。$y$ は $x$ に特殊化しないので、
アフィン開集合 $U \subset X$ であって $x \in U$ かつ $y \not \in U$ を満たすものを選べる。
すると $X = U \amalg (X \setminus U)$
は開かつ閉な部分集合への分解であり、$x$ と $y$ が $X$ の同じ連結成分に属さないことが分かる。
次に、$U \subset X$ を準コンパクト開集合とする。アフィン開集合の合併として
$U = U_1 \cup \ldots \cup U_n$ と書く。$n$ に関する帰納法で $U$ がアフィンであることを
証明する。これは直ちに $n = 2$ の場合に帰着する。この場合
$U = (U_1 \setminus U_2) \amalg (U_1 \cap U_2) \amalg (U_2 \setminus U_1)$
であり、上の議論により各部分はアフィンである。
\end{proof}

\begin{lemma}
\label{properties-lemma-isolated-dim-0-locally-noetherian}
$x$ を局所 Noether スキーム $X$ の点とする。このとき $\dim_x(X) = 0$ であるための
必要十分条件は、$x$ が $X$ の孤立点であることである。
\end{lemma}

\begin{proof}
$x$ が孤立点ならば、$\{x\}$ は $X$ の開部分集合であり（位相の章の定義
\ref{topology-definition-isolated-point}）、従って定義により
$\dim_x(X) = \dim(\{x\}) = 0$ である。逆に、$\dim_x(X) = 0$ ならば、
開近傍 $U \subset X$ であって、$x$ を含み $\dim(U) = 0$ を満たすものが存在する
（位相の章の定義 \ref{topology-definition-Krull}）。補題
\ref{properties-lemma-locally-Noetherian-dimension-0} により、$U$ の位相は離散であり、従って
$x$ は孤立点である。
\end{proof}





\section{鎖状スキーム}
\label{properties-section-catenary}

\noindent
位相空間 $X$ が {\it 鎖状}であるとは、既約閉部分集合の任意の対
$T \subset T'$ に対して、既約閉部分集合の極大鎖
$$
T = T_0 \subset T_1 \subset \ldots \subset T_e = T'
$$
が存在し、かつそのようなすべての鎖が同じ長さを持つことをいう。位相の章の定義
\ref{topology-definition-catenary} を参照せよ。

\begin{definition}
\label{properties-definition-catenary}
$S$ をスキームとする。$S$ の台位相空間が鎖状であるとき、$S$ は {\it 鎖状}であるという。
\end{definition}

\noindent
環 $A$ が {\it 鎖状}であるとは、素イデアルの任意の対
$\mathfrak p \subset \mathfrak q$ に対して素イデアルの極大鎖
$$
\mathfrak p =
\mathfrak p_0 \subset \ldots \subset \mathfrak p_e
= \mathfrak q
$$
が存在し、それらがすべて同じ長さを持つことをいう。代数の章の定義
\ref{algebra-definition-catenary} を参照せよ。

\begin{lemma}
\label{properties-lemma-catenary-local}
$S$ をスキームとする。次の条件は同値である：
\begin{enumerate}
\item $S$ は鎖状である。
\item $S$ の開被覆であって、そのすべての要素が鎖状スキームであるものが存在する。
\item 任意のアフィン開集合 $\Spec(R) = U \subset S$ に対して環 $R$ は鎖状である。
\item アフィン開被覆 $S = \bigcup U_i$ であって、各 $U_i$ が鎖状環のスペクトルであるものが
存在する。
\end{enumerate}
さらに、この場合 $S$ の任意の局所閉部分スキームも鎖状である。
\end{lemma}

\begin{proof}
位相の章の補題 \ref{topology-lemma-catenary} と代数の章の補題
\ref{algebra-lemma-catenary} を組み合わせよ。
\end{proof}

\begin{lemma}
\label{properties-lemma-catenary-dimension-function}
$S$ を局所 Noether スキームとする。次の条件は同値である：
\begin{enumerate}
\item $S$ は鎖状である。
\item Zariski 位相に関して局所的に、$S$ 上の次元関数が存在する
（位相の章の定義 \ref{topology-definition-dimension-function} を参照）。
\end{enumerate}
\end{lemma}

\begin{proof}
これは位相の章の補題
\ref{topology-lemma-catenary}、
\ref{topology-lemma-dimension-function-catenary}、
\ref{topology-lemma-locally-dimension-function}、
スキームの章の補題 \ref{schemes-lemma-scheme-sober}、
および最後に補題 \ref{properties-lemma-Noetherian-topology} から従う。
\end{proof}

\noindent
実際、スキームが鎖状であるための必要十分条件は、その局所環が鎖状であることである。

\begin{lemma}
\label{properties-lemma-catenary-local-rings-catenary}
$X$ をスキームとする。次の条件は同値である：
\begin{enumerate}
\item $X$ は鎖状である。
\item 任意の $x \in X$ に対して局所環 $\mathcal{O}_{X, x}$ は鎖状である。
\end{enumerate}
\end{lemma}

\begin{proof}
$X$ が鎖状であると仮定する。$x \in X$ とする。補題 \ref{properties-lemma-catenary-local} により、
$X$ を $x$ のアフィン開近傍で置き換えてよく、そのとき
$\Gamma(X, \mathcal{O}_X)$ は鎖状環である。代数の章の補題
\ref{algebra-lemma-localization-catenary} により、鎖状環の任意の局所化は鎖状である。
従って $\mathcal{O}_{X, x}$ は鎖状である。

\medskip\noindent
逆に、$X$ のすべての局所環が鎖状であると仮定する。$Y \subset Y'$ を $X$ の既約閉部分集合の
包含とする。$\xi \in Y$ を生成点とする。$\mathfrak p \subset \mathcal{O}_{X, \xi}$ を
$Y'$ の生成点に対応する素イデアルとする。スキームの章の補題
\ref{schemes-lemma-specialize-points} を参照せよ。同じ補題により、$X$ の既約閉部分集合であって
$Y$ と $Y'$ の間にあるものは、素イデアル $\mathfrak q \subset \mathcal{O}_{X, \xi}$ であって
$\mathfrak p \subset \mathfrak q \subset \mathfrak m_{\xi}$ を満たすものに対応する。
従って、$\mathcal{O}_{X, \xi}$ は鎖状環なので、これらのすべての極大鎖は有限で同じ長さを持つ。
\end{proof}






\section{Serre の条件}
\label{properties-section-Rk}

\noindent
ここでは、しばしば有用となる二つの技術的概念を導入する。
スキームのコホモロジーの章の節 \ref{coherent-section-depth} も参照せよ。

\begin{definition}
\label{properties-definition-Rk}
$X$ を局所 Noether スキームとし、$k \geq 0$ とする。
\begin{enumerate}
\item $X$ が {\it 余次元 $k$ で正則}である、または $X$ が性質 {\it $(R_k)$} を持つとは、
任意の $x \in X$ に対して
$$
\dim(\mathcal{O}_{X, x}) \leq k
\Rightarrow
\mathcal{O}_{X, x}\text{ is regular}
$$
が成り立つことをいう。
\item $X$ が性質 {\it $(S_k)$} を持つとは、任意の $x \in X$ に対して
$\text{depth}(\mathcal{O}_{X, x}) \geq \min(k, \dim(\mathcal{O}_{X, x}))$ が成り立つことをいう。
\end{enumerate}
\end{definition}

\noindent
「余次元 $k$ で正則」という表現に意味があるのは、節 \ref{properties-section-catenary} で、
$Y \subset X$ が生成点 $x$ を持つ既約閉部分集合ならば
$\dim(\mathcal{O}_{X, x}) = \text{codim}(Y, X)$ であることを見たからである。
例えば条件 $(R_0)$ は、任意の生成点 $\eta \in X$（$X$ の既約成分のもの）に対して局所環
$\mathcal{O}_{X, \eta}$ が体であることを意味する。しかし一般の Noether スキームでは
$X$ の正則軌跡が悪い振る舞いをすることがあるので、注意が必要である。

\begin{lemma}
\label{properties-lemma-scheme-regular-iff-all-Rk}
$X$ を局所 Noether スキームとする。このとき $X$ が正則であるための必要十分条件は、
$X$ が $(R_k)$ をすべての $k \geq 0$ に対して持つことである。
\end{lemma}

\begin{proof}
補題 \ref{properties-lemma-characterize-regular} と定義から従う。
\end{proof}

\begin{lemma}
\label{properties-lemma-scheme-CM-iff-all-Sk}
$X$ を局所 Noether スキームとする。このとき $X$ が Cohen--Macaulay であるための必要十分条件は、
$X$ が $(S_k)$ をすべての $k \geq 0$ に対して持つことである。
\end{lemma}

\begin{proof}
補題 \ref{properties-lemma-characterize-Cohen-Macaulay} により局所環を調べることに帰着する。
従って、Noether 局所環が Cohen--Macaulay であるための必要十分条件は、その深さが次元に等しい
ことであるから、補題は成り立つ。
\end{proof}

\begin{lemma}
\label{properties-lemma-criterion-reduced}
$X$ を局所 Noether スキームとする。このとき $X$ が被約であるための必要十分条件は、
$X$ が性質 $(S_1)$ と $(R_0)$ を持つことである。
\end{lemma}

\begin{proof}
これは代数の章の補題 \ref{algebra-lemma-criterion-reduced} である。
\end{proof}

\begin{lemma}
\label{properties-lemma-criterion-normal}
$X$ を局所 Noether スキームとする。このとき $X$ が正規であるための必要十分条件は、
$X$ が性質 $(S_2)$ と $(R_1)$ を持つことである。
\end{lemma}

\begin{proof}
これは代数の章の補題 \ref{algebra-lemma-criterion-normal} である。
\end{proof}

\begin{lemma}
\label{properties-lemma-normal-dimension-1-regular}
$X$ を、正規かつ次元 $\leq 1$ の局所 Noether スキームとする。このとき $X$ は正則である。
\end{lemma}

\begin{proof}
補題 \ref{properties-lemma-criterion-normal} と定義から従う。
\end{proof}

\begin{lemma}
\label{properties-lemma-normal-dimension-2-Cohen-Macaulay}
$X$ を、正規かつ次元 $\leq 2$ の局所 Noether スキームとする。このとき $X$ は
Cohen--Macaulay である。
\end{lemma}

\begin{proof}
補題 \ref{properties-lemma-criterion-normal} と定義から従う。
\end{proof}






\section{日本スキームと永田スキーム}
\label{properties-section-nagata}

\noindent
この節で扱う概念は、EGA では目立った形で定義されていない。
「普遍的日本スキーム」は \cite[IV Corollary 5.11.4]{EGA} で言及され、定義されている。
「日本スキーム」は \cite[IV Remark 10.4.14 (ii)]{EGA} で言及されているが、
定義は与えられていない。以下で定義する永田スキームは、文献のいくつかの箇所に現れる
（例えば \cite[Definition 8.2.30]{Liu} および \cite[Page 142]{Greco} を参照）。

\medskip\noindent
簡単に復習しておく。整域 $R$ が {\it 日本環} であるとは、その商体の任意の有限次拡大における
$R$ の整閉包が $R$ 上有限であることをいう。環 $R$ が {\it 普遍的日本環} であるとは、
任意の有限型環準同型 $R \to S$ に対して、$S$ が整域ならば $S$ が日本環であることをいう。
環 $R$ が {\it 永田環} であるとは、Noether 環であり、$R/\mathfrak p$ が任意の素イデアル
$\mathfrak p$（$R$ の素イデアル）に対して日本環であることをいう。

\begin{definition}
\label{properties-definition-nagata}
$X$ をスキームとする。
\begin{enumerate}
\item $X$ が整であると仮定する。$X$ を {\it 日本スキーム} と呼ぶのは、任意の
$x \in X$ に対して、アフィン開近傍 $x \in U \subset X$ であって、その環
$\mathcal{O}_X(U)$ が日本環であるものが存在するときである
（代数の章の定義 \ref{algebra-definition-N} を参照）。
\item $X$ を {\it 普遍的日本スキーム} と呼ぶのは、任意の $x \in X$ に対して、
アフィン開近傍 $x \in U \subset X$ であって、その環 $\mathcal{O}_X(U)$ が普遍的日本環である
ものが存在するときである
（代数の章の定義 \ref{algebra-definition-nagata} を参照）。
\item $X$ を {\it 永田スキーム} と呼ぶのは、任意の $x \in X$ に対して、
アフィン開近傍 $x \in U \subset X$ であって、その環 $\mathcal{O}_X(U)$ が永田環である
ものが存在するときである
（代数の章の定義 \ref{algebra-definition-nagata} を参照）。
\end{enumerate}
\end{definition}

\noindent
永田スキームであることは、局所 Noether かつ普遍的日本スキームであることと同値である。
補題 \ref{properties-lemma-nagata-universally-Japanese} を参照せよ。

\begin{remark}
\label{properties-remark-non-integral-Japanese}
\cite{Hoobler-finite} では、（局所 Noether）スキーム $X$ が日本スキームであるとは、
任意の $x \in X$ と任意の随伴素イデアル $\mathfrak p$（$\mathcal{O}_{X, x}$ のもの）に対して、
環 $\mathcal{O}_{X, x}/\mathfrak p$ が日本環であることをいう。この定義は用いない。
実際、局所環が優秀（従って特に日本環）であるにもかかわらず、その正規化が有限でない
1 次元 Noether 整域が存在する。\cite[Example 1]{Hochster-loci}、\cite{Heinzer-Levy}、
または \cite[Expos\'e XIX]{Traveaux} を参照せよ。
一方、スキーム $X$ を日本スキームと呼ぶ条件として、任意のアフィン開集合
$\Spec(A) \subset X$ に対し、$A/\mathfrak p$ が日本環であることを、任意の随伴素イデアル
$\mathfrak p$（$A$ のもの）に要求すれば、この問題を回避できる。
\end{remark}

\begin{lemma}
\label{properties-lemma-nagata-locally-Noetherian}
永田スキームは局所 Noether である。
\end{lemma}

\begin{proof}
これは、永田環が定義により Noether 環であることから従う。
\end{proof}

\begin{lemma}
\label{properties-lemma-locally-Japanese}
$X$ を整スキームとする。次の条件は同値である：
\begin{enumerate}
\item スキーム $X$ は日本スキームである。
\item 任意のアフィン開集合 $U \subset X$ に対して、整域 $\mathcal{O}_X(U)$ は日本環である。
\item アフィン開被覆 $X = \bigcup U_i$ であって、各 $\mathcal{O}_X(U_i)$ が日本環であるものが存在する。
\item 開被覆 $X = \bigcup X_j$ であって、各開部分スキーム $X_j$ が日本スキームであるものが存在する。
\end{enumerate}
さらに、$X$ が日本スキームならば、すべての開部分スキームも日本スキームである。
\end{lemma}

\begin{proof}
これは補題 \ref{properties-lemma-locally-P} と代数の章の補題 \ref{algebra-lemma-localize-N} および
\ref{algebra-lemma-Japanese-local} から従う。
\end{proof}

\begin{lemma}
\label{properties-lemma-locally-universally-Japanese}
$X$ をスキームとする。次の条件は同値である：
\begin{enumerate}
\item スキーム $X$ は普遍的日本スキームである。
\item 任意のアフィン開集合 $U \subset X$ に対して、環 $\mathcal{O}_X(U)$ は普遍的日本環である。
\item アフィン開被覆 $X = \bigcup U_i$ であって、各 $\mathcal{O}_X(U_i)$ が普遍的日本環であるものが存在する。
\item 開被覆 $X = \bigcup X_j$ であって、各開部分スキーム $X_j$ が普遍的日本スキームであるものが存在する。
\end{enumerate}
さらに、$X$ が普遍的日本スキームならば、すべての開部分スキームも普遍的日本スキームである。
\end{lemma}

\begin{proof}
これは補題 \ref{properties-lemma-locally-P} と代数の章の補題 \ref{algebra-lemma-universally-japanese} および
\ref{algebra-lemma-nagata-local} から従う。
\end{proof}

\begin{lemma}
\label{properties-lemma-locally-nagata}
$X$ をスキームとする。次の条件は同値である：
\begin{enumerate}
\item スキーム $X$ は永田スキームである。
\item 任意のアフィン開集合 $U \subset X$ に対して、環 $\mathcal{O}_X(U)$ は永田環である。
\item アフィン開被覆 $X = \bigcup U_i$ であって、各 $\mathcal{O}_X(U_i)$ が永田環であるものが存在する。
\item 開被覆 $X = \bigcup X_j$ であって、各開部分スキーム $X_j$ が永田スキームであるものが存在する。
\end{enumerate}
さらに、$X$ が永田スキームならば、すべての開部分スキームも永田スキームである。
\end{lemma}

\begin{proof}
これは補題 \ref{properties-lemma-locally-P} と代数の章の補題 \ref{algebra-lemma-nagata-localize} および
\ref{algebra-lemma-nagata-local} から従う。
\end{proof}

\begin{lemma}
\label{properties-lemma-characterize-nagata}
$X$ を局所 Noether スキームとする。このとき $X$ が永田スキームであるための必要十分条件は、
すべての整閉部分スキーム $Z \subset X$ が日本スキームであることである。
\end{lemma}

\begin{proof}
$X$ が永田スキームであると仮定する。$Z \subset X$ を整閉部分スキームとし、$z \in Z$ とする。
$\Spec(A) = U \subset X$ を $z$ を含むアフィン開集合で、$A$ が永田環であるものとして取る。
このとき、$Z \cap U \cong \Spec(A/\mathfrak p)$ がある素イデアル $\mathfrak p$ に対して成り立つ。
スキームの章の補題 \ref{schemes-lemma-closed-subspace-scheme}（および定義
\ref{properties-definition-integral}）を参照せよ。代数の章の定義 \ref{algebra-definition-nagata} により、
$A/\mathfrak p$ は日本環である。従って、定義により $Z$ は日本スキームである。

\medskip\noindent
$X$ のすべての整閉部分スキームが日本スキームであると仮定する。
$\Spec(A) = U \subset X$ を任意のアフィン開集合とする。$X$ は局所 Noether なので、
$A$ は Noether 環である（補題 \ref{properties-lemma-locally-Noetherian}）。$\mathfrak p \subset A$ を素イデアルとする。
$A/\mathfrak p$ が日本環であることを示せばよい。$T \subset U$ を閉部分集合
$V(\mathfrak p) \subset \Spec(A)$ とする。$\overline{T} \subset X$ をその閉包とする。
すると $\overline{T}$ は既約部分集合の閉包なので既約である。従って、$\overline{T}$ により定まる
被約閉部分スキームは整閉部分スキームである（これも $\overline{T}$ と書く）。
スキームの章の補題 \ref{schemes-lemma-reduced-closed-subscheme} を参照せよ。
言い換えると、$\Spec(A/\mathfrak p)$ は $X$ のある整閉部分スキームのアフィン開集合である。
仮定によりこの部分スキームは日本スキームであり、補題 \ref{properties-lemma-locally-Japanese} により
$A/\mathfrak p$ は日本環である。
\end{proof}

\begin{lemma}
\label{properties-lemma-nagata-universally-Japanese}
$X$ をスキームとする。次の条件は同値である：
\begin{enumerate}
\item $X$ は永田スキームである。
\item $X$ は局所 Noether かつ普遍的日本スキームである。
\end{enumerate}
\end{lemma}

\begin{proof}
これは代数の章の命題 \ref{algebra-proposition-nagata-universally-japanese} である。
\end{proof}

\noindent
この議論は射の章の節 \ref{morphisms-section-nagata} で続ける。



\section{G-スキーム}
\label{properties-section-G-scheme}

\noindent
環 $R$ が G-環であるとは、$R$ が Noether 環であり、任意の素イデアル
$\mathfrak p$（$R$ のもの）に対して環準同型 $R_\mathfrak p \to (R_\mathfrak p)^\wedge$ が
正則であることをいう、と注意しておく。

\begin{definition}
\label{properties-definition-G-scheme}
$X$ をスキームとする。$X$ を {\it G-スキーム}\footnote{これは標準的でない用語かもしれない。
$G$ が有限群または群スキームであるとき、$G$-スキームはしばしば $G$ の作用を備えたスキームを
意味する。} と呼ぶのは、任意の $x \in X$ に対して、アフィン開近傍
$x \in U \subset X$ であって、その環 $\mathcal{O}_X(U)$ が G-環となるものが存在するときである
（発展代数の章の定義
\ref{more-algebra-definition-G-ring} を参照）。
\end{definition}

\begin{lemma}
\label{properties-lemma-characterize-G-scheme}
$X$ をスキームとする。次の条件は同値である：
\begin{enumerate}
\item $X$ は G-スキームである。
\item $X$ は局所 Noether であり、すべての $x \in X$ に対して環準同型
$\mathcal{O}_{X, x} \to \mathcal{O}_{X, x}^\wedge$ は正則である。
\item $X$ は局所 Noether であり、任意の閉点 $x \in X$ に対して環準同型
$\mathcal{O}_{X, x} \to \mathcal{O}_{X, x}^\wedge$ は正則である。
\end{enumerate}
\end{lemma}

\begin{proof}
(1) が (2) を含意し、(2) が (3) を含意することは明らかである。
(3) を仮定し、$x \in X$ を任意の点とする。補題
\ref{properties-lemma-locally-Noetherian-closed-point} により、特殊化 $x \leadsto x'$ であって、$x'$ が
$X$ の閉点となるものが存在する。このとき $\mathcal{O}_{X, x'}$ は、発展代数の章の補題
\ref{more-algebra-lemma-check-G-ring-maximal-ideals} により G-環である。
$\mathcal{O}_{X, x}$ は $\mathcal{O}_{X, x'}$ のある素イデアルにおける局所化なので
（スキームの章の補題 \ref{schemes-lemma-specialize-points}）、
$\mathcal{O}_{X, x} \to \mathcal{O}_{X, x}^\wedge$ は正則環準同型である。
$x \in X$ は任意であったから、任意のアフィン開集合 $U \subset X$ に対して、Noether 環
（補題 \ref{properties-lemma-locally-Noetherian}）$\mathcal{O}_X(U)$ は G-環である。従って (1) が成り立つ。
\end{proof}

\begin{lemma}
\label{properties-lemma-locally-G-scheme}
$X$ をスキームとする。次の条件は同値である：
\begin{enumerate}
\item スキーム $X$ は G-スキームである。
\item 任意のアフィン開集合 $U \subset X$ に対して、環 $\mathcal{O}_X(U)$ は G-環である。
\item アフィン開被覆 $X = \bigcup U_i$ であって、各 $\mathcal{O}_X(U_i)$ が G-環であるものが存在する。
\item 開被覆 $X = \bigcup X_j$ であって、各開部分スキーム $X_j$ が G-スキームであるものが存在する。
\end{enumerate}
さらに、$X$ が G-スキームならば、すべての開部分スキームも G-スキームである。
\end{lemma}

\begin{proof}
補題 \ref{properties-lemma-locally-Noetherian} と \ref{properties-lemma-characterize-G-scheme} を組み合わせよ。
\end{proof}




\section{特異点集合}
\label{properties-section-singular-locus}

\noindent
定義は次のとおりである。

\begin{definition}
\label{properties-definition-singular-locus}
$X$ を局所 Noether スキームとする。{\it 正則点集合} $\text{Reg}(X)$（$X$ のもの）とは、
点 $x \in X$ であって $\mathcal{O}_{X, x}$ が正則局所環となるものの集合である。{\it 特異点集合}
$\text{Sing}(X)$ とは補集合 $X \setminus \text{Reg}(X)$、すなわち
点 $x \in X$ であって $\mathcal{O}_{X, x}$ が正則局所環でないものの集合である。
\end{definition}

\noindent
局所 Noether スキームの正則点集合は一般化の下で安定である。代数の章の定義
\ref{algebra-definition-regular} の前の議論を参照せよ。しかし一般の局所 Noether スキームでは、
正則点集合が開であるとは限らない。これが成り立つためのいくつかの判定条件は、発展代数の章の節
\ref{more-algebra-section-singular-locus} にある。射の章の節
\ref{morphisms-section-singular-locus} でさらに論じる。

\section{局所既約性}
\label{properties-section-unibranch}

\noindent
発展代数の章の節 \ref{more-algebra-section-unibranch} で、（幾何的）単枝局所環の概念を
導入したことを思い出そう。

\begin{definition}
\label{properties-definition-unibranch}
\begin{reference}
\cite[Chapter IV (6.15.1)]{EGA4}
\end{reference}
$X$ をスキームとし、$x \in X$ とする。$X$ が {\it $x$ において単枝} であるとは、
局所環 $\mathcal{O}_{X, x}$ が単枝であることをいう。$X$ が
{\it $x$ において幾何的単枝} であるとは、局所環 $\mathcal{O}_{X, x}$ が幾何的単枝で
あることをいう。$X$ がそのすべての点で
単枝であるとき、$X$ は {\it 単枝} であるという。$X$ がそのすべての点で幾何的単枝で
あるとき、$X$ は {\it 幾何的単枝} であるという。
\end{definition}

\noindent
念のため述べると、局所環 $A$ が（発展代数の章の定義
\ref{more-algebra-definition-unibranch} の意味で）幾何的単枝であっても、スキーム
$\Spec(A)$ が定義 \ref{properties-definition-unibranch} の意味で幾何的単枝でないことはありうる。
例えば、通常二重点特異点（節点）を持つ既約平面曲線上の錐の頂点における局所環を $A$ と
すれば、このことが起こる。

\begin{lemma}
\label{properties-lemma-normal-geometrically-unibranch}
正規スキームは幾何的単枝である。
\end{lemma}

\begin{proof}
これは定義から従う。実際、スキームが正規であるとは、その局所環が正規整域であることをいう。
局所正規整域が幾何的単枝であることは、発展代数の章の定義
\ref{more-algebra-definition-unibranch} から直ちに分かる。
\end{proof}

\begin{lemma}
\label{properties-lemma-geometrically-unibranch}
\begin{reference}
\cite[Proposition 2.3]{Etale-coverings} と比較せよ。
\end{reference}
$X$ を Noether スキームとする。次は同値である。
\begin{enumerate}
\item $X$ は幾何的単枝である（定義 \ref{properties-definition-unibranch}）。
\item 任意の点 $x \in X$ であって $X$ の既約成分の生成点ではないものに対して、強 Hensel 化
$\mathcal{O}_{X, x}^{sh}$ の閉点を除いたスペクトルは連結である。
\end{enumerate}
\end{lemma}

\begin{proof}
発展代数の章の補題 \ref{more-algebra-lemma-geometrically-unibranch} により、(1) ならば
(2) の閉点を除いたスペクトルは既約であり、特に連結である。

\medskip\noindent
(2) を仮定する。$x \in X$ とする。$\mathcal{O}_{X, x}$ が幾何的単枝であることを
示さなければならない。$\dim(\mathcal{O}_{X, x})$ に関する帰納法により、$x$ の任意の
非自明な一般化について結論が成り立つと仮定してよい。$X$ を
$\Spec(\mathcal{O}_{X, x})$ で置き換えてよい。言い換えると、$X = \Spec(A)$ であり
$A$ は局所環、さらに $A_\mathfrak p$ は各非極大素イデアル
$\mathfrak p \subset A$ に対して幾何的単枝であると仮定してよい。

\medskip\noindent
$A^{sh}$ を $A$ の強 Hensel 化とする。$\mathfrak q \subset A^{sh}$ が
$\mathfrak p \subset A$ の上にある素イデアルならば、$A_\mathfrak p \to A^{sh}_\mathfrak q$
はエタール代数のフィルター余極限である。したがって、$A_\mathfrak p$ と
$A^{sh}_\mathfrak q$ の強 Hensel 化は同型である。よって発展代数の章の補題
\ref{more-algebra-lemma-geometrically-unibranch} から、局所環 $A^{sh}_\mathfrak q$ は、
非極大素イデアル $\mathfrak q$ が $A^{sh}$ を走るとき、常に唯一の最小素イデアルを持つと
結論できる。

\medskip\noindent
$\mathfrak q_1, \ldots, \mathfrak q_r$ を $A^{sh}$ の最小素イデアルとする。$r = 1$ を
示せばよい。上の議論から、$V(\mathfrak q_1) \cap V(\mathfrak q_j) = \{\mathfrak m^{sh}\}$
であり、ここで $j = 2, \ldots, r$ である。したがって
$V(\mathfrak q_1) \setminus \{\mathfrak m^{sh}\}$ は $A^{sh}$ の閉点を除いたスペクトルの
開かつ閉な部分集合である。このスペクトルが連結であるという仮定に反する。ただし
$r = 1$ の場合を除く。
\end{proof}

\begin{definition}
\label{properties-definition-number-of-branches}
$X$ をスキームとし、$x \in X$ とする。{\it $X$ の $x$ における枝数} とは、発展代数の章の
定義 \ref{more-algebra-definition-number-of-branches} で定めた局所環
$\mathcal{O}_{X, x}$ の枝数である。{\it $X$ の $x$ における幾何的枝数} とは、発展代数の章の
定義 \ref{more-algebra-definition-number-of-branches} で定めた局所環
$\mathcal{O}_{X, x}$ の幾何的枝数である。
\end{definition}

\noindent
これを完備局所環の枝と比較したいことが多いが、一般には比較は単純ではない。この話題については、
発展代数の章の節 \ref{more-algebra-section-branches-completion} にいくらかの情報がある。

\begin{lemma}
\label{properties-lemma-number-of-branches-irreducible-components}
$X$ をスキーム、$x \in X$ とする。$X_i$（$i \in I$）を、$X$ の $x$ を通る既約成分とする。
$X$ の $x$ における（幾何的）枝数は、$i \in I$ にわたる、$X_i$ の $x$ における
（幾何的）枝数の和である。
\end{lemma}

\begin{proof}
$X_i$ を $X$ の整閉部分スキームとみなす。スキームの章の定義
\ref{schemes-definition-reduced-induced-scheme} および補題
\ref{properties-lemma-characterize-integral} を参照せよ。$X_i$ の $x$ における（幾何的）枝数は
少なくとも $1$ である（本質的には定義による）。これはすべての $i$ について成り立つ。
$X_i$ は $1$ 対 $1$ に、その局所環の最小素イデアル
$\mathfrak p_i \subset \mathcal{O}_{X, x}$ と対応することを思い出そう。代数の章の補題
\ref{algebra-lemma-irreducible-components-containing-x} を参照せよ。したがって $I$ が無限ならば、
$\mathcal{O}_{X, x}$ は無限個の最小素イデアルを持つ。それゆえ
$\mathcal{O}_{X, x}^h$ と $\mathcal{O}_{X, x}^{sh}$ の双方も無限個の最小素イデアルを持つ
（代数の章の補題 \ref{algebra-lemma-injective-minimal-primes-in-image} と
\ref{algebra-lemma-minimal-prime-image-minimal-prime}、および写像
$\mathcal{O}_{X, x} \to  \mathcal{O}_{X, x}^h \to \mathcal{O}_{X, x}^{sh}$ の単射性を
組み合わせよ）。この場合、$X$ の $x$ における（幾何的）枝数は $\infty$ と定義され、
和についても同じである。したがって $I$ は有限と仮定してよい。
$A'$ を、$\mathcal{O}_{X, x}$ の、全商環 $Q$（$(\mathcal{O}_{X, x})_{red}$ のもの）における
整閉包とする。$A'_i$ を、$\mathcal{O}_{X, x}/\mathfrak p_i$ の、全商環 $Q_i$
（$\mathcal{O}_{X, x}/\mathfrak p_i$ のもの）における整閉包とする。代数の章の補題
\ref{algebra-lemma-total-ring-fractions-no-embedded-points} により
$Q = \prod_{i \in I} Q_i$ である。したがって $A' = \prod A'_i$ である。これで、
（幾何的）枝数を $A'$ の極大イデアルによって表す発展代数の章の補題
\ref{more-algebra-lemma-number-of-branches-1} から、補題の等式が従う。
\end{proof}

\begin{lemma}
\label{properties-lemma-number-of-branches-1}
$X$ をスキームとし、$x \in X$ とする。
\begin{enumerate}
\item $X$ の $x$ における枝数が $1$ であるための必要十分条件は、$X$ が $x$ において
単枝であることである。
\item $X$ の $x$ における幾何的枝数が $1$ であるための必要十分条件は、$X$ が $x$ において
幾何的単枝であることである。
\end{enumerate}
\end{lemma}

\begin{proof}
この補題は、定義と環についての対応する結果から直ちに従う。発展代数の章の補題
\ref{more-algebra-lemma-number-of-branches-1} を参照せよ。
\end{proof}






\section{有限型加群と有限表示加群の特徴づけ}
\label{properties-section-characterizing-finite-type-presentation}

\noindent
$X$ をスキームとする。$\mathcal{F}$ を準連接 $\mathcal{O}_X$-加群とする。
次の補題によれば、$\mathcal{F}$ が有限型であること
（加群の章の定義 \ref{modules-definition-finite-type} を参照）と、$\mathcal{F}$ が
各アフィン開集合 $\Spec(A) = U \subset X$ 上で $\widetilde M$ の形をし、ここで
$A$-加群 $M$ が有限生成であることとは同値である。同様に、$\mathcal{F}$ が有限表示で
あること（加群の章の定義 \ref{modules-definition-finite-presentation} を参照）と、
$\mathcal{F}$ が各アフィン開集合 $\Spec(A) = U \subset X$ 上で $\widetilde M$ の形をし、
ここで $A$-加群 $M$ が有限表示であることとは同値である。

\begin{lemma}
\label{properties-lemma-finite-type-module}
$X = \Spec(R)$ をアフィンスキームとする。$\mathcal{O}_X$-加群の準連接層
$\widetilde M$ が有限型 $\mathcal{O}_X$-加群であるための必要十分条件は、$M$ が有限生成
$R$-加群であることである。
\end{lemma}

\begin{proof}
$\widetilde M$ が有限型 $\mathcal{O}_X$-加群であると仮定する。これは、$X$ のある開被覆が
存在し、その各要素への $\widetilde M$ の制限が有限個の切断により大域生成されることを
意味する。したがって、標準開被覆
$X = \bigcup_{i = 1, \ldots, n} D(f_i)$ であって、$\widetilde M|_{D(f_i)}$ が有限個の
切断により生成されるものも存在する。従って $M_{f_i}$ は各 $i$ に対して有限生成である。
ゆえに代数の章の補題 \ref{algebra-lemma-cover} から結論が従う。
\end{proof}

\begin{lemma}
\label{properties-lemma-finite-presentation-module}
$X = \Spec(R)$ をアフィンスキームとする。$\mathcal{O}_X$-加群の準連接層
$\widetilde M$ が有限表示 $\mathcal{O}_X$-加群であるための必要十分条件は、$M$ が有限表示
$R$-加群であることである。
\end{lemma}

\begin{proof}
$\widetilde M$ が有限表示 $\mathcal{O}_X$-加群であると仮定する。補題
\ref{properties-lemma-finite-type-module} により、$M$ は有限生成 $R$-加群である。全射
$R^n \to M$ を選び、その核を $K$ とする。スキームの章の補題
\ref{schemes-lemma-spec-sheaves} により、短完全列
$$
0 \to \widetilde{K} \to
\bigoplus \mathcal{O}_X^{\oplus n} \to
\widetilde{M} \to 0
$$
がある。加群の章の補題
\ref{modules-lemma-kernel-surjection-finite-free-onto-finite-presentation}
により、$\widetilde{K}$ は有限型 $\mathcal{O}_X$-加群である。したがって補題
\ref{properties-lemma-finite-type-module} を再び用いると、$K$ は有限生成 $R$-加群である。
従って $M$ は有限表示 $R$-加群である。
\end{proof}







\section{主開集合上の切断}
\label{properties-section-principal-opens}

\noindent
これはこの種の典型的な結果である。後の補題では、より素朴だが直接的な証明法を用いる。

\begin{lemma}
\label{properties-lemma-invert-f-sections}
\begin{slogan}
準連接層の切断は無限遠で有理型特異点しか持たない。
\end{slogan}
$X$ をスキームとする。$f \in \Gamma(X, \mathcal{O}_X)$ とする。
$X_f \subset X$ を、$f$ が可逆となる開集合と書く。スキームの章の補題
\ref{schemes-lemma-f-open} を参照せよ。
$X$ が準コンパクトかつ準分離的ならば、標準写像
$$
\Gamma(X, \mathcal{O}_X)_f \longrightarrow \Gamma(X_f, \mathcal{O}_X)
$$
は同型である。さらに、$\mathcal{F}$ が準連接な $\mathcal{O}_X$-加群の層ならば、写像
$$
\Gamma(X, \mathcal{F})_f \longrightarrow \Gamma(X_f, \mathcal{F})
$$
は同型である。
\end{lemma}

\begin{proof}
$R = \Gamma(X, \mathcal{O}_X)$ と書く。スキームの標準射
$$
\varphi : X \longrightarrow \Spec(R)
$$
を考える。スキームの章の補題
\ref{schemes-lemma-morphism-into-affine} を参照せよ。右辺の標準開集合 $D(f)$ の逆像は、
左辺の $X_f$ である。さらに $X$ は準コンパクトかつ準分離的と仮定したので、射 $\varphi$ も
準コンパクトかつ準分離的である。スキームの章の補題
\ref{schemes-lemma-quasi-compact-affine} および
\ref{schemes-lemma-compose-after-separated} を参照せよ。従ってスキームの章の補題
\ref{schemes-lemma-push-forward-quasi-coherent} により、$\varphi_*\mathcal{F}$ は準連接である。
ゆえに、$\varphi_*\mathcal{F} = \widetilde M$ であり、
$M = \Gamma(X, \mathcal{F})$ は $R$-加群である。従って
$$
\Gamma(X_f, \mathcal{F}) =
\Gamma(D(f), \varphi_*\mathcal{F}) =
\Gamma(D(f), \widetilde M) = M_f
$$
となり、これはまさに補題の内容である。補題の最初の表示同型は
$\mathcal{F} = \mathcal{O}_X$ と取れば従う。
\end{proof}

\noindent
$X$ をスキームとし、$\mathcal{L}$ を $X$ 上の可逆層とする。
$\mathcal{O}_X$-加群の層 $\mathcal{F}$ に対して、このとき次数付き環
$\Gamma_*(X, \mathcal{L}) =
\bigoplus\nolimits_{n \geq 0} \Gamma(X, \mathcal{L}^{\otimes n})$
と、次数付き $\Gamma_*(X, \mathcal{L})$-加群
$\Gamma_*(X, \mathcal{L}, \mathcal{F}) =
\bigoplus\nolimits_{n \in \mathbf{Z}}
\Gamma(X, \mathcal{F} \otimes_{\mathcal{O}_X} \mathcal{L}^{\otimes n})$
を得る。加群の章の定義 \ref{modules-definition-gamma-star} を参照せよ。さらに切断
$s \in \Gamma(X, \mathcal{L})$ が与えられると、写像
\begin{equation}
\label{properties-equation-module-invert-s}
\Gamma_*(X, \mathcal{L}, \mathcal{F})_{(s)}
\longrightarrow
\Gamma(X_s, \mathcal{F}|_{X_s})
\end{equation}
を得る。これは $t/s^n$（ここで
$t \in \Gamma(X, \mathcal{F} \otimes_{\mathcal{O}_X} \mathcal{L}^{\otimes n})$）を
$t|_{X_s} \otimes s|_{X_s}^{-n}$ に送る。定義により $X_s \subset X$ は $s$ が逆元を持つ
開集合なので、これは意味を持つ。加群の章の補題 \ref{modules-lemma-s-open} を参照せよ。

\begin{lemma}
\label{properties-lemma-invert-s-sections}
\begin{reference}
\cite[Section 6.8]{EGA1-second}
\end{reference}
$X$ をスキームとする。$\mathcal{L}$ を $X$ 上の可逆層とする。
$s \in \Gamma(X, \mathcal{L})$ とする。$\mathcal{F}$ を準連接
$\mathcal{O}_X$-加群とする。
\begin{enumerate}
\item $X$ が準コンパクトならば、(\ref{properties-equation-module-invert-s}) は単射である。
\item $X$ が準コンパクトかつ準分離的ならば、
(\ref{properties-equation-module-invert-s}) は同型である。
\end{enumerate}
特に、標準写像
$$
\Gamma_*(X, \mathcal{L})_{(s)}
\longrightarrow
\Gamma(X_s, \mathcal{O}_X),\quad
a/s^n \longmapsto a \otimes s^{-n}
$$
は、$X$ が準コンパクトかつ準分離的ならば同型である。
\end{lemma}

\begin{proof}
$X$ が準コンパクトであると仮定する。有限アフィン開被覆
$X = U_1 \cup \ldots \cup U_m$ であって、$U_j$ はアフィンであり
$\mathcal{L}|_{U_j} \cong \mathcal{O}_{U_j}$ となるものを選ぶ。この同型を通じて、像
$s|_{U_j}$ はある $f_j \in \Gamma(U_j, \mathcal{O}_{U_j})$ に対応する。このとき
$X_s \cap U_j = D(f_j)$ である。

\medskip\noindent
(1) を証明する。(\ref{properties-equation-module-invert-s}) の核の元を $t/s^n$ とする。このとき
$t|_{X_s} = 0$ である。従って $(t|_{U_j})|_{D(f_j)} = 0$ である。補題
\ref{properties-lemma-invert-f-sections} により、$f_j^{e_j} t|_{U_j} = 0$ となるある
$e_j \geq 0$ が存在する。$e = \max(e_j)$ と置く。このとき $t \otimes s^e$ は
すべての $U_j$ 上（すべての $j$ について）で零に制限されるので、零である。$t/s^n$ は
$t \otimes s^e/s^{n + e}$ に等しく、これは
$\Gamma_*(X, \mathcal{L}, \mathcal{F})_{(s)}$ における等式である。従って、望む通り $t/s^n = 0$ を得る。

\medskip\noindent
(2) を証明する。$X$ が準コンパクトかつ準分離的であると仮定する。このとき
$U_j \cap U_{j'}$ はすべての組 $j, j'$ に対して準コンパクトである。スキームの章の補題
\ref{schemes-lemma-characterize-quasi-separated} を参照せよ。(1) により
(\ref{properties-equation-module-invert-s}) が単射であることは分かっている。
$t' \in \Gamma(X_s, \mathcal{F}|_{X_s})$ とする。各 $j$ に対して、整数
$e_j \geq 0$ と $t'_j \in \Gamma(U_j, \mathcal{F}|_{U_j})$ が存在し、補題
\ref{properties-lemma-invert-f-sections} の同型を通じて $t'|_{D(f_j)}$ は
$t'_j/f_j^{e_j}$ に対応する。$e = \max(e_j)$ と置き、
$$
t_j = f_j^{e - e_j} t'_j \otimes q_j^e \in
\Gamma(U_j,
(\mathcal{F} \otimes_{\mathcal{O}_X} \mathcal{L}^{\otimes e})|_{U_j})
$$
とする。ここで $q_j \in \Gamma(U_j, \mathcal{L}|_{U_j})$ は同型
$\mathcal{L}|_{U_j} \cong \mathcal{O}_{U_j}$ から来る自明化切断である。特に
$s|_{U_j} = f_j q_j$ である。これを用いて計算すると、
$t_j|_{U_j \cap U_{j'}}$ と $t_{j'}|_{U_j \cap U_{j'}}$ は、
$\mathcal{F}$ の $U_j \cap U_{j'} \cap X_s$ 上の同じ切断に写る。
$U_j \cap U_{j'}$ の準コンパクト性と (1) により、整数 $e' \geq 0$ が存在して
$$
t_j|_{U_j \cap U_{j'}} \otimes s^{e'}|_{U_j \cap U_{j'}} =
t_{j'}|_{U_j \cap U_{j'}} \otimes s^{e'}|_{U_j \cap U_{j'}}
$$
が $\mathcal{F} \otimes \mathcal{L}^{\otimes e + e'}$ の $U_j \cap U_{j'}$ 上の切断として
成り立つ。同じ $e'$ がすべての組 $j, j'$ に対して働くように選べる。すると層条件により、
切断 $t \in \Gamma(X, \mathcal{F} \otimes \mathcal{L}^{\otimes e + e'})$ であって、
$U_j$ への制限が $t_j \otimes s^{e'}|_{U_j}$ となるものが存在する。簡単な計算により、
$t/s^{e + e'}$ は望む通り $t'$ に写る。
\end{proof}

\noindent
$X$ をスキームとする。$\mathcal{L}$ を可逆 $\mathcal{O}_X$-加群とする。
$\mathcal{F}$ と $\mathcal{G}$ を準連接 $\mathcal{O}_X$-加群とする。次数付き
$\Gamma_*(X, \mathcal{L})$-加群
$$
M = \bigoplus\nolimits_{n \in \mathbf{Z}} \Hom_{\mathcal{O}_X}(\mathcal{F},
\mathcal{G} \otimes_{\mathcal{O}_X} \mathcal{L}^{\otimes n})
$$
を考える。次に、切断 $s \in \Gamma(X, \mathcal{L})$ を取る。このとき標準写像
\begin{equation}
\label{properties-equation-hom-invert-s}
M_{(s)} \longrightarrow
\Hom_{\mathcal{O}_{X_s}}(\mathcal{F}|_{X_s}, \mathcal{G}|_{X_s})
\end{equation}
があり、$\alpha/s^n$ を写像 $\alpha|_{X_s} \otimes s|_{X_s}^{-n}$ に送る。次の補題と補題
\ref{properties-lemma-extend-finite-presentation} を合わせると、大まかに言って、$X$ が準コンパクトかつ
準分離的ならば、有限表示 $\mathcal{O}_{X_s}$-加群の圏は、有限表示
$\mathcal{O}_X$-加群の圏において写像の乗法系
$s^n: \mathcal{F} \to
\mathcal{F} \otimes_{\mathcal{O}_X} \mathcal{L}^{\otimes n}$ を可逆にした圏である。

\begin{lemma}
\label{properties-lemma-section-maps-backwards}
$X$ をスキームとする。$\mathcal{L}$ を可逆 $\mathcal{O}_X$-加群とする。
$s \in \Gamma(X, \mathcal{L})$ を切断とする。
$\mathcal{F}$、$\mathcal{G}$ を準連接 $\mathcal{O}_X$-加群とする。
\begin{enumerate}
\item $X$ が準コンパクトで $\mathcal{F}$ が有限型ならば、
(\ref{properties-equation-hom-invert-s}) は単射である。
\item $X$ が準コンパクトかつ準分離的で $\mathcal{F}$ が有限表示ならば、
(\ref{properties-equation-hom-invert-s}) は全単射である。
\end{enumerate}
\end{lemma}

\begin{proof}
まず $X = \Spec(A)$ がアフィンで、$\mathcal{L} = \mathcal{O}_X$ の場合に補題を証明する。
この場合、$s$ は元 $f \in A$ に対応する。$\mathcal{F} = \widetilde{M}$ および
$\mathcal{G} = \widetilde{N}$ と書く（ここで $A$-加群 $M$ と $N$ を取る）。このとき補題は
（補題 \ref{properties-lemma-finite-type-module} および
\ref{properties-lemma-finite-presentation-module} を通じて）次の代数的主張に翻訳される。
\begin{enumerate}
\item $M$ が有限生成 $A$-加群で、$\varphi : M \to N$ が $A$-加群写像であり、誘導写像
$M_f \to N_f$ が零ならば、$f^n\varphi = 0$ となる $n$ が存在する。
\item $M$ が有限表示 $A$-加群ならば、
$\Hom_A(M, N)_f = \Hom_{A_f}(M_f, N_f)$ である。
\end{enumerate}
第二の主張は代数の章の補題 \ref{algebra-lemma-hom-from-finitely-presented} である。第一の主張の
証明は省略する。

\medskip\noindent
次に一般の $X$ に対して (1) を証明する。$X$ が準コンパクトであると仮定し、有限アフィン開被覆
$X = U_1 \cup \ldots \cup U_m$ であって、$U_j$ はアフィンかつ
$\mathcal{L}|_{U_j} \cong \mathcal{O}_{U_j}$ となるものを選ぶ。この同型を通じて、像
$s|_{U_j}$ はある $f_j \in \Gamma(U_j, \mathcal{O}_{U_j})$ に対応する。このとき
$X_s \cap U_j = D(f_j)$ である。(\ref{properties-equation-hom-invert-s}) の核の元を
$\alpha/s^n$ とする。このとき $\alpha|_{X_s} = 0$ である。従って
$(\alpha|_{U_j})|_{D(f_j)} = 0$ である。上で扱ったアフィンの場合により、
$f_j^{e_j} \alpha|_{U_j} = 0$ となるある $e_j \geq 0$ が存在する。$e = \max(e_j)$ と置く。
すると $\alpha \otimes s^e$ はすべての $U_j$ 上（すべての $j$ に対して）で零に制限されるので、零である。
$\alpha/s^n$ は $\alpha \otimes s^e/s^{n + e}$ と $M_{(s)}$ において等しいので、
望む通り $\alpha/s^n = 0$ を得る。

\medskip\noindent
(2) を証明する。$\mathcal{F}$ は有限表示なので、層
$\SheafHom_{\mathcal{O}_X}(\mathcal{F}, \mathcal{G})$ は準連接である。スキームの章の節
\ref{schemes-section-quasi-coherent} を参照せよ。さらに、すべての $n$ に対して明らかに
$$
\SheafHom_{\mathcal{O}_X}(\mathcal{F},
\mathcal{G} \otimes_{\mathcal{O}_X} \mathcal{L}^{\otimes n}) =
\SheafHom_{\mathcal{O}_X}(\mathcal{F}, \mathcal{G})
\otimes_{\mathcal{O}_X} \mathcal{L}^{\otimes n}
$$
である。従ってこの場合、$\SheafHom_{\mathcal{O}_X}(\mathcal{F}, \mathcal{G})$ に補題
\ref{properties-lemma-invert-s-sections} を適用すれば主張が従う。
\end{proof}





\section{準アフィンスキーム}
\label{properties-section-quasi-affine}

\begin{definition}
\label{properties-definition-quasi-affine}
スキーム $X$ が {\it 準アフィン} であるとは、準コンパクトであり、かつアフィンスキームの開部分スキームと同型であることをいう。
\end{definition}

\begin{lemma}
\label{properties-lemma-quasi-coherent-quasi-affine}
環 $A$ を取り、$U \subset \Spec(A)$ を準コンパクトな開部分スキームとする。
準連接層 $\mathcal{F}$ を $U$ 上に取ると、標準写像
$$
\widetilde{H^0(U, \mathcal{F})}|_U \to \mathcal{F}
$$
は同型である。
\end{lemma}

\begin{proof}
$j : U \to \Spec(A)$ を包含射と記す。このとき
$H^0(U, \mathcal{F}) = H^0(\Spec(A), j_*\mathcal{F})$ であり、
$j_*\mathcal{F}$ は、スキームの章の補題
\ref{schemes-lemma-push-forward-quasi-coherent} により準連接である。
従って、スキームの章の補題
\ref{schemes-lemma-equivalence-quasi-coherent} により
$j_*\mathcal{F} = \widetilde{H^0(U, \mathcal{F})}$ である。
$U$ に制限すれば補題を得る。
\end{proof}

\begin{lemma}
\label{properties-lemma-invert-f-affine}
$X$ をスキームとする。$f \in \Gamma(X, \mathcal{O}_X)$ とする。
$X$ が準コンパクトかつ準分離的であり、$X_f$ がアフィンであると仮定する。
スキームの章の補題
\ref{schemes-lemma-morphism-into-affine} にある標準射
$$
j : X \longrightarrow \Spec(\Gamma(X, \mathcal{O}_X))
$$
は、$X_f = j^{-1}(D(f))$ を標準アフィン開集合
$D(f) \subset \Spec(\Gamma(X, \mathcal{O}_X))$ へ同型に写す。
\end{lemma}

\begin{proof}
これは、上の補題 \ref{properties-lemma-invert-f-sections} により、$j$ が環の同型
$\Gamma(X, \mathcal{O}_X)_f \to \mathcal{O}_X(X_f)$ を誘導することから明らかである。
\end{proof}

\begin{lemma}
\label{properties-lemma-quasi-affine}
$X$ をスキームとする。このとき $X$ が準アフィンであることと、
スキームの章の補題 \ref{schemes-lemma-morphism-into-affine} にある標準射
$$
X \longrightarrow \Spec(\Gamma(X, \mathcal{O}_X))
$$
が準コンパクトな開埋め込みであることは同値である。
\end{lemma}

\begin{proof}
表示された射が準コンパクトな開埋め込みならば、$X$ は
$\Spec(\Gamma(X, \mathcal{O}_X))$ の準コンパクトな開部分スキームと同型であり、
従って明らかに $X$ は準アフィンである。

\medskip\noindent
$X$ が準アフィンであり、$X \subset \Spec(R)$ が準コンパクトな開集合であると仮定する。
特に $X$ は分離的である。これはスキームの章の補題
\ref{schemes-lemma-subscheme-of-separated-scheme} による。
$A = \Gamma(X, \mathcal{O}_X)$ と置く。
環写像 $R \to A$ を、$R = \Gamma(\Spec(R), \mathcal{O}_{\Spec(R)})$ と
$\mathcal{O}_{\Spec(R)}$ の制限写像から考える。
スキームの章の補題 \ref{schemes-lemma-morphism-into-affine} により、
包含射の因子分解
$$
X \longrightarrow
\Spec(A) \longrightarrow
\Spec(R)
$$
を得る。$x \in X$ とする。$r \in R$ を選び、$x \in D(r)$ かつ
$D(r) \subset X$ とする。$f \in A$ と記す。これは $r$ の $A$ における像である。
上の補題 \ref{properties-lemma-invert-f-sections} の開集合 $X_f$ は
$D(r) \subset X$ に等しく、従って同補題の結論により $A_f \cong R_r$ である。
従って $D(r) \to \Spec(A)$ は標準アフィン開集合
$D(f)$ の $\Spec(A)$ への同型である。$X$ はこのようなアフィン開集合 $D(f)$ で覆えるので、
証明が完了する。
\end{proof}

\begin{lemma}
\label{properties-lemma-cartesian-diagram-quasi-affine}
準アフィンスキームの開埋め込み $U \to V$ を取る。このとき
$$
\xymatrix{
U \ar[d] \ar[rr]_-j & & \Spec(\Gamma(U, \mathcal{O}_U)) \ar[d] \\
U \ar[r] & V \ar[r]^-{j'} & \Spec(\Gamma(V, \mathcal{O}_V))
}
$$
はカルテジアンである。
\end{lemma}

\begin{proof}
この図式は、スキームの章の補題
\ref{schemes-lemma-morphism-into-affine} により可換である。
$A = \Gamma(U, \mathcal{O}_U)$ および $B = \Gamma(V, \mathcal{O}_V)$ と書く。
$g \in B$ を、$V_g$ がアフィンで $U$ に含まれるように取る。
これは、$f$ を $g$ の $A$ における像とすれば $U_f = V_g$ であることを意味する。
補題 \ref{properties-lemma-invert-f-affine} により、$j'$ は $V_g$ を
標準開集合 $D(g)$ の $\Spec(B)$ への同型に写す。従って
$V_g \times_{\Spec(B)} \Spec(A) \to \Spec(A)$ は
$D(f) \subset \Spec(A)$ への同型である。再び補題
\ref{properties-lemma-invert-f-affine} により、$j$ は $U_f$ を $D(f)$ へ同型に写す。
従って
$U_f = U_f \times_{\Spec(B)} \Spec(A)$ である。補題
\ref{properties-lemma-quasi-affine} により、$U$ を上記のような $V_g = U_f$ で覆えるので、
$U \to U \times_{\Spec(B)} \Spec(A)$ は同型である。
\end{proof}

\begin{lemma}
\label{properties-lemma-quasi-affine-presentation}
$X$ を準アフィンスキームとする。ある整数 $n \geq 0$、アフィンスキーム $T$、
および射 $T \to X$ が存在し、任意の射 $X' \to X$ で $X'$ がアフィンなものに対して、
ファイバー積 $X' \times_X T$ は $\mathbf{A}^n_{X'}$ と $X'$ 上で同型である。
\end{lemma}

\begin{proof}
定義により、ある環 $A$ が存在して、$X$ は準コンパクトな開部分スキーム
$U \subset \Spec(A)$ と同型である。
標準開集合 $D(f) \subset \Spec(A)$ が位相の基底をなすことを思い出そう。
これは代数の章の節 \ref{algebra-section-spectrum-ring} による。
$U$ は準コンパクトなので、$f_1, \ldots, f_n \in A$ を選んで
$U = D(f_1) \cup \ldots \cup D(f_n)$ とできる。従って
$X = \Spec(A) \setminus V(I)$ と仮定してよい。ここで
$I = (f_1, \ldots, f_n)$ である。次を置く。
$$
T = \Spec(A[t, x_1, \ldots, x_n]/(f_1 x_1 + \ldots + f_n x_n - 1))
$$
$T \to \Spec(A)$ の構造射は開集合 $X$ を経由するので、射
$T \to X$ を与える。$X' = \Spec(A')$ とし、$X' \to X$ が環写像
$A \to A'$ に対応するとする。このとき
$f'_1, \ldots, f'_n \in A'$ は、$f_1, \ldots, f_n$ の像であり、
$A'$ の単位イデアルを生成する。
次のように書く。
$1 = f'_1 a'_1 + \ldots + f'_n a'_n$.
基底変換 $X' \times_X T$ は次の環のスペクトルである。
$A'[t, x_1, \ldots, x_n]/(f'_1 x_1 + \ldots + f'_n x_n - 1)$.
$A'$-代数準同型
$$
\varphi :
A'[y_1, \ldots, y_n]
\longrightarrow
A'[t, x_1, \ldots, x_n, x_{n + 1}]/(f'_1 x_1 + \ldots + f'_n x_n - 1)
$$
を考え、$y_i$ を $a'_i t + x_i$ に送る。この写像は同型であると主張する。
これで補題の証明は完了する。$\varphi$ の逆は次の $A'$-代数準同型で与えられる。
$$
\psi :
A'[t, x_1, \ldots, x_n, x_{n + 1}]/(f'_1 x_1 + \ldots + f'_n x_n - 1)
\longrightarrow
A'[y_1, \ldots, y_n]
$$
$t$ を $-1 + f'_1 y_1 + \ldots + f'_n y_n$ に送り、また
$x_i$ を $y_i + a'_i - a'_i(f'_1 y_1 + \ldots + f'_n y_n)$ に送る（$i = 1, \ldots, n$）。
これは、$\sum f'_ix_i$ が次へ写ることから意味を持つ。
$$
\begin{matrix}
\sum f'_i(y_i + a'_i - a'_i(\sum f'_j y_j)) =
(\sum f'_iy_i) + 1 - (\sum f'_j y_j) = 1
\end{matrix}
$$
二つの写像が互いに逆であることは、次の計算から分かる。
$$
\begin{matrix}
\varphi(\psi(t) = \varphi(-1 +  \sum f'_i y_i) =
-1 + \sum f'_i (a'_i t + x_i) = t \\
\varphi(\psi(x_i)) = \varphi(y_i + a'_i - a'_i(\sum f'_j y_j)) =
a'_i t + x_i + a'_i - a'_i(\sum f'_ja'_jt + f'_jx_j) = x_i \\
\psi(\varphi(y_i)) = \psi(a'_i t + x_i) =
a'_i(-1 + \sum f'_j y_j) + y_i + a'_i - a'_i(\sum f'_j y_j) = y_i
\end{matrix}
$$
証明が完了する。
\end{proof}

\section{平坦加群}
\label{properties-section-flat}

\noindent
任意の環付き空間 $(X, \mathcal{O}_X)$ について、
$\mathcal{O}_X$-加群が（点において）平坦であることの意味は既知である。これは
加群の章の定義 \ref{modules-definition-flat}
（定義 \ref{modules-definition-flat-at-point}）を参照せよ。
アフィンスキーム上の準連接層について、これは代数の章で定義された概念と一致する。

\begin{lemma}
\label{properties-lemma-flat-module}
\begin{slogan}
加群と層では平坦性は同じである。
\end{slogan}
アフィンスキーム $X = \Spec(R)$ を取る。
$\mathcal{F} = \widetilde{M}$ とする。$R$-加群 $M$ を取る。
準連接層 $\mathcal{F}$ が平坦な
$\mathcal{O}_X$-加群であることと、$M$ が平坦な $R$-加群であることは同値である。
\end{lemma}

\begin{proof}
$\mathcal{F}$ の平坦性は茎上で調べられる。加群の章の補題
\ref{modules-lemma-flat-stalks-flat} を参照せよ。
環上の加群の場合も同様であり、代数の章の補題
\ref{algebra-lemma-flat-localization} による。
そして、$\mathcal{F}_x = M_{\mathfrak p}$ であり、$x$ が
$\mathfrak p$ に対応するので、補題が従う。
\end{proof}

\section{局所自由加群}
\label{properties-section-finite-locally-free}

\noindent
任意の環付き空間について、$\mathcal{O}_X$-加群が（有限）局所自由であることの意味は
既知である。アフィンスキーム上では、これは代数の章で定義された概念と一致する。

\begin{lemma}
\label{properties-lemma-locally-free-module}
アフィンスキーム $X = \Spec(R)$ を取る。
$\mathcal{F} = \widetilde{M}$ とする。ある $R$-加群 $M$ を選ぶ。
準連接層 $\mathcal{F}$ が（有限）局所自由な
$\mathcal{O}_X$-加群であることと、$M$ が（有限）局所自由な
$R$-加群であることは同値である。
\end{lemma}

\begin{proof}
定義から従う。加群の章の定義
\ref{modules-definition-locally-free} および
代数の章の定義 \ref{algebra-definition-locally-free} を参照せよ。
\end{proof}

\noindent
有限局所自由加群は、さまざまな異なる方法で特徴づけられる。

\begin{lemma}
\label{properties-lemma-finite-locally-free}
$X$ をスキームとする。
$\mathcal{F}$ を準連接 $\mathcal{O}_X$-加群とする。
次は同値である。
\begin{enumerate}
\item $\mathcal{F}$ は有限表示の平坦な $\mathcal{O}_X$-加群である。
\item $\mathcal{F}$ は有限表示の $\mathcal{O}_X$-加群であり、
すべての $x \in X$ に対して茎 $\mathcal{F}_x$ は
$\mathcal{O}_{X, x}$-加群として自由である。
\item $\mathcal{F}$ は有限型の局所自由な $\mathcal{O}_X$-加群である。
\item $\mathcal{F}$ は有限局所自由な $\mathcal{O}_X$-加群であり、かつ
\item $\mathcal{F}$ は有限型の $\mathcal{O}_X$-加群であり、
すべての $x \in X$ に対して茎 $\mathcal{F}_x$ は
$\mathcal{O}_{X, x}$-加群として自由であり、関数
$$
\rho_\mathcal{F} : X \to \mathbf{Z}, \quad
x \longmapsto
\dim_{\kappa(x)} \mathcal{F}_x \otimes_{\mathcal{O}_{X, x}} \kappa(x)
$$
は $X$ 上のザリスキー位相で局所定数である。
\end{enumerate}
\end{lemma}

\begin{proof}
この補題は直ちにアフィンの場合へ帰着する。
この場合、補題は代数の章の補題
\ref{algebra-lemma-finite-projective} の言い換えである。
この翻訳では補題 \ref{properties-lemma-finite-type-module}、
\ref{properties-lemma-finite-presentation-module}、
\ref{properties-lemma-flat-module}、および
\ref{properties-lemma-locally-free-module} を用いる。
\end{proof}

\begin{lemma}
\label{properties-lemma-finite-locally-free-reduced}
$X$ を被約スキームとする。$\mathcal{F}$ を準連接
$\mathcal{O}_X$-加群とする。このとき、補題
\ref{properties-lemma-finite-locally-free} の同値な条件は、次の条件とも同値である。
\begin{enumerate}
\item[(6)] $\mathcal{F}$ は有限型の $\mathcal{O}_X$-加群であり、関数
$$
\rho_\mathcal{F} : X \to \mathbf{Z}, \quad
x \longmapsto
\dim_{\kappa(x)} \mathcal{F}_x \otimes_{\mathcal{O}_{X, x}} \kappa(x)
$$
は $X$ 上のザリスキー位相で局所定数である。
\end{enumerate}
\end{lemma}

\begin{proof}
この補題は直ちにアフィンの場合へ帰着する。
この場合、補題は代数の章の補題
\ref{algebra-lemma-finite-projective-reduced} の言い換えである。
\end{proof}

\section{局所射影加群}
\label{properties-section-locally-projective}

\noindent
代数の章で行った議論の帰結として、局所射影加群を次のように定義できる。

\begin{definition}
\label{properties-definition-locally-projective}
$X$ をスキームとする。$\mathcal{F}$ を準連接
$\mathcal{O}_X$-加群とする。$\mathcal{F}$ を {\it 局所射影} と呼ぶのは、
任意のアフィン開集合 $U \subset X$ に対して $\mathcal{O}_X(U)$-加群
$\mathcal{F}(U)$ が射影である場合である。
\end{definition}

\begin{lemma}
\label{properties-lemma-locally-projective}
$X$ をスキームとする。
$\mathcal{F}$ を準連接 $\mathcal{O}_X$-加群とする。
次は同値である。
\begin{enumerate}
\item $\mathcal{F}$ は局所射影であり、かつ
\item アフィン開被覆 $X = \bigcup U_i$ が存在し、$\mathcal{O}_X(U_i)$-加群
$\mathcal{F}(U_i)$ がすべての $i$ について射影である。
\end{enumerate}
特に、$X = \Spec(A)$ かつ $\mathcal{F} = \widetilde{M}$ のとき、
$\mathcal{F}$ が局所射影であることと $M$ が射影 $A$-加群であることは同値である。
\end{lemma}

\begin{proof}
まず、$M$ が射影 $A$-加群で $A \to B$ が環写像ならば、
$M \otimes_A B$ は射影 $B$-加群であることに注意する。代数の章の補題
\ref{algebra-lemma-ascend-properties-modules} を参照せよ。
従って、アフィン開集合 $U$ で $\mathcal{F}(U)$ が射影
$\mathcal{O}_X(U)$-加群であるものに対して、標準開集合 $D(f)$ は
$\mathcal{F}(D(f))$ が射影 $\mathcal{O}_X(D(f))$-加群となるアフィン開集合であり、
すべての $f \in \mathcal{O}_X(U)$ についてこれが成り立つ。

(2) を仮定する。$U \subset X$ を任意のアフィン開集合とする。
開被覆 $U = \bigcup_{j = 1, \ldots, m} D(f_j)$ を有限個の標準開集合
$D(f_j)$ により取り、各 $j$ について $D(f_j)$ がある $U_i$ の標準開集合となるようにできる。これはスキームの章の補題
\ref{schemes-lemma-standard-open-two-affines} による。
従って、$A = \mathcal{O}_X(U)$ と置き、$M$ を $A$-加群として取り、
$\mathcal{F}|_U$ が $M$ に対応するとする。このとき
$M_{f_j}$ は射影 $A_{f_j}$-加群である。従って
$A \to B = \prod A_{f_j}$ は忠実平坦な環写像であり、
$M \otimes_A B$ は射影 $B$-加群である。
ゆえに、代数の章の定理
\ref{algebra-theorem-ffdescent-projectivity} により $M$ は射影である。
\end{proof}

\begin{lemma}
\label{properties-lemma-locally-projective-pullback}
スキームの射 $f : X \to Y$ を取る。
$\mathcal{G}$ を準連接 $\mathcal{O}_Y$-加群とする。
$\mathcal{G}$ が $Y$ 上局所射影ならば、$f^*\mathcal{G}$ は
$X$ 上局所射影である。
\end{lemma}

\begin{proof}
代数の章の補題
\ref{algebra-lemma-ascend-properties-modules} および補題
\ref{properties-lemma-locally-projective} から従う。
\end{proof}

\section{準連接層の延長}
\label{properties-section-extending-quasi-coherent-sheaves}

\noindent
開部分スキーム上の与えられた準連接層をスキーム全体へ延長できることを示せると便利な場合がある。

\begin{lemma}
\label{properties-lemma-extend-trivial}
スキームの準コンパクトな開埋め込み $j : U \to X$ を取る。
\begin{enumerate}
\item $U$ 上の任意の準連接層は $X$ 上の準連接層へ延長できる。
\item $\mathcal{F}$ を $X$ 上の準連接層とする。
 $\mathcal{G} \subset \mathcal{F}|_U$ を準連接部分層とする。
 $\mathcal{H}$ であって、$\mathcal{F}$ の準連接部分層かつ $\mathcal{H}|_U = \mathcal{G}$ が
$\mathcal{F}|_U$ の部分層として成り立つものが存在する。
\item $\mathcal{F}$ を $X$ 上の準連接層とする。
$\mathcal{G}$ を $U$ 上の準連接層とする。
$\varphi : \mathcal{G} \to \mathcal{F}|_U$ を
$\mathcal{O}_U$-加群の射とする。準連接層 $\mathcal{H}$ の
$\mathcal{O}_X$-加群および射 $\psi : \mathcal{H} \to \mathcal{F}$ で、
$\mathcal{H}|_U = \mathcal{G}$ かつ
$\psi|_U = \varphi$ となるものが存在する。
\end{enumerate}
\end{lemma}

\begin{proof}
埋め込みは分離的である（スキームの章の補題
\ref{schemes-lemma-immersions-monomorphisms} を参照）し、仮定により
$j$ は準コンパクトである。
従って任意の準連接層 $\mathcal{G}$ に対して、それが $U$ 上にあれば層
$j_*\mathcal{G}$ は $X$ への延長である。スキームの章の補題
\ref{schemes-lemma-push-forward-quasi-coherent} および
層の章の節 \ref{sheaves-section-open-immersions} を参照せよ。

\medskip\noindent
(2) の $\mathcal{F}$ と $\mathcal{G}$ を仮定する。
すると $j_*\mathcal{G}$ は $X$ 上の準連接層である（上記参照）。
これは $j_*j^*\mathcal{F}$ の部分層である。従って核
$$
\mathcal{H} =
\Ker(\mathcal{F} \oplus j_* \mathcal{G}
\longrightarrow j_*j^*\mathcal{F})
$$
は同様に準連接である。スキームの章の節
\ref{schemes-section-quasi-coherent} を参照せよ。
$\mathcal{H} \subset \mathcal{F}$ および
$\mathcal{H}|_U = \mathcal{G}$（再び層の章の節
\ref{sheaves-section-open-immersions} の内容を用いる）が形式的に確認できる。

\medskip\noindent
(3) は (2) と同じ方法で証明される。すなわち
$\mathcal{H} = \Ker(\mathcal{F} \oplus j_* \mathcal{G}
\to j_*j^*\mathcal{F})$
と置き、$\mathcal{F}$ と $\mathcal{G}$ の $U$ 上での明らかな射および
明らかな同一視を取ればよい。
\end{proof}

\begin{lemma}
\label{properties-lemma-extend}
$X$ を準コンパクトかつ準分離的なスキームとする。
$U \subset X$ を準コンパクトな開集合とする。
$\mathcal{F}$ を準連接 $\mathcal{O}_X$-加群とする。
$\mathcal{G} \subset \mathcal{F}|_U$ を有限型の準連接
$\mathcal{O}_U$-部分加群とする。このとき
$\mathcal{G}' \subset \mathcal{F}$ なる有限型の準連接部分加群で
$\mathcal{G}'|_U = \mathcal{G}$ を満たすものが存在する。
\end{lemma}

\begin{proof}
最小個数を $n$ とする。これは、アフィン開集合 $U_i \subset X$（$i = 1, \ldots , n$）を選んで
$X = U \cup \bigcup U_i$ と書ける最小個数である。（ここでは $X$ の準コンパクト性を用いた。）
$n = 1$ の場合に補題を証明できると仮定する。このとき $\mathcal{G}$ を
$\mathcal{G}_1$ へ $U \cup U_1$ 上で、次に
$\mathcal{G}_2$ へ $U \cup U_1 \cup U_2$ 上で、
$\mathcal{G}_3$ へ $U \cup U_1 \cup U_2 \cup U_3$ 上で、というように
逐次延長できる。従って $n = 1$ の場合へ帰着する。

\medskip\noindent
従って $X = U \cup V$ で $V$ がアフィンであると仮定してよい。
$X$ は準分離的で、$U$ と $V$ は準コンパクト開なので、
$U \cap V$ は準コンパクト開である。系
$(V, U \cap V, \mathcal{F}|_V, \mathcal{G}|_{U \cap V})$
に対して補題を証明すれば十分である。得られた層
$\mathcal{G}'$ を $V$ 上で、与えられた層 $\mathcal{G}$ と $U$ 上で貼り合わせ、
$U \cap V$ 上の共通の値に沿って貼り合わせられるからである。
従って $X$ がアフィンの場合へ帰着する。

\medskip\noindent
$X = \Spec(R)$ と仮定する。$\mathcal{F} = \widetilde M$ と書く。
ここで $R$-加群 $M$ を取る。上の補題
\ref{properties-lemma-extend-trivial} により、$\mathcal{H} \subset \mathcal{F}$ なる準連接部分層で、
$\mathcal{G}$ に $U$ 上で制限されるものを取れる。
$\mathcal{H} = \widetilde N$ と書く。ここで $R$-加群 $N$ を取る。
任意の $u \in U$ に対して、$f \in R$ で $u \in D(f) \subset U$ かつ
$N_f$ が有限生成となるものが存在する。補題
\ref{properties-lemma-finite-type-module} を参照せよ。
$U$ は準コンパクトなので、有限個の $D(f_i)$ で覆え、各 $N_{f_i}$ は有限個の元、例えば
$x_{i, 1}/f_i^N, \ldots, x_{i, r_i}/f_i^N$ により生成される。
$N' \subset N$ を元 $x_{i, j}$ が生成する部分加群とする。このとき
部分層
$\mathcal{G}' = \widetilde{N'} \subset \mathcal{H} \subset \mathcal{F}$
が求めるものである。
\end{proof}

\begin{lemma}
\label{properties-lemma-quasi-coherent-colimit-finite-type}
$X$ を準コンパクトかつ準分離的なスキームとする。
$\mathcal{O}_X$-加群の任意の準連接層は、その有限型な
$\mathcal{O}_X$-準連接部分加群の有向余極限である。
\end{lemma}

\begin{proof}
余極限が有向であるのは、$\mathcal{G}_1$ と $\mathcal{G}_2$ が
有限型の準連接部分層ならば、
$\mathcal{G}_1 \oplus \mathcal{G}_2 \to \mathcal{F}$ の像が
有限型の準連接部分加群となるからである。
$U \subset X$ を任意のアフィン開集合とし、
$s \in \Gamma(U, \mathcal{F})$ を任意の切断とする。
$\mathcal{G} \subset \mathcal{F}|_U$ を $s$ が生成する部分層とする。
すると明らかに $\mathcal{G}$ は準連接であり、
$\mathcal{O}_U$-加群として有限型である。
補題 \ref{properties-lemma-extend} により、$\mathcal{G}$ は有限型の準連接部分層
$\mathcal{G}' \subset \mathcal{F}$ の制限である。
$X$ の位相がアフィン開集合からなる基底を持つことから、
$\mathcal{F}$ のすべての局所切断は、有限型の準連接部分加群に
局所的に含まれる。これで証明が終わる。
\end{proof}

\begin{lemma}
\label{properties-lemma-extend-finite-presentation}
$X$ を準コンパクトかつ準分離的なスキームとする。
$\mathcal{F}$ を準連接 $\mathcal{O}_X$-加群とする。
$U \subset X$ を準コンパクトな開集合とする。
$\mathcal{G}$ を有限表示の $\mathcal{O}_U$-加群とする。
$\varphi : \mathcal{G} \to \mathcal{F}|_U$ を
$\mathcal{O}_U$-加群の射とする。
このとき、有限表示の $\mathcal{O}_X$-加群
$\mathcal{G}'$ と、$\mathcal{O}_X$-加群の射
$\varphi' : \mathcal{G}' \to \mathcal{F}$ で
$\mathcal{G}'|_U = \mathcal{G}$ かつ
$\varphi'|_U = \varphi$ を満たすものが存在する。
\end{lemma}

\begin{proof}
証明の冒頭は補題 \ref{properties-lemma-extend} の冒頭の繰り返しである。
ここでも注意深く書き下す。

\medskip\noindent
最小個数を $n$ とする。これは、アフィン開集合 $U_i \subset X$（$i = 1, \ldots , n$）を選んで
$X = U \cup \bigcup U_i$ と書ける最小個数である。
（ここでは $X$ の準コンパクト性を用いた。）
$n = 1$ の場合に補題を証明できると仮定する。このとき対
$(\mathcal{G}, \varphi)$ を $(\mathcal{G}_1, \varphi_1)$ へ $U \cup U_1$ 上で、次に
$(\mathcal{G}_2, \varphi_2)$ へ $U \cup U_1 \cup U_2$ 上で、
$(\mathcal{G}_3, \varphi_3)$ へ $U \cup U_1 \cup U_2 \cup U_3$ 上で逐次延長できる。
従って $n = 1$ の場合へ帰着する。

\medskip\noindent
従って $X = U \cup V$ で $V$ がアフィンであると仮定してよい。
$X$ は準分離的で $U$ は準コンパクトなので、
$U \cap V \subset V$ は準コンパクトである。
系
$(V, U \cap V, \mathcal{F}|_V, \mathcal{G}|_{U \cap V}, \varphi|_{U \cap V})$
に対して補題を証明し、$(\mathcal{G}', \varphi')$ を $V$ 上で得たとする。
すると $\mathcal{G}'$ を $V$ 上で与えられた層 $\mathcal{G}$ と $U$ 上で貼り合わせ、
$U \cap V$ 上の共通の値に沿って貼り合わせられ、同様に
$\varphi'$ を $\varphi$ と $U \cap V$ 上の共通の値に沿って貼り合わせられる。
従って $X$ がアフィンの場合へ帰着する。

\medskip\noindent
$X = \Spec(R)$ と仮定する。上の補題
\ref{properties-lemma-extend-trivial} により、$\mathcal{H}$ を取り、
$\psi : \mathcal{H} \to \mathcal{F}$ を持つ $X$ 上の準連接層で、
$\mathcal{G}$ と $\varphi$ に $U$ 上で制限されるものとする。
補題 \ref{properties-lemma-extend} により、有限型の準連接
$\mathcal{O}_X$-部分加群
$\mathcal{H}' \subset \mathcal{H}$ で
$\mathcal{H}'|_U = \mathcal{G}$ を満たすものを取れる。
従って $\mathcal{H}$ を $\mathcal{H}'$ に置き換え、
$\psi$ を $\psi$ の $\mathcal{H}'$ への制限に置き換えることで、
$\mathcal{H}$ は有限型であると仮定してよい。
補題 \ref{properties-lemma-finite-presentation-module} により、
$\mathcal{H} = \widetilde{N}$ で $N$ が有限生成 $R$-加群であると分かる。
従って、準連接 $\mathcal{O}_X$-加群の次の短完全列にあるような全射が存在する。
$$
0 \to \mathcal{K} \to \mathcal{O}_X^{\oplus n} \to \mathcal{H} \to 0
$$
ここで $\mathcal{K}$ は核として定義される。
$\mathcal{G}$ は有限表示で、補題
\ref{modules-lemma-kernel-surjection-finite-free-onto-finite-presentation} により
$\mathcal{H}|_U = \mathcal{G}$ から
$\mathcal{K}|_U$ は有限型の $\mathcal{O}_U$-加群である。
従って再び補題 \ref{properties-lemma-extend} により、
有限型の準連接 $\mathcal{O}_X$-部分加群
$\mathcal{K}' \subset \mathcal{K}$ で
$\mathcal{K}'|_U = \mathcal{K}|_U$ を満たすものが存在する。
補題で問われた問題の解は
$$
\mathcal{G}' = \mathcal{O}_X^{\oplus n}/\mathcal{K}'
$$
と置くことで得られる。これは明らかに有限表示であり、
$\mathcal{G}$ に $U$ 上で制限され、$\varphi'$ は合成
$$
\mathcal{G}' = \mathcal{O}_X^{\oplus n}/\mathcal{K}'
\to \mathcal{O}_X^{\oplus n}/\mathcal{K} = \mathcal{H} \xrightarrow{\psi}
\mathcal{F}.
$$
に等しい。
これで補題の証明が完了する。
\end{proof}

\begin{lemma}
\label{properties-lemma-lift-finite-presentation}
$X$ を準コンパクトかつ準分離的なスキームとする。$U \subset X$ を
準コンパクトな開集合とする。$\mathcal{G}$ を $\mathcal{O}_U$-加群とする。
\begin{enumerate}
\item $\mathcal{G}$ が有限型の準連接加群ならば、有限型の準連接
$\mathcal{O}_X$-加群 $\mathcal{G}'$ で
$\mathcal{G}'|_U = \mathcal{G}$ を満たすものが存在する。
\item $\mathcal{G}$ が有限表示ならば、有限表示の
$\mathcal{O}_X$-加群 $\mathcal{G}'$ で
$\mathcal{G}'|_U = \mathcal{G}$ を満たすものが存在する。
\end{enumerate}
\end{lemma}

\begin{proof}
(2) は補題 \ref{properties-lemma-extend-finite-presentation} において
$\mathcal{F} = 0$ とした特別な場合である。(1) については、まず
$\mathcal{G} = \mathcal{F}|_U$ と、補題 \ref{properties-lemma-extend-trivial} による
ある準連接 $\mathcal{O}_X$-加群について書き、次に
$\mathcal{G} = \mathcal{F}|_U$ として補題 \ref{properties-lemma-extend} を適用する。
\end{proof}

\noindent
次の補題は、準コンパクトかつ準分離的なスキーム上の任意の準連接層が、
有限表示の $\mathcal{O}$-加群のフィルター余極限であることを述べる。
実際、次の補題ではこれを有向集合上の有向余極限という
（おそらくより馴染み深い）言葉で言い換える。

\begin{lemma}
\label{properties-lemma-directed-colimit-diagram-finite-presentation}
\begin{slogan}
準コンパクトかつ準分離的なスキーム上の準連接加群は、
有限表示加群のフィルター余極限である。
\end{slogan}
$X$ をスキームとし、$X$ は準コンパクトかつ準分離的であると仮定する。
$\mathcal{F}$ を準連接 $\mathcal{O}_X$-加群とする。次が存在する。
\begin{enumerate}
\item フィルター添字圏 $\mathcal{I}$（圏の章の定義
\ref{categories-definition-directed} を参照）、
\item 図式 $\mathcal{I} \to \textit{Mod}(\mathcal{O}_X)$（圏の章の節
\ref{categories-section-limits} を参照）、
$i \mapsto \mathcal{F}_i$、
\item $\mathcal{O}_X$-加群の射
$\varphi_i : \mathcal{F}_i \to \mathcal{F}$。
\end{enumerate}
ここで各 $\mathcal{F}_i$ は有限表示であり、射 $\varphi_i$ は同型
$$
\colim_i \mathcal{F}_i
=
\mathcal{F}.
$$
を誘導する。
\end{lemma}

\begin{proof}
集合 $I$ を選び、各 $i \in I$ に対して、有限表示の
$\mathcal{O}_X$-加群と $\mathcal{O}_X$-加群の準同型
$\varphi_i : \mathcal{F}_i \to \mathcal{F}$ を次の性質を満たすように選ぶ。
任意の $\psi : \mathcal{G} \to \mathcal{F}$ で $\mathcal{G}$ が有限表示のものに対し、
$i \in I$ が存在して、同型 $\alpha : \mathcal{F}_i \to \mathcal{G}$ が存在し、
$\varphi_i = \psi \circ \alpha$ を満たす。
そのような集合の存在は、加群の章の補題
\ref{modules-lemma-set-isomorphism-classes-finite-type-modules}
（およびその証明）から明らかである。
$\mathcal{I}$ を $\Ob(\mathcal{I}) = I$ を持つ圏と記し、
$i, i' \in I$ が与えられたとき、
$$
\Mor_\mathcal{I}(i, i') =
\{\alpha : \mathcal{F}_i \to \mathcal{F}_{i'} \mid
\alpha \circ \varphi_{i'} = \varphi_i
\}.
$$
と置く。$\mathcal{I}$ はフィルター圏であり、
$\mathcal{F} = \colim_i \mathcal{F}_i$ であると主張する。

\medskip\noindent
$i, i' \in I$ とする。このとき射
$$
\mathcal{F}_i \oplus \mathcal{F}_{i'} \longrightarrow \mathcal{F}
$$
を考えられる。これは $\varphi_i$ と $\varphi_{i'}$ の直和である。
有限表示 $\mathcal{O}_X$-加群の直和は有限表示なので、ある $i'' \in I$ が存在し、
$\varphi_{i''} : \mathcal{F}_{i''} \to \mathcal{F}$ は上に表示した
$\mathcal{F}$ への射と同型になる。可換図式
$$
\xymatrix{
\mathcal{F}_i \ar[r] \ar[d] & \mathcal{F} \ar@{=}[d] \\
\mathcal{F}_i \oplus \mathcal{F}_{i'} \ar[r] & \mathcal{F}
}
\quad
\text{and}
\quad
\xymatrix{
\mathcal{F}_{i'} \ar[r] \ar[d] & \mathcal{F} \ar@{=}[d] \\
\mathcal{F}_i \oplus \mathcal{F}_{i'} \ar[r] & \mathcal{F}
}
$$
があるので、射 $i \to i''$ および $i' \to i''$ が $\mathcal{I}$ に存在する。
次に、$i, i' \in I$ と射 $\alpha, \beta : i \to i'$
（$\mathcal{O}_X$-加群写像
$\alpha, \beta : \mathcal{F}_i \to \mathcal{F}_{i'}$ に対応する）があるとする。
この場合、余等化子
$$
\mathcal{G} =
\Coker(
\mathcal{F}_i \xrightarrow{\alpha - \beta} \mathcal{F}_{i'}
)
$$
を考える。$\mathcal{G}$ は有限表示の $\mathcal{O}_X$-加群であることに注意する。
圏 $\mathcal{I}$ の射の定義により
$\varphi_{i'} \circ \alpha = \varphi_{i'} \circ \beta$ なので、誘導射
$\psi : \mathcal{G} \to \mathcal{F}$ を得る。従って再び、対
$(\mathcal{G}, \psi)$ は対 $(\mathcal{F}_{i''}, \varphi_{i''})$ と、ある $i''$ について同型である。
従って射 $i' \to i''$ が $\mathcal{I}$ に存在し、$\alpha$ と $\beta$ を等化する。
以上により圏 $\mathcal{I}$ がフィルター圏であることが示された。

\medskip\noindent
図式の余極限が $\mathcal{F}$ であることをなお示さなければならない。
余極限の定義と圏 $\mathcal{I}$ の定義により、標準射
$$
\varphi :
\colim_i \mathcal{F}_i
\longrightarrow
\mathcal{F}.
$$
がある。$x \in X$ を取る。$\varphi_x$ が同型であることを示そう。
$$
(\colim_i \mathcal{F}_i)_x
=
\colim_i \mathcal{F}_{i, x},
$$
であることを思い出す。層の章の節
\ref{sheaves-section-limits-sheaves} を参照せよ。
まず写像 $\varphi_x$ が単射であることを示す。
$s \in \mathcal{F}_{i, x}$ を、$s$ が $\mathcal{F}_x$ で零へ写る元とする。
このとき、準コンパクトな開集合 $U$ が存在し、$s$ は
$s \in \mathcal{F}_i(U)$ から来て、$\varphi_i(s) = 0$ が
$\mathcal{F}(U)$ で成り立つ。補題 \ref{properties-lemma-extend} により、有限型の準連接部分層
$\mathcal{K} \subset \Ker(\varphi_i)$ で、準連接 $\mathcal{O}_U$-部分加群として
$\mathcal{F}_i$ の中で $s$ が生成するものへ制限されるものを取れる：
$\mathcal{K}|_U = \mathcal{O}_U\cdot s \subset \mathcal{F}_i|_U$。
明らかに $\mathcal{F}_i/\mathcal{K}$ は有限表示であり、写像 $\varphi_i$ は商写像
$\mathcal{F}_i \to \mathcal{F}_i/\mathcal{K}$ を経由する。従って、
$i' \in I$ と射 $\alpha : \mathcal{F}_i \to \mathcal{F}_{i'}$ で $\mathcal{I}$ に属し、
商写像 $\mathcal{F}_i \to \mathcal{F}_i/\mathcal{K}$ と同一視できるものが存在する。
すると切断 $s$ は $\mathcal{F}_{i'}(U)$ で、従って特に
$(\colim_i \mathcal{F}_i)_x =
\colim_i \mathcal{F}_{i, x}$ で零へ写る。これで単射性が従う。
最後に写像 $\varphi_x$ が全射であることを示す。
$s \in \mathcal{F}_x$ を取る。準コンパクトな開近傍
$U \subset X$ を $x$ に対して、$s$ が切断 $s \in \mathcal{F}(U)$ に対応するように選ぶ。
写像 $s : \mathcal{O}_U \to \mathcal{F}$（$s$ による乗法）を考える。
補題 \ref{properties-lemma-extend-finite-presentation} により、有限表示の
$\mathcal{O}_X$-加群 $\mathcal{G}$ と $\mathcal{O}_X$-加群写像
$\mathcal{G} \to \mathcal{F}$ が存在し、$\mathcal{G}|_U \to \mathcal{F}|_U$ は
$s : \mathcal{O}_U \to \mathcal{F}$ と同一視される。
再び $\mathcal{I}$ の定義により $i \in I$ が存在し、
$\mathcal{G} \to \mathcal{F}$ は $\varphi_i : \mathcal{F}_i \to \mathcal{F}$ と同型である。
明らかに切断 $s' \in \mathcal{F}_i(U)$ が存在し、$s \in \mathcal{F}(U)$ へ写る。
これで全射性が示され、補題の証明が完了する。
\end{proof}

\begin{lemma}
\label{properties-lemma-directed-colimit-finite-presentation}
$X$ をスキームとし、$X$ は準コンパクトかつ準分離的であると仮定する。
$\mathcal{F}$ を準連接 $\mathcal{O}_X$-加群とする。次が存在する。
\begin{enumerate}
\item 有向集合 $I$（圏の章の定義
\ref{categories-definition-directed-set} を参照）、
\item 系 $(\mathcal{F}_i, \varphi_{ii'})$ で、$I$ 上かつ
$\textit{Mod}(\mathcal{O}_X)$ におけるもの（圏の章の定義
\ref{categories-definition-system-over-poset} を参照）、
\item $\mathcal{O}_X$-加群の射
$\varphi_i : \mathcal{F}_i \to \mathcal{F}$。
\end{enumerate}
ここで各 $\mathcal{F}_i$ は有限表示であり、射 $\varphi_i$ は同型
$$
\colim_i \mathcal{F}_i
=
\mathcal{F}.
$$
を誘導する。
\end{lemma}

\begin{proof}
これは補題 \ref{properties-lemma-directed-colimit-diagram-finite-presentation} と
圏の章の補題 \ref{categories-lemma-directed-category-system} の直接の帰結である
（$\mathcal{O}_X$-加群の層の圏で余極限が存在するという事実も併用する。
層の章の節 \ref{sheaves-section-limits-sheaves} を参照）。
\end{proof}

\begin{lemma}
\label{properties-lemma-finite-directed-colimit-surjective-maps}
$X$ をスキームとし、$X$ は準コンパクトかつ準分離的であると仮定する。
$\mathcal{F}$ を有限型の準連接 $\mathcal{O}_X$-加群とする。
このとき $\mathcal{F} = \colim \mathcal{F}_i$ と書け、$\mathcal{F}_i$ は有限表示で、
すべての遷移写像 $\mathcal{F}_i \to \mathcal{F}_{i'}$ は全射である。
\end{lemma}

\begin{proof}
$\mathcal{F} = \colim \mathcal{G}_i$ を、有限表示 $\mathcal{O}_X$-加群の
フィルター余極限として書く（補題
\ref{properties-lemma-directed-colimit-finite-presentation}）。
ある $\mathcal{G}_i \to \mathcal{F}$ が、ある $i$ について全射であると主張する。
実際、有限アフィン開被覆 $X = U_1 \cup \ldots \cup U_m$ を選ぶ。
切断 $s_{jl} \in \mathcal{F}(U_j)$ を、$\mathcal{F}|_{U_j}$ を生成するように選ぶ。
補題 \ref{properties-lemma-finite-type-module} を参照せよ。
層の章の補題 \ref{sheaves-lemma-directed-colimits-sections} により、
$s_{jl}$ は $\mathcal{G}_i \to \mathcal{F}$ の像に十分大きい $i$ について入る。
従って $\mathcal{G}_i \to \mathcal{F}$ は、十分大きい $i$ に対して全射である。
そのような $i$ を選び、$\mathcal{K} \subset \mathcal{G}_i$ を写像
$\mathcal{G}_i \to \mathcal{F}$ の核とする。
$\mathcal{K} = \colim \mathcal{K}_a$ を、その有限型準連接部分加群の
フィルター余極限として書く（補題
\ref{properties-lemma-quasi-coherent-colimit-finite-type}）。すると
$\mathcal{F} = \colim \mathcal{G}_i/\mathcal{K}_a$ が補題で問われた問題の解である。
\end{proof}

\begin{lemma}
\label{properties-lemma-application-directed-colimit}
$X$ を準コンパクトかつ準分離的なスキームとする。
$\mathcal{F}$ を有限型の準連接 $\mathcal{O}_X$-加群とする。
$U \subset X$ を準コンパクトな開集合とし、$\mathcal{F}|_U$ は有限表示であるとする。
このとき $\mathcal{O}_X$-加群の写像 $\varphi : \mathcal{G} \to \mathcal{F}$ で、
(a) $\mathcal{G}$ は有限表示、
(b) $\varphi$ は全射、かつ
(c) $\varphi|_U$ は同型
となるものが存在する。
\end{lemma}

\begin{proof}
$\mathcal{F} = \colim \mathcal{F}_i$ を、各 $\mathcal{F}_i$ が有限表示である
有向余極限として書く。補題
\ref{properties-lemma-directed-colimit-finite-presentation} を参照せよ。
有限アフィン開被覆 $X = \bigcup V_j$ を選び、有限個の切断
$s_{jl} \in \mathcal{F}(V_j)$ を、$\mathcal{F}|_{V_j}$ を生成するように選ぶ。補題
\ref{properties-lemma-finite-type-module} を参照せよ。
層の章の補題 \ref{sheaves-lemma-directed-colimits-sections} により、
$s_{jl}$ は $\mathcal{F}_i \to \mathcal{F}$ の像に十分大きい $i$ について入る。
従って $\mathcal{F}_i \to \mathcal{F}$ は、十分大きい $i$ に対して全射である。
そのような $i$ を選び、$\mathcal{K} \subset \mathcal{F}_i$ を写像
$\mathcal{F}_i \to \mathcal{F}$ の核とする。$\mathcal{F}_U$ は有限表示なので、
$\mathcal{K}|_U$ は有限型である。加群の章の補題
\ref{modules-lemma-kernel-surjection-finite-free-onto-finite-presentation} を参照せよ。
従って有限型の準連接部分加群 $\mathcal{K}' \subset \mathcal{K}$ で
$\mathcal{K}'|_U = \mathcal{K}|_U$ を満たすものを取れる。補題
\ref{properties-lemma-extend} を参照せよ。すると
$\mathcal{G} = \mathcal{F}_i/\mathcal{K}'$
と与えられた写像 $\mathcal{G} \to \mathcal{F}$ が解である。
\end{proof}

\noindent
$X$ をスキームとする。次の補題では、{\it 有限表示の準連接
$\mathcal{O}_X$-代数 $\mathcal{A}$} という概念を用いる。これは、任意のアフィン開集合
$\Spec(R) \subset X$ に対して $\mathcal{A} = \widetilde{A}$ となり、
$A$ が（可換）$R$-代数で、$R$-代数として有限表示であることを意味する。

\begin{lemma}
\label{properties-lemma-algebra-directed-colimit-finite-presentation}
$X$ をスキームとし、$X$ は準コンパクトかつ準分離的であると仮定する。
$\mathcal{A}$ を準連接 $\mathcal{O}_X$-代数とする。次が存在する。
\begin{enumerate}
\item 有向集合 $I$（圏の章の定義
\ref{categories-definition-directed-set} を参照）、
\item 系 $(\mathcal{A}_i, \varphi_{ii'})$ で、$I$ 上かつ
$\mathcal{O}_X$-代数の圏におけるもの、
\item $\mathcal{O}_X$-代数の射
$\varphi_i : \mathcal{A}_i \to \mathcal{A}$。
\end{enumerate}
ここで各 $\mathcal{A}_i$ は有限表示の準連接 $\mathcal{O}_X$-代数であり、
射 $\varphi_i$ は同型
$$
\colim_i \mathcal{A}_i
=
\mathcal{A}.
$$
を誘導する。
\end{lemma}

\begin{proof}
まず補題 \ref{properties-lemma-directed-colimit-finite-presentation} のように、
$\mathcal{A} = \colim_i \mathcal{F}_i$ を有限表示準連接層の有向余極限として書く。
各 $i$ に対して $\mathcal{B}_i = \text{Sym}(\mathcal{F}_i)$ を、
$\mathcal{F}_i$ の $\mathcal{O}_X$ 上の対称代数とする。
$\mathcal{I}_i = \Ker(\mathcal{B}_i \to \mathcal{A})$ と書く。
$\mathcal{I}_i = \colim_j \mathcal{F}_{i, j}$ と書く。ここで
$\mathcal{F}_{i, j}$ は $\mathcal{I}_i$ の有限型準連接部分加群である。補題
\ref{properties-lemma-quasi-coherent-colimit-finite-type} を参照せよ。
$\mathcal{I}_{i, j} \subset \mathcal{I}_i$ を、$\mathcal{B}_i$-イデアルで
$\mathcal{F}_{i, j}$ が生成するものに等しいと置く。
$\mathcal{A}_{i, j} = \mathcal{B}_i/\mathcal{I}_{i, j}$ と置く。
すると $\mathcal{A}_{i, j}$ は有限表示の準連接 $\mathcal{O}_X$-代数である。
$(i, j) \leq (i', j')$ を、$i \leq i'$ であり、かつ写像 $\mathcal{B}_i \to \mathcal{B}_{i'}$ がイデアル
$\mathcal{I}_{i, j}$ をイデアル $\mathcal{I}_{i', j'}$ に写すとき、
成り立つと定義する。このとき明らかに
$\mathcal{A} = \colim_{i, j} \mathcal{A}_{i, j}$ である。
\end{proof}

\noindent
$X$ をスキームとする。次の補題では、{\it 有限型の準連接
$\mathcal{O}_X$-代数 $\mathcal{A}$} という概念を用いる。これは、任意のアフィン開集合
$\Spec(R) \subset X$ に対して $\mathcal{A} = \widetilde{A}$ となり、
$A$ が（可換）$R$-代数で、$R$-代数として有限型であることを意味する。

\begin{lemma}
\label{properties-lemma-algebra-directed-colimit-finite-type}
$X$ をスキームとし、$X$ は準コンパクトかつ準分離的であると仮定する。
$\mathcal{A}$ を準連接 $\mathcal{O}_X$-代数とする。
すると $\mathcal{A}$ は、その有限型準連接
$\mathcal{O}_X$-部分代数の有向余極限である。
\end{lemma}

\begin{proof}
$\mathcal{A}_1, \mathcal{A}_2 \subset \mathcal{A}$ が有限型の準連接
$\mathcal{O}_X$-部分代数ならば、
$\mathcal{A}_1 \otimes_{\mathcal{O}_X} \mathcal{A}_2 \to \mathcal{A}$ の像も
有限型の準連接 $\mathcal{O}_X$-部分代数であり（細部の一部は省略する）、
$\mathcal{A}_1$ と $\mathcal{A}_2$ の両方を含む。
このようにして系が有向であることが分かる。
$\mathcal{A}$ がこの系の余極限であることを示すため、補題
\ref{properties-lemma-algebra-directed-colimit-finite-presentation} のように
$\mathcal{A} = \colim_i \mathcal{A}_i$ を有限表示の準連接
$\mathcal{O}_X$-代数の有向余極限として書く。
すると像 $\mathcal{A}'_i = \Im(\mathcal{A}_i \to \mathcal{A})$ は
$\mathcal{A}$ の有限型準連接部分代数である。
$\mathcal{A}$ はこれらの余極限なので、結果が従う。
\end{proof}

\noindent
$X$ をスキームとする。次の補題では、{\it 有限な（それぞれ整な）準連接
$\mathcal{O}_X$-代数 $\mathcal{A}$} という概念を用いる。これは任意のアフィン開集合
$\Spec(R) \subset X$ に対し $\mathcal{A} = \widetilde{A}$ となり、
$A$ が（可換）$R$-代数で、$R$-代数として有限（それぞれ整）であることを意味する。

\begin{lemma}
\label{properties-lemma-finite-algebra-directed-colimit-finite-finitely-presented}
$X$ をスキームとし、$X$ は準コンパクトかつ準分離的であると仮定する。
$\mathcal{A}$ を有限な準連接 $\mathcal{O}_X$-代数とする。
すると $\mathcal{A} = \colim \mathcal{A}_i$ は、有限かつ有限表示の準連接
$\mathcal{O}_X$-代数の有向余極限であり、すべての遷移写像
$\mathcal{A}_{i'} \to \mathcal{A}_i$ は全射である。
\end{lemma}

\begin{proof}
補題 \ref{properties-lemma-finite-directed-colimit-surjective-maps} により、有限表示の
$\mathcal{O}_X$-加群 $\mathcal{F}$ と全射 $\mathcal{F} \to \mathcal{A}$ が存在する。
代数構造を用いて全射
$$
\text{Sym}^*_{\mathcal{O}_X}(\mathcal{F}) \longrightarrow \mathcal{A}
$$
を得る。$\mathcal{J}$ を核とする。$\mathcal{J} = \colim \mathcal{E}_i$ を、有限型
$\mathcal{O}_X$-部分加群 $\mathcal{E}_i$ のフィルター余極限として書く
（補題 \ref{properties-lemma-quasi-coherent-colimit-finite-type}）。次のように置く。
$$
\mathcal{A}_i = \text{Sym}^*_{\mathcal{O}_X}(\mathcal{F})/(\mathcal{E}_i)
$$
ここで $(\mathcal{E}_i)$ は $\mathcal{E}_i \to \text{Sym}^*_{\mathcal{O}_X}(\mathcal{F})$
の像が生成するイデアル層を表す。すると各 $\mathcal{A}_i$ は有限表示の
$\mathcal{O}_X$-代数であり、遷移写像は全射で、
$\mathcal{A} = \colim \mathcal{A}_i$ である。証明を終えるには、$\mathcal{A}_i$ が有限な
$\mathcal{O}_X$-代数であることを、十分大きい $i$ に対しなお示さなければならない。
そのためにアフィン開被覆 $X = U_1 \cup \ldots \cup U_m$ を選ぶ。
生成元 $f_{j, 1}, \ldots, f_{j, N_j} \in \Gamma(U_i, \mathcal{F})$ を取る。
$\mathcal{A}(U_j)$ は有限な $\mathcal{O}_X(U_j)$-代数なので、各 $k$ に対して
モニック多項式 $P_{j, k} \in \mathcal{O}(U_j)[T]$ が存在し、
$P_{j, k}(f_{j, k})$ は $\mathcal{A}(U_j)$ で零となる。
構成により $\mathcal{A} = \colim \mathcal{A}_i$ なので、
$P_{j, k}(f_{j, k}) = 0$ が $\mathcal{A}_i(U_j)$ で、十分大きいすべての $i$ に対して成り立つ。
そのような $i$ に対して代数 $\mathcal{A}_i$ は有限である。
\end{proof}

\begin{lemma}
\label{properties-lemma-integral-algebra-directed-colimit-finite}
$X$ をスキームとし、$X$ は準コンパクトかつ準分離的であると仮定する。
$\mathcal{A}$ を整な準連接 $\mathcal{O}_X$-代数とする。このとき
\begin{enumerate}
\item $\mathcal{A}$ は、その有限な準連接
$\mathcal{O}_X$-部分代数の有向余極限であり、かつ
\item $\mathcal{A}$ は、有限かつ有限表示の準連接
$\mathcal{O}_X$-代数の直接余極限である。
\end{enumerate}
\end{lemma}

\begin{proof}
補題 \ref{properties-lemma-algebra-directed-colimit-finite-type} により
$\mathcal{A} = \colim \mathcal{A}_i$ であり、ここで
$\mathcal{A}_i \subset \mathcal{A}$ は有限型の準連接
$\mathcal{O}_X$-代数すべてを走る。
$\mathcal{O}_X$ 上の任意の有限型準連接部分代数で、$\mathcal{A}$ の部分代数であるものは有限である
（$\mathcal{A}_i(U) \subset \mathcal{A}(U)$ へ、アフィン開集合 $U$ を $X$ の中で取り、代数の章の補題
\ref{algebra-lemma-characterize-finite-in-terms-of-integral} を適用せよ）。
これで (1) が示された。

\medskip\noindent
(2) を示すため、補題 \ref{properties-lemma-directed-colimit-finite-presentation} を用いて
$\mathcal{A} = \colim \mathcal{F}_i$ を有限表示 $\mathcal{O}_X$-加群の余極限として書く。
各 $i$ に対し、写像
$$
\text{Sym}^*_{\mathcal{O}_X}(\mathcal{F}_i) \longrightarrow \mathcal{A}
$$
の核を $\mathcal{J}_i$ とする。$i' \geq i$ に対して誘導写像
$\mathcal{J}_i \to \mathcal{J}_{i'}$ があり、
$\mathcal{A} =
\colim \text{Sym}^*_{\mathcal{O}_X}(\mathcal{F}_i)/\mathcal{J}_i$ である。
さらに、準連接 $\mathcal{O}_X$-代数
$\text{Sym}^*_{\mathcal{O}_X}(\mathcal{F}_i)/\mathcal{J}_i$ は有限である（上記参照）。
$\mathcal{J}_i = \colim \mathcal{E}_{ik}$ を有限表示 $\mathcal{O}_X$-加群の余極限として書く。
$i' \geq i$ と $k$ が与えられれば、ある $k'$ が存在し、写像
$\mathcal{E}_{ik} \to \mathcal{E}_{i'k'}$ を取って図式
$$
\xymatrix{
\mathcal{J}_i \ar[r] & \mathcal{J}_{i'} \\
\mathcal{E}_{ik} \ar[u] \ar[r] & \mathcal{E}_{i'k'} \ar[u]
}
$$
を可換にできる。これは加群の章の補題
\ref{modules-lemma-finite-presentation-quasi-compact-colimit} から従う。
これにより写像
$$
\mathcal{A}_{ik} =
\text{Sym}^*_{\mathcal{O}_X}(\mathcal{F}_i)/(\mathcal{E}_{ik})
\longrightarrow
\text{Sym}^*_{\mathcal{O}_X}(\mathcal{F}_{i'})/(\mathcal{E}_{i'k'}) =
\mathcal{A}_{i'k'}
$$
が誘導される。ここで $(\mathcal{E}_{ik})$ は $\mathcal{E}_{ik}$ が生成するイデアルを表す。
準連接 $\mathcal{O}_X$-代数 $\mathcal{A}_{ki}$ は有限表示であり、十分大きい $k$ に対して有限である
（補題 \ref{properties-lemma-finite-algebra-directed-colimit-finite-finitely-presented} の証明を参照）。
最後に
$$
\colim \mathcal{A}_{ik} = \colim \mathcal{A}_i = \mathcal{A}
$$
である。実際、第一の等式は補題
\ref{properties-lemma-finite-algebra-directed-colimit-finite-finitely-presented} の証明で示され、
第二の等式は $\mathcal{A}$ が加群 $\mathcal{F}_i$ の余極限だからである。
\end{proof}

\section{Gabber の結果}
\label{properties-section-gabber}

\noindent
この節では Gabber の結果を証明する。それは、任意のスキーム上にある基数
$\kappa$ が存在して、任意の準連接加群 $\mathcal{F}$ がその準連接な
$\kappa$-生成部分層の合併となることを保証する。従って、スキーム上の準連接層の圏は、
極限と十分な入射対象を持つ Grothendieck アーベル圏である\footnote{Akhil Mathew による
\href{https://amathew.wordpress.com/2011/07/30/quasi-coherent-sheaves-presentable-categories-and-a-result-of-gabber/}{ブログ記事}
に明快な説明がある。}。

\begin{definition}
\label{properties-definition-kappa-generated}
$(X, \mathcal{O}_X)$ を環付き空間とする。$\kappa$ を無限基数とする。
$\mathcal{O}_X$-加群の層 $\mathcal{F}$ が {\it $\kappa$-生成} であるとは、
開被覆 $X = \bigcup U_i$ が存在し、$\mathcal{F}|_{U_i}$ が
部分集合 $R_i \subset \mathcal{F}(U_i)$ により生成され、その濃度が
高々 $\kappa$ であることをいう。
\end{definition}

\noindent
高々 $\kappa$ 個の $\kappa$-生成加群の直和は再び $\kappa$-生成であることに注意する。
これは $\kappa \otimes \kappa = \kappa$ だからである。集合の章の節
\ref{sets-section-cardinals} を参照せよ。特に、二つの $\kappa$-生成加群の直和について成り立つ。
さらに、$\kappa$-生成層の商は $\kappa$-生成である。
（しかし部分加群について同じことが成り立つとは限らない。）

\begin{lemma}
\label{properties-lemma-set-of-iso-classes}
$(X, \mathcal{O}_X)$ を環付き空間とする。$\kappa$ を基数とする。
集合 $T$ と、族 $(\mathcal{F}_t)_{t \in T}$ で $\kappa$-生成
$\mathcal{O}_X$-加群からなるものが存在し、任意の $\kappa$-生成
$\mathcal{O}_X$-加群はある $\mathcal{F}_t$ と同型である。
\end{lemma}

\begin{proof}
$X$ の被覆全体は集合をなす（重複を許さないものとする）。
$X = \bigcup U_i$ を被覆とし、$\mathcal{F}_i$ を
$\mathcal{O}_{U_i}$-加群とする。このとき、貼り合わせ写像全体が集合をなすので、
$\mathcal{O}_X$-加群 $\mathcal{F}$ で、
$\mathcal{F}|_{U_i} \cong \mathcal{F}_i$ という性質を持つものの同型類全体は集合をなす。
従って、商の（同型類の）集合
$\oplus_{k \in \kappa} \mathcal{O}_X \to \mathcal{F}$ が任意の環付き空間 $X$ に対し
存在することを示せばよい。
これは明らかである。
\end{proof}

\noindent
これが、この節の表題が指す結果である。

\begin{lemma}
\label{properties-lemma-colimit-kappa}
$X$ をスキームとする。ある基数 $\kappa$ が存在して、任意の準連接加群
$\mathcal{F}$ は、その準連接な $\kappa$-生成部分加群の有向余極限である。
\end{lemma}

\begin{proof}
アフィン開被覆 $X = \bigcup_{i \in I} U_i$ を選ぶ。各対
$i, j$ に対しアフィン開被覆
$U_i \cap U_j = \bigcup_{k \in I_{ij}} U_{ijk}$ を選ぶ。
$U_i = \Spec(A_i)$ および $U_{ijk} = \Spec(A_{ijk})$ と書く。
$\kappa$ を任意の無限基数で、$\geq$ が集合 $I$, $I_{ij}$ のいずれの濃度との間にも
成り立つものとする。

\medskip\noindent
$\mathcal{F}$ を準連接層とする。$M_i = \mathcal{F}(U_i)$ および
$M_{ijk} = \mathcal{F}(U_{ijk})$ と置く。次に注意する。
$$
M_i \otimes_{A_i} A_{ijk} = M_{ijk} = M_j \otimes_{A_j} A_{ijk}.
$$
スキームの章の補題 \ref{schemes-lemma-widetilde-pullback} を参照せよ。
選択公理を用いて写像
$$
(i, j, k, m) \mapsto S(i, j, k, m)
$$
を選ぶ。これはすべての $i, j \in I$、$k \in I_{ij}$、$m \in M_i$ に対し、
有限部分集合 $S(i, j, k, m) \subset M_j$ を対応させ、$M_{ijk}$ 内で、ある
$a_{m'} \in A_{ijk}$ について
$$
m \otimes 1 = \sum\nolimits_{m' \in S(i, j, k, m)} m' \otimes a_{m'}
$$
が成り立つようにする。さらに、$S(i, i, k, m) = \{m\}$ と、
上記のすべての $i, j = i, k, m$ について約束する。そのような写像を一つ固定する。

\medskip\noindent
部分集合の族 $\mathcal{S} = (S_i)_{i \in I}$ で、
$S_i \subset M_i$ の濃度が高々 $\kappa$ であるものが与えられたとき、
$\mathcal{S}' = (S'_i)$ を次で定める。
$$
S'_j = \bigcup\nolimits_{(i, k, m)\text{ such that }m \in S_i}
S(i, j, k, m)
$$
$S_i \subset S'_i$ に注意する。$S'_i$ は、高々 $\kappa$ の濃度の集合を添字として
有限集合を合併したものなので、その濃度は高々 $\kappa$ である。
$\mathcal{S}^{(0)} = \mathcal{S}$、$\mathcal{S}^{(1)} = \mathcal{S}'$ と置き、帰納的に
$\mathcal{S}^{(n + 1)} = (\mathcal{S}^{(n)})'$ と置く。次に
$\mathcal{S}^{(\infty)} = \bigcup_{n \geq 0} \mathcal{S}^{(n)}$ と置く。
$\mathcal{S}^{(\infty)} = (S^{(\infty)}_i)$ と書くと、任意の元
$m \in S^{(\infty)}_i$ に対し、$m$ の $M_{ijk}$ における像は有限和
$\sum m' \otimes a_{m'}$ と書け、ここで $m' \in S_j^{(\infty)}$ である。
このようにして
$$
N_i = A_i\text{-submodule of }M_i\text{ generated by }S^{(\infty)}_i
$$
と置くと
$$
N_i \otimes_{A_i} A_{ijk} = N_j \otimes_{A_j} A_{ijk}.
$$
が $M_{ijk}$ の部分加群として成り立つ。従って準連接部分層
$\mathcal{G} \subset \mathcal{F}$ で $\mathcal{G}(U_i) = N_i$ を満たすものが存在する。
さらに構成により、層 $\mathcal{G}$ は $\kappa$-生成である。

\medskip\noindent
$\{\mathcal{G}_t\}_{t \in T}$ を $\kappa$-生成準連接部分層の集合とする。
$t, t' \in T$ ならば、$\mathcal{G}_t + \mathcal{G}_{t'}$ も $\kappa$-生成準連接部分層である。
これは写像 $\mathcal{G}_t \oplus \mathcal{G}_{t'} \to \mathcal{F}$ の像だからである。
従って（包含で順序づけた）系は有向である。
上の議論により、$\mathcal{F}$ の $U_i$ 上の任意の切断は、ある
$\mathcal{G}_t$ に入る（$\mathcal{S}$ から始め、与えられた切断が
$S_i$ の元となるようにすればよい）。従って望むように
$\colim_t \mathcal{G}_t \to \mathcal{F}$ は単射かつ全射である。
\end{proof}

\begin{proposition}
\label{properties-proposition-coherator}
$X$ をスキームとする。
\begin{enumerate}
\item 圏 $\QCoh(\mathcal{O}_X)$ は Grothendieck アーベル圏である。
従って $\QCoh(\mathcal{O}_X)$ は十分な入射対象とすべての極限を持つ。
\item 包含関手
$\QCoh(\mathcal{O}_X) \to \textit{Mod}(\mathcal{O}_X)$ は
右随伴を持つ\footnote{この関手は {\it コヒーレーター} と呼ばれることがある。}
$$
Q : \textit{Mod}(\mathcal{O}_X) \longrightarrow \QCoh(\mathcal{O}_X)
$$
そして任意の準連接層 $\mathcal{F}$ に対して随伴写像
$Q(\mathcal{F}) \to \mathcal{F}$ は同型である。
\end{enumerate}
\end{proposition}

\begin{proof}
(1) は、$\QCoh(\mathcal{O}_X)$ が (a) すべての余極限を持ち、
(b) フィルター余極限が完全であり、(c) 生成対象を持つことを意味する。
入射対象の章の節 \ref{injectives-section-grothendieck-conditions} を参照せよ。
スキームの章の節 \ref{schemes-section-quasi-coherent} により、
$\QCoh(\mathcal{O}_X)$ における余極限は存在し、
$\textit{Mod}(\mathcal{O}_X)$ における余極限と一致する。
加群の章の補題 \ref{modules-lemma-limits-colimits} により、
フィルター余極限は完全である。従って (a) と (b) が成り立つ。
生成対象 $U$ を構成するため、補題 \ref{properties-lemma-colimit-kappa} のような基数
$\kappa$ を選ぶ。補題 \ref{properties-lemma-set-of-iso-classes} のような集まり
$(\mathcal{F}_t)_{t \in T}$ を、$\kappa$-生成準連接層から選ぶ。
$U = \bigoplus_{t \in T} \mathcal{F}_t$ と置く。
$\QCoh(\mathcal{O}_X)$ のすべての対象は、$\kappa$-生成準連接加群、すなわち
$\mathcal{F}_t$ と同型な対象のフィルター余極限なので、$U$ が生成対象であることは明らかである。
極限と入射対象に関する主張は任意の Grothendieck アーベル圏で成り立つ。
入射対象の章の定理
\ref{injectives-theorem-injective-embedding-grothendieck} および補題
\ref{injectives-lemma-grothendieck-products} を参照せよ。

\medskip\noindent
(2) の証明。$Q$ を構成するため、次の一般的な手続きを用いる。
$\mathcal{F}$ を $\textit{Mod}(\mathcal{O}_X)$ の対象として与えたとき、関手
$$
\QCoh(\mathcal{O}_X)^{opp} \longrightarrow \textit{Sets},\quad
\mathcal{G} \longmapsto \Hom_X(\mathcal{G}, \mathcal{F})
$$
を考える。この関手は余極限を極限へ移すので、表現可能である。
入射対象の章の補題 \ref{injectives-lemma-grothendieck-brown} を参照せよ。
従って準連接層 $Q(\mathcal{F})$ と関手的同型
$\Hom_X(\mathcal{G}, \mathcal{F}) = \Hom_X(\mathcal{G}, Q(\mathcal{F}))$
が存在する。ここで $\mathcal{G}$ は $\QCoh(\mathcal{O}_X)$ の対象である。
Yoneda の補題（圏の章の補題 \ref{categories-lemma-yoneda}）により、構成
$\mathcal{F} \leadsto Q(\mathcal{F})$ は $\mathcal{F}$ に関して関手的である。
構成により $Q$ は包含関手の右随伴である。
$Q(\mathcal{F}) \to \mathcal{F}$ が同型であることは、$\mathcal{F}$ が準連接ならば、
包含関手 $\QCoh(\mathcal{O}_X) \to \textit{Mod}(\mathcal{O}_X)$ が
充満忠実であることの形式的帰結である。
\end{proof}

\section{閉部分集合に台を持つ切断}
\label{properties-section-sections-with-support-in-closed}

\noindent
任意の位相空間 $X$、閉部分集合 $Z \subset X$、およびアーベル層
$\mathcal{F}$ が与えられれば、台が $Z$ に含まれる切断の部分層を取れる。
$X$ がスキーム、$Z$ が閉部分スキーム、$\mathcal{F}$ が準連接加群ならば、
$Z$ 上にスキーム論的に台を持つ切断を取る変種がある。
ただしスキームの場合、得られる $\mathcal{O}_X$-加群が準連接とは限らないため注意が必要である。

\begin{lemma}
\label{properties-lemma-quasi-coherent-finite-type-ideals}
$X$ を準コンパクトかつ準分離的なスキームとする。
$U \subset X$ を開部分スキームとする。次は同値である。
\begin{enumerate}
\item $U$ は $X$ において逆コンパクトである。
\item $U$ は準コンパクトである。
\item $U$ は有限個のアフィン開集合の合併である。
\item 有限型の準連接イデアル層 $\mathcal{I} \subset \mathcal{O}_X$ が存在し、
$X \setminus U = V(\mathcal{I})$ が（集合論的に）成り立つ。
\end{enumerate}
\end{lemma}

\begin{proof}
(1)、(2)、(3) の同値性は補題
\ref{properties-lemma-quasi-separated-quasi-compact-open-retrocompact} から従う。
(1)、(2)、(3) を仮定する。$T = X \setminus U$ と置く。
スキームの章の補題 \ref{schemes-lemma-reduced-closed-subscheme} により、
一意な準連接イデアル層 $\mathcal{J}$ が存在し、$T$ 上の被約誘導閉部分スキーム構造を切り出す。
$\mathcal{J}|_U = \mathcal{O}_U$ であり、これは有限型の
$\mathcal{O}_U$-加群であることに注意する。補題 \ref{properties-lemma-extend} により、
有限型の準連接部分層 $\mathcal{I} \subset \mathcal{J}$ で、
$\mathcal{I}|_U = \mathcal{J}|_U$ という性質を持つものが存在する。
すると $X \setminus U = V(\mathcal{I})$ であり、(4) を得る。逆に、
$\mathcal{I}$ が (4) のようで、$W = \Spec(R) \subset X$ がアフィン開ならば、
$\mathcal{I}|_W = \widetilde{I}$ となる有限生成イデアル $I \subset R$ が存在する。
補題 \ref{properties-lemma-finite-type-module} を参照せよ。従って
$U \cap W = \Spec(R) \setminus V(I)$ は準コンパクトである。代数の章の補題
\ref{algebra-lemma-qc-open} を参照せよ。ゆえに $U \subset X$ は補題
\ref{properties-lemma-retrocompact} により逆コンパクトである。
\end{proof}

\begin{lemma}
\label{properties-lemma-sections-annihilated-by-ideal}
$X$ をスキームとする。
$\mathcal{I} \subset \mathcal{O}_X$ を準連接イデアル層とする。
$\mathcal{F}$ を準連接 $\mathcal{O}_X$-加群とする。
$\mathcal{O}_X$-加群の層 $\mathcal{F}'$ を、各開集合 $U \subset X$ に
$$
\mathcal{F}'(U)
=
\{s \in \mathcal{F}(U) \mid
\mathcal{I}s = 0\}
$$
を対応させるものとして考える。$\mathcal{I}$ は有限型であると仮定する。このとき
\begin{enumerate}
\item $\mathcal{F}'$ は準連接 $\mathcal{O}_X$-加群の層である。
\item 任意のアフィン開集合 $U \subset X$ 上で
$\mathcal{F}'(U) = \{s \in \mathcal{F}(U) \mid \mathcal{I}(U)s = 0\}$ である。
\item $\mathcal{F}'_x = \{s \in \mathcal{F}_x \mid \mathcal{I}_x s = 0\}$ である。
\end{enumerate}
\end{lemma}

\begin{proof}
$\mathcal{F}'$ を定義する規則が $\mathcal{F}$ の部分層を与えることは明らかである
（層条件は容易に確認できる）。従って他の主張は $X$ 上局所的に確認してよい。
すなわち $X = \Spec(A)$、$\mathcal{F} = \widetilde{M}$、
$\mathcal{I} = \widetilde{I}$ と仮定してよい。この場合、明らかに
$\mathcal{F}'(U) = \{x \in M \mid Ix = 0\} =: M'$ である。
$\widetilde{I}$ はその大域切断 $I$ により生成されるからであり、これで (2) が示される。
$\mathcal{F}'$ が準連接であることを示すには、任意の $f \in A$ に対し
$\{x \in M_f \mid I_f x = 0\} = (M')_f$ を示せば十分である。
$I = (g_1, \ldots, g_t)$ と書く。$\mathcal{I}$ は有限型なのでこれは可能である。
補題 \ref{properties-lemma-finite-type-module} を参照せよ。
$x = y/f^n$ かつ $I_fx = 0$ ならば、任意の $i$ に対し、ある
$m \geq 0$ が存在して $f^mg_ix = 0$ となることを意味する。
一つの $m$ を選び、すべての $i$ に使えるようにできる（ここで $I$ が有限生成であることを用いる）。
すると $f^mx \in M'$ かつ $x/f^n = f^mx/f^{n + m}$ であり、
これは $(M')_f$ 内で成り立つことが分かり、望みどおりである。
(3) の証明も同様なので省略する。
\end{proof}

\begin{definition}
\label{properties-definition-subsheaf-sections-annihilated-by-ideal}
$X$ をスキームとする。
$\mathcal{I} \subset \mathcal{O}_X$ を有限型の準連接イデアル層とする。
$\mathcal{F}$ を準連接 $\mathcal{O}_X$-加群とする。
上の補題 \ref{properties-lemma-sections-annihilated-by-ideal} で定義した部分層
$\mathcal{F}' \subset \mathcal{F}$ を、{\it $\mathcal{I}$ により零化される切断の部分層} と呼ぶ。
\end{definition}

\begin{lemma}
\label{properties-lemma-push-sections-annihilated-by-ideal}
$f : X \to Y$ をスキームの準コンパクトかつ準分離的な射とする。
$\mathcal{I} \subset \mathcal{O}_Y$ を有限型の準連接イデアル層とする。
$\mathcal{F}$ を準連接 $\mathcal{O}_X$-加群とする。
$\mathcal{F}' \subset \mathcal{F}$ を $f^{-1}\mathcal{I}\mathcal{O}_X$ により
零化される切断の部分層とする。すると $f_*\mathcal{F}' \subset f_*\mathcal{F}$ は
$\mathcal{I}$ により零化される切断の部分層である。
\end{lemma}

\begin{proof}
省略する。（ヒント：$f$ が準コンパクトかつ準分離的であるという仮定により
$f_*\mathcal{F}$ は準連接なので、補題 \ref{properties-lemma-sections-annihilated-by-ideal} を
$\mathcal{I}$ と $f_*\mathcal{F}$ に適用できる。）
\end{proof}

\noindent
位相空間上のアーベル層について、閉部分集合に台を持つ切断の部分層を
加群の章の注 \ref{modules-remark-sections-support-in-closed} で論じた。
準連接加群について、この部分加群は必ずしも準連接ではないが、
閉部分集合の補集合が逆コンパクトならば準連接である。

\begin{lemma}
\label{properties-lemma-sections-supported-on-closed-subset}
$X$ をスキームとする。$Z \subset X$ を閉部分集合とする。
$\mathcal{F}$ を準連接 $\mathcal{O}_X$-加群とする。
$\mathcal{O}_X$-加群の層 $\mathcal{F}'$ を、各開集合 $U \subset X$ に
$$
\mathcal{F}'(U)
=
\{s \in \mathcal{F}(U) \mid
\text{the support of }s\text{ is contained in }Z \cap U\}
$$
を対応させるものとして考える。$X \setminus Z$ が $X$ の逆コンパクトな開集合ならば、
\begin{enumerate}
\item アフィン開集合 $U \subset X$ に対し、有限生成イデアル
$I \subset \mathcal{O}_X(U)$ が存在して $Z \cap U = V(I)$ となる。
\item (1) のような $U$ と $I$ に対して
$\mathcal{F}'(U) = \{x \in \mathcal{F}(U) \mid
I^nx = 0 \text{ for some } n\}$ である。
\item $\mathcal{F}'$ は準連接 $\mathcal{O}_X$-加群の層である。
\end{enumerate}
\end{lemma}

\begin{proof}
(1) は代数の章の補題 \ref{algebra-lemma-qc-open} である。
$U = \Spec(A)$ とし、$I$ を (1) のように取る。
すると $\mathcal{F}|_U$ は、ある $A$-加群 $M$ に付随する準連接層である。次が成り立つ。
$$
\mathcal{F}'(U) = \{x \in M \mid x = 0\text{ in }M_\mathfrak p
\text{ for all }\mathfrak p \not \in Z\}.
$$
これは加群の章の定義 \ref{modules-definition-support} による。従って
$x \in \mathcal{F}'(U)$ であることと $V(\text{Ann}(x)) \subset V(I)$ であることは同値である。
代数の章の補題 \ref{algebra-lemma-support-element} を参照せよ。
$I$ は有限生成なので、これは $I^n x = 0$ がある $n$ に対して成り立つことと同値である。
これで (2) が示された。

\medskip\noindent
(3) の証明。開集合 $U \subset X$ が与えられたとき、完全列
$$
0 \to \mathcal{F}'(U) \to \mathcal{F}(U) \to \mathcal{F}(U \setminus Z)
$$
があることに注意する。包含射を $j : X \setminus Z \to X$ と記すと、
$\mathcal{F}(U \setminus Z)$ は加群 $j_*(\mathcal{F}|_{X \setminus Z})$ の
$U$ 上の切断である。従って完全列
$$
0 \to \mathcal{F}' \to \mathcal{F} \to j_*(\mathcal{F}|_{X \setminus Z})
$$
がある。制限 $\mathcal{F}|_{X \setminus Z}$ は準連接である。
従って $j_*(\mathcal{F}|_{X \setminus Z})$ は、スキームの章の補題
\ref{schemes-lemma-push-forward-quasi-coherent} と $j$ が準コンパクトであるという仮定
（任意の開埋め込みは分離的である）により準連接である。
従って $\mathcal{F}'$ は準連接加群の写像の核として準連接である。
スキームの章の節 \ref{schemes-section-quasi-coherent} を参照せよ。
\end{proof}

\begin{definition}
\label{properties-definition-subsheaf-sections-supported-on-closed}
$X$ をスキームとする。
$T \subset X$ を、その補集合が $X$ において逆コンパクトである閉部分集合とする。
$\mathcal{F}$ を準連接 $\mathcal{O}_X$-加群とする。
補題 \ref{properties-lemma-sections-supported-on-closed-subset} で定義した準連接部分層
$\mathcal{F}' \subset \mathcal{F}$ を、{\it $T$ 上に台を持つ切断の部分層} と呼ぶ。
\end{definition}

\begin{lemma}
\label{properties-lemma-push-sections-supported-on-closed-subset}
$f : X \to Y$ をスキームの準コンパクトかつ準分離的な射とする。
$Z \subset Y$ を閉部分集合とし、$Y \setminus Z$ は $Y$ において逆コンパクトであるとする。
$\mathcal{F}$ を準連接 $\mathcal{O}_X$-加群とする。
$\mathcal{F}' \subset \mathcal{F}$ を $f^{-1}Z$ に台を持つ切断の部分層とする。
すると $f_*\mathcal{F}' \subset f_*\mathcal{F}$ は $Z$ に台を持つ切断の部分層である。
\end{lemma}

\begin{proof}
省略する。（ヒント：まず $X \setminus f^{-1}Z$ が $X$ において逆コンパクトであることを、
$Y \setminus Z$ が $Y$ において逆コンパクトであることから示す。
従って補題 \ref{properties-lemma-sections-supported-on-closed-subset} を
$f^{-1}Z$ と $\mathcal{F}$ に適用できる。$f$ は準コンパクトかつ準分離的なので
$f_*\mathcal{F}$ は準連接である。従って補題
\ref{properties-lemma-sections-supported-on-closed-subset} を $Z$ と $f_*\mathcal{F}$ に適用できる。
最後に層を直接照合する。）
\end{proof}

\section{準連接層の切断}
\label{properties-section-sections-quasi-coherent}

\noindent
ここでは、アフィンスペクトルの準コンパクト開集合上の準連接層の切断を計算する。

\begin{lemma}
\label{properties-lemma-sections-over-quasi-compact-open-in-affine}
$A$ を環とする。
$I \subset A$ を有限生成イデアルとする。
$M$ を $A$-加群とする。
このとき標準写像
$$
\colim_n \Hom_A(I^n, M)
\longrightarrow
\Gamma(\Spec(A) \setminus V(I), \widetilde{M}).
$$
が存在する。この写像は常に単射である。
すべての $x \in M$ に対して $Ix = 0 \Rightarrow x = 0$ ならば、
この写像は同型である。一般には
$M_n = \{x \in M \mid I^nx = 0\}$ とおくと、同型
$$
\colim_n \Hom_A(I^n, M/M_n)
\longrightarrow
\Gamma(\Spec(A) \setminus V(I), \widetilde{M}).
$$
が存在する。
\end{lemma}

\begin{proof}
$I^{n + 1} \subset I^n$ かつ $M_n \subset M_{n + 1}$ なので、
これらの写像との合成により標準的な $A$-加群の写像
$$
\Hom_A(I^n, M)
\longrightarrow
\Hom_A(I^{n + 1}, M)
$$
および
$$
\Hom_A(I^n, M/M_n)
\longrightarrow
\Hom_A(I^{n + 1}, M/M_{n + 1})
$$
を得る。これらを各系の遷移写像として用いる。
$A$-加群の写像 $\varphi : I^n \to M$ が与えられると、層の写像
$\widetilde{\varphi} : \widetilde{I^n} \to \widetilde{M}$ が得られ、
これを開集合 $\Spec(A) \setminus V(I)$ に制限できる。
$\widetilde{I^n}$ をこの開集合に制限すると構造層になるので、
$\Gamma(\Spec(A) \setminus V(I), \widetilde{M})$ の元を得る。
これが系 $\Hom_A(I^n, M)$ の遷移写像と両立することの確認は省略する。
これにより最初の写像が得られる。二つ目の写像については、
$\widetilde{M}$ と $\widetilde{M/M_n}$ が開集合
$\Spec(A) \setminus V(I)$ 上で一致することに注意する。実際、層
$\widetilde{M_n}$ は明らかに $V(I)$ 上に台を持つ。
従って前と同じ仕組みを用いることができる。

\medskip\noindent
次に、この写像を代数的に定義する方法を具体的に述べる。
$I = (f_1, \ldots, f_t)$ とする。このとき
$\Spec(A) \setminus V(I) = \bigcup_{i = 1, \ldots, t} D(f_i)$ である。
従って
$$
0 \to
\Gamma(\Spec(A) \setminus V(I), \widetilde{M}) \to
\bigoplus\nolimits_i M_{f_i} \to
\bigoplus\nolimits_{i, j} M_{f_if_j}
$$
は完全である。$\varphi : I^n \to M$ を $A$-加群の写像とする。
元のベクトル $\varphi(f_i^n)/f_i^n \in M_{f_i}$ を考える。
このベクトルが上の完全列の二つ目の直和において零へ写ることは容易に分かる。
従って $\Gamma(\Spec(A) \setminus V(I), \widetilde{M})$ の元を得る。
この記述が上の記述と一致することの確認は省略する。

\medskip\noindent
この記述を用いて最初の写像が単射であることを示そう。
実際、$\varphi$ が零に写るなら、各 $i$ について元
$\varphi(f_i^n)/f_i^n$ は $M_{f_i}$ において零である。言い換えると、
各 $i$ について $f_i^m\varphi(f_i^n) = 0$ が、ある
$m \geq 0$ に対して成り立つ。一つの $m$ を選び、これがすべての
$i$ に対して使えるようにできる。このとき $\varphi(f_i^{n + m}) = 0$ であり、これは
すべての $i$ について成り立つ。また、これは容易に
$\varphi|_{I^{t(n + m - 1) + 1}} = 0$ を意味する。言い換えると、
$\varphi$ は余極限の第 $t(n + m - 1) + 1$ 項において零に写る。
従って単射性が従う。

\medskip\noindent
各 $M_n = 0$ であることに注意する。ただしこれは
$Ix = 0 \Rightarrow x = 0$ が $x \in M$ に対して成り立つ場合である。
従って補題の証明を終えるには、
二つ目の写像が同型であることを示せば十分である。

\medskip\noindent
補題の二つ目の写像の逆写像を構成してみよう。
$s \in \Gamma(\Spec(A) \setminus V(I), \widetilde{M})$ とする。
これは上の完全列の最初の直和に属するベクトル
$x_i/f_i^n$ に対応し、ここで $x_i \in M$ である。
従って各 $i, j$ に対して $m \geq 0$ が存在し、
$f_i^m f_j^m (f_j^n x_i - f_i^n x_j) = 0$ が $M$ において成り立つ。
すべての組に対して使える一つの $m$ を選ぶことができる。
ここでいう組は $i, j$ のすべての組である。
$x_i$ を $f_i^mx_i$ で、$n$ を
$n + m$ で置き換えると、$f_j^nx_i = f_i^nx_j$ が $M$ において成り立ち、これはすべての
$i, j$ に対する主張である。ここで
$$
K_n = \{x \in M \mid f_1^nx = \ldots = f_t^nx = 0\}
$$
とおく。$A$-加群の写像
$$
\varphi :
I^{t(n - 1) + 1}
\longrightarrow
M/K_n
$$
であって、単項式 $f_1^{e_1} \ldots f_t^{e_t}$（ここで
$\sum e_i = t(n - 1) + 1$）を、$K_n$ を法とする式
$f_1^{e_1} \ldots f_i^{e_i - n} \ldots f_t^{e_t}x_i$ の類へ
写すものが存在すると主張する。ここで $i$ は $e_i \geq n$ を満たすように選ぶ
（そのような $i$ は少なくとも一つ存在する）。
これを確かめるために、
$$
\sum\nolimits_{E = (e_1, \ldots, e_t), |E| = t(n - 1) + 1}
a_E f_1^{e_1} \ldots f_t^{e_t} = 0
$$
を、係数 $a_E$ が $A$ に属するこれら単項式の間の関係とする。
この関係は
$$
z =
\sum\nolimits_{E = (e_1, \ldots, e_t), |E| = t(n - 1) + 1}
a_E f_1^{e_1} \ldots f_{i(E)}^{e_{i(E)} - n} \ldots f_t^{e_t}x_{i(E)}
$$
へ写る。ここで各多重指数 $E$ に対して特定の $i(E)$ を、
$e_{i(E)} \geq n$ を満たすように選んでいる。これを $f_j^n$ 倍すると
零になることに注意する。ここで $j$ は任意である。
実際、上の関係 $f_j^nx_i = f_i^nx_j$ により
\begin{align*}
f_j^nz & = \sum\nolimits_{E = (e_1, \ldots, e_t), |E| = t(n - 1) + 1}
a_E f_1^{e_1} \ldots f_j^{e_j + n}
\ldots f_{i(E)}^{e_{i(E)} - n} \ldots f_t^{e_t}x_{i(E)} \\
& =
\sum\nolimits_{E = (e_1, \ldots, e_t), |E| = t(n - 1) + 1}
a_E f_1^{e_1} \ldots f_t^{e_t}x_j
= 0.
\end{align*}
を得る。従って $z \in K_n$ であり、すべての関係が $M/K_n$ において
零へ写ることが分かる。これで主張が証明された。

\medskip\noindent
$K_n \subset M_{t(n - 1) + 1}$ であることに注意する。従って写像
$\varphi$ は特に $A$-加群の写像
$I^{t(n - 1) + 1} \to M/M_{t(n - 1) + 1}$ を誘導する。
これにより補題の二つ目の写像が全射であることが証明された。
単射性の証明は省略する。
\end{proof}

\begin{example}
\label{properties-example-does-not-work-in-general}
補題 \ref{properties-lemma-sections-over-quasi-compact-open-in-affine} の最初の表示写像が
同型でないことを示す二つの例を挙げる。

\medskip\noindent
$k$ を体とする。環
$$
A = k[x, y, z_1, z_2, \ldots]/(x^nz_n).
$$
を考える。$I = (x)$ とおき、$M = A$ とする。このとき元 $y/x$ は、
$\Spec(A)$ の構造層の $D(x) = \Spec(A) \setminus V(I)$ 上の切断を定める。
$y/x$ は標準写像
$\colim \Hom_A(I^n, A) \to A_x = \mathcal{O}(D(x))$ の像に属さないと主張する。
実際、像に属するなら、準同型 $\varphi : I^n \to A$ から、ある
$n$ について生じるはずである。$a = \varphi(x^n)$ とおく。
このとき $x^m(xa - x^ny) = 0$ となるような、ある $m > 0$ が存在するはずである。
従って $x^{m + 1}a = x^{m + n}y$ となり、さらに
$\varphi(x^{n + m + 1}) = x^{m + n}y$ となる。これは
$$
0 = \varphi(0) = \varphi(z_{n + m + 1} x^{n + m + 1}) =
x^{m + n}y z_{n + m + 1}
$$
を意味するので矛盾である。この等式は環 $A$ において成り立たない。

\medskip\noindent
$k$ を体とする。環
$$
A = k[f, g, x, y, \{a_n, b_n\}_{n \geq 1}]/
(fy - gx, \{a_nf^n + b_ng^n\}_{n \geq 1}).
$$
を考える。$I = (f, g)$ とおき、$M = A$ とする。このとき
$x/f \in A_f$ と $y/g \in A_g$ は $A_{fg}$ の同じ元へ写る。
従って、これらは切断 $s$、すなわち $\Spec(A)$ の構造層の
$D(f) \cup D(g) = \Spec(A) \setminus V(I)$ 上の切断を定める。しかし、補題
\ref{properties-lemma-sections-over-quasi-compact-open-in-affine} の最初の表示写像の
始域に関して、$n \geq 0$ であって、$s$ が $A$-加群の写像
$\varphi : I^n \to A$ から生じるようなものは存在しない。
実際、そのような加群の写像が与えられたとして
$x_n = \varphi(f^n)$ および $y_n = \varphi(g^n)$ とおく。
このとき $f^mx_n = f^{n + m - 1}x$ および
$g^my_n = g^{n + m - 1}y$ が、ある $m \geq 0$ に対して成り立つ
（補題の証明を参照）。しかしこのとき
$0 = \varphi(0) =
\varphi(a_{n + m}f^{n + m} + b_{n + m}g^{n + m}) =
a_{n + m}f^{n + m - 1}x + b_{n + m}g^{n + m - 1}y$ となるはずだが、
これは環 $A$ において成り立たない。
\end{example}

\noindent
Noether の場合には次の補題を改良する。《スキームのコホモロジー》の補題
\ref{coherent-lemma-homs-over-open} を参照せよ。

\begin{lemma}
\label{properties-lemma-sections-over-quasi-compact-open}
$X$ を準コンパクトスキームとする。
$\mathcal{I} \subset \mathcal{O}_X$ を有限型の準連接イデアル層とする。
$Z \subset X$ を $\mathcal{I}$ により定まる閉部分スキームとし、
$U = X \setminus Z$ とおく。
$\mathcal{F}$ を準連接 $\mathcal{O}_X$-加群とする。
標準写像
$$
\colim_n \Hom_{\mathcal{O}_X}(\mathcal{I}^n,
\mathcal{F})
\longrightarrow
\Gamma(U, \mathcal{F})
$$
は単射である。さらに $X$ が準分離的であると仮定する。
$\mathcal{F}_n \subset \mathcal{F}$ を $\mathcal{I}^n$ により零化される
切断の部分層とする。標準写像
$$
\colim_n \Hom_{\mathcal{O}_X}(\mathcal{I}^n,
\mathcal{F}/\mathcal{F}_n)
\longrightarrow
\Gamma(U, \mathcal{F})
$$
は同型である。
\end{lemma}

\begin{proof}
$\Spec(A) = W \subset X$ をアフィン開集合とする。
$\mathcal{F}|_W = \widetilde{M}$ と書く。この表示では $A$-加群 $M$ を用いる。
また $\mathcal{I}|_W = \widetilde{I}$ と書く。ここで $I \subset A$ は
有限型イデアルである。
補題の最初の表示写像を $W$ に制限すると、補題
\ref{properties-lemma-sections-over-quasi-compact-open-in-affine} の最初の表示写像を得る。
$X$ は有限個のアフィン開集合で被覆できるので、これにより補題の最初の
表示写像が単射であることが証明される。

\medskip\noindent
$\mathcal{F}_n|_W = \widetilde{M_n}$ であり、ここで
$M_n \subset M$ は補題
\ref{properties-lemma-sections-over-quasi-compact-open-in-affine} と同様に定義される
（詳細は省略する）。この補題により、任意のそのようなアフィン開集合 $W$ に対して
全単射
$$
\colim_n  \Hom_{\mathcal{O}_W}(
\mathcal{I}^n|_W, (\mathcal{F}/\mathcal{F}_n)|_W)
\longrightarrow
\Gamma(U \cap W, \mathcal{F})
$$
が存在する。

\medskip\noindent
補題の二つ目の表示写像が全単射であることを示すため、有限アフィン開被覆
$X = \bigcup_{j = 1, \ldots, m} W_j$ を選ぶ。単射性は上の結果と被覆の
有限性から直ちに従う。$X$ が準分離的ならば、各組 $j, j'$ に対して
有限アフィン開被覆
$$
W_j \cap W_{j'} = \bigcup\nolimits_{k = 1, \ldots, m_{jj'}} W_{jj'k}.
$$
を選ぶ。$s \in \Gamma(U, \mathcal{F})$ とする。上で見たように、各 $j$ に対して
$n_j$ と写像
$\varphi_j : \mathcal{I}^{n_j}|_{W_j} \to
(\mathcal{F}/\mathcal{F}_{n_j})|_{W_j}$
が存在し、これは $s|_{U \cap W_j}$ に対応する。
同様に、各三つ組 $(j, j', k)$ に対して整数 $n_{jj'k}$ が存在し、
$\varphi_j$ と $\varphi_{j'}$ を写像
$\mathcal{I}^{n_{jj'k}} \to \mathcal{F}/\mathcal{F}_{n_{jj'k}}$ とみなしたときの制限は
$W_{jj'k}$ 上で一致する。$n = \max\{n_j, n_{jj'k}\}$ とおくと、
$\varphi_j$ は写像 $\mathcal{I}^n \to \mathcal{F}/\mathcal{F}_n$ として
$X$ 上で貼り合わさる。
これにより写像の全射性が証明された。
\end{proof}




\section{豊富な可逆層}
\label{properties-section-ample}

\noindent
加群の章の補題 \ref{modules-lemma-s-open} を思い出そう。
可逆層 $\mathcal{L}$ が局所環付き空間 $X$ 上に与えられ、さらに大域切断
$s$ が $\mathcal{L}$ に与えられると、集合
$X_s = \{x \in X \mid s \not \in \mathfrak m_x\mathcal{L}_x\}$ は開である。
一般に $X_s \cap X_{s'} = X_{ss'}$ である。ここで $ss'$ は切断
$s \otimes s' \in \Gamma(X, \mathcal{L} \otimes \mathcal{L}')$ を表す。

\begin{definition}
\label{properties-definition-ample}
\begin{reference}
\cite[II Definition 4.5.3]{EGA}
\end{reference}
$X$ をスキームとする。
$\mathcal{L}$ を可逆 $\mathcal{O}_X$-加群とする。
$\mathcal{L}$ が {\it 豊富} であるとは、次の条件が成り立つことをいう。
\begin{enumerate}
\item $X$ は準コンパクトである。
\item すべての $x \in X$ に対して、$n \geq 1$ と
$s \in \Gamma(X, \mathcal{L}^{\otimes n})$ が存在し、
$x \in X_s$ かつ $X_s$ はアフィンである。
\end{enumerate}
\end{definition}

\begin{lemma}
\label{properties-lemma-ample-power-ample}
\begin{reference}
\cite[II Proposition 4.5.6(i)]{EGA}
\end{reference}
$X$ をスキームとする。$\mathcal{L}$ を可逆 $\mathcal{O}_X$-加群とする。
$n \geq 1$ とする。このとき $\mathcal{L}$ が豊富であることと、
$\mathcal{L}^{\otimes n}$ が豊富であることは同値である。
\end{lemma}

\begin{proof}
$X_{s^n} = X_s$ であることから従う。
\end{proof}

\begin{lemma}
\label{properties-lemma-ample-on-closed}
$X$ をスキームとする。
$\mathcal{L}$ を豊富な可逆 $\mathcal{O}_X$-加群とする。
任意の閉部分スキーム $Z \subset X$ に対して、$\mathcal{L}$ の $Z$ への制限は豊富である。
\end{lemma}

\begin{proof}
準コンパクト空間の閉部分集合は準コンパクトであり、アフィンスキームの閉部分スキームは
アフィンであるから明らかである（スキームの章の補題
\ref{schemes-lemma-closed-immersion-affine-case} を参照）。
\end{proof}

\begin{lemma}
\label{properties-lemma-affine-cap-s-open}
$X$ をスキームとする。$\mathcal{L}$ を可逆 $\mathcal{O}_X$-加群とする。
$s \in \Gamma(X, \mathcal{L})$ とする。任意のアフィン開集合 $U \subset X$ に対して、
共通部分 $U \cap X_s$ はアフィンである。
\end{lemma}

\begin{proof}
これは次の代数の問題に翻訳される。
$R$ を環とする。$N$ を可逆 $R$-加群
（すなわち階数 1 の局所自由加群）とする。$s \in N$ を元とする。
このとき $V = \{\mathfrak p \mid s \not \in \mathfrak p N\}$ は
$\Spec(R)$ のアフィン開部分集合である。

\medskip\noindent
$A = \bigoplus_{n \geq 0} A_n$ を $N$ の対称代数（可換である）とし、
$s$ を $A_1$ の元とみなす。$B = A/(s - 1)A$ とおく。
これは $R$-代数であり、その構成は任意の基底変換 $R \to R'$ と可換である。
従って $B' = B \otimes_R R'$ は、$s$ が $N' = N \otimes_R R'$ において
零へ写るならば零環である。これより、$x \in \Spec(R) \setminus V$ ならば
$B \otimes_R \kappa(x) = 0$ である。従って
$\Spec(B) \to \Spec(R)$ は $V$ を経由する。ここでファイバーが空なのは
$x \not \in V$ に対してである。一方、
$\Spec(R') \subset V$ がアフィン開集合ならば、$s$ は $N'$ の基底元に写り、
$B' = R'[s]/(s - 1) \cong R'$ となる。従って
$\Spec(B) \to V$ は同型であり、$V$ は確かにアフィンである。
\end{proof}

\begin{lemma}
\label{properties-lemma-ample-tensor-globally-generated}
\begin{reference}
\cite[II Proposition 4.5.6(ii)]{EGA}
\end{reference}
$X$ をスキームとする。$\mathcal{L}$ と $\mathcal{M}$ を可逆
$\mathcal{O}_X$-加群とする。次を仮定する。
\begin{enumerate}
\item $\mathcal{L}$ は豊富である。
\item 開集合 $X_t$（ここで $t \in \Gamma(X, \mathcal{M}^{\otimes m})$ かつ
$m > 0$）は $X$ を被覆する。
\end{enumerate}
このとき $\mathcal{L} \otimes \mathcal{M}$ は豊富である。
\end{lemma}

\begin{proof}
定義 \ref{properties-definition-ample} の条件を確認する。
$\mathcal{L}$ は豊富なので $X$ は準コンパクトである。
$x \in X$ とする。$n \geq 1$、$m \geq 1$、
$s \in \Gamma(X, \mathcal{L}^{\otimes n})$、および
$t \in \Gamma(X, \mathcal{M}^{\otimes m})$ を、
$x \in X_s$、$x \in X_t$ かつ $X_s$ がアフィンとなるように選ぶ。
このとき $s^mt^n \in \Gamma(X, (\mathcal{L} \otimes \mathcal{M})^{\otimes nm})$、
$x \in X_{s^mt^n}$ であり、補題 \ref{properties-lemma-affine-cap-s-open} により
$X_{s^mt^n}$ はアフィンである。
\end{proof}

\begin{lemma}
\label{properties-lemma-affine-s-opens}
$X$ をスキームとする。$\mathcal{L}$ を可逆 $\mathcal{O}_X$-加群とする。
$X_s$ という形の開集合（ここで $s \in \Gamma(X, \mathcal{L}^{\otimes n})$ かつ
$n \geq 1$）が $X$ の位相の基底をなすと仮定する。このとき、それらのうちアフィンである
開集合 $X_s$ は $X$ の位相の基底をなす。
\end{lemma}

\begin{proof}
$x \in X$ とする。アフィン開近傍
$\Spec(R) = U \subset X$ を $x$ に対して選ぶ。仮定により、
$n \geq 1$ と $s \in \Gamma(X, \mathcal{L}^{\otimes n})$ が存在して
$X_s \subset U$ となる。上の補題
\ref{properties-lemma-affine-cap-s-open} により、共通部分 $X_s = U \cap X_s$ はアフィンである。
$U$ はいくらでも小さく選べるので結論を得る。
\end{proof}

\begin{lemma}
\label{properties-lemma-affine-s-opens-cover-quasi-separated}
$X$ をスキームとし、$\mathcal{L}$ を可逆 $\mathcal{O}_X$-加群とする。
$x$ を $X$ の任意の点とする。この点に対して $n \geq 1$ と
$s \in \Gamma(X, \mathcal{L}^{\otimes n})$ が存在し、
$x \in X_s$ かつ $X_s$ はアフィンであると仮定する。このとき $X$ は分離的である。
\end{lemma}

\begin{proof}
まず $X$ が準分離的であることを示す。仮定により、$X$ は
$X_s$ という形のアフィン開集合で被覆できる。補題
\ref{properties-lemma-affine-cap-s-open} により、このような二つの集合の共通部分はアフィンである。
従ってスキームの章の補題
\ref{schemes-lemma-characterize-quasi-separated} から $X$ は準分離的である。

\medskip\noindent
$X$ が分離的であることを示すため、スキームの章の補題
\ref{schemes-lemma-valuative-criterion-separatedness} の付値判定法を用いる。
$A$ を分数体 $K$ を持つ付値環とし、二つの射
$f, g : \Spec(A) \to X$ であって、二つの合成
$\Spec(K) \to \Spec(A) \to X$ が一致するものを考える。
$A$ は局所環なので、$p, q \ge 1$、
$s \in \Gamma(X,
\mathcal{L}^{\otimes p})$、および
$t \in \Gamma(X,
\mathcal{L}^{\otimes q})$ が存在し、
$X_s$ と $X_t$ はアフィンで、$f(\Spec A) \subseteq X_s$、かつ
$g(\Spec A)
\subseteq X_t$ となる。ここで $s$ を $s^q$ に、$t$ を $t^p$ に、
$\mathcal{L}$ を $\mathcal{L}^{\otimes pq}$ に置き換える。
$X_s = X_{s^q}$ かつ $X_t = X_{t^p}$ なのでこれは無害であり、いまや $s$ と
$t$ は同じ層 $\mathcal{L}$ の切断である。

\medskip\noindent
準連接加群 $f^*\mathcal{L}$ は $A$-加群 $M$ に対応し、
$g^*\mathcal{L}$ は $A$-加群 $N$ に対応する。これはアフィンスキーム上の
準連接加群の分類による（スキームの章の補題
\ref{schemes-lemma-quasi-coherent-affine}）。$A$-加群 $M$ と $N$ は
階数 $1$ の局所自由加群であり（補題 \ref{properties-lemma-locally-free-module}）、
$A$ は局所環なので自由である（代数の章の補題
\ref{algebra-lemma-K0-local}）。従って $M$ と $N$ をそれぞれ $A$-部分加群として
$M \otimes_A K$ と $N
\otimes_A K$ の中に同一視できる。等式
$f|_{\Spec(K)} = g|_{\Spec(K)}$ は同型
$\phi \colon M \otimes_A K \to N
\otimes_A K$ を定める。

\medskip\noindent
$x \in M$ と $y \in N$ を、それぞれ $s$ の $f$ および $g$ に沿う引き戻しに
対応する元とする。これらは $\phi(x \otimes 1) = y \otimes 1$ を満たす。
$f$ の像は $X_s$ に含まれるので $x \not\in \mathfrak{m}_A M$、すなわち
$x$ は $M$ を生成する。従って $\phi$ は $M$ と、$N$ のうち $y$ により生成される
部分加群との同型を定める。$t$ を用いて対称的に論じると、
$\phi^{-1}$ は $N$ と $M$ のある部分加群との同型を定める。
従って $\phi$ は $M$ と $N$ の同型に制限される。$x$ は $M$ を生成するので、
その像 $y$ は $N$ を生成し、$y \not\in \mathfrak{m}_A N$ となる。
従って $g(\Spec(A)) \subseteq X_s$ である。$X_s$ はアフィンなので、
スキームの章の補題 \ref{schemes-lemma-affine-separated} により分離的である。
従って $f = g$ を得る。
\end{proof}

\begin{lemma}
\label{properties-lemma-ample-separated}
$X$ をスキームとする。$X$ 上に豊富な可逆層が存在するならば、$X$ は分離的である。
\end{lemma}

\begin{proof}
補題 \ref{properties-lemma-affine-s-opens-cover-quasi-separated} と定義
\ref{properties-definition-ample} から直ちに従う。
\end{proof}

\begin{lemma}
\label{properties-lemma-map-into-proj}
$X$ をスキームとする。
$\mathcal{L}$ を可逆 $\mathcal{O}_X$-加群とする。
$S = \Gamma_*(X, \mathcal{L})$ を次数付き環とする。
$X$ の各点が開部分スキーム $X_s$ のいずれかに含まれ、ここで
$s \in S_{+}$ は斉次元であるとする。このとき斉次スペクトル $S$ へのスキームの標準射
$$
f : X \longrightarrow Y = \text{Proj}(S),
$$
が存在する（構成の章の節 \ref{constructions-section-proj} を参照）。
この射は次の性質を持つ。
\begin{enumerate}
\item $f^{-1}(D_{+}(s)) = X_s$ は任意の斉次元 $s \in S_{+}$ に対して成り立つ。
\item 乗法写像と両立する $\mathcal{O}_X$-加群の写像
$f^*\mathcal{O}_Y(n) \to \mathcal{L}^{\otimes n}$ が存在する。構成の章の式
(\ref{constructions-equation-multiply}) を参照せよ。
\item 合成
$S_n \to \Gamma(Y, \mathcal{O}_Y(n)) \to \Gamma(X, \mathcal{L}^{\otimes n})$
は恒等写像である。
\item すべての $x \in X$ に対して、整数 $d \geq 1$ と開近傍
$U \subset X$ が存在し、これは $x$ の近傍である。また、写像
$f^*\mathcal{O}_Y(dn)|_U \to \mathcal{L}^{\otimes dn}|_U$ はすべての
$n \in \mathbf{Z}$ に対して同型である。
\end{enumerate}
\end{lemma}

\begin{proof}
$\psi : S \to \Gamma_*(X, \mathcal{L})$ を恒等写像とする。
構成の章の補題 \ref{constructions-lemma-invertible-map-into-proj} の三つ組
$(U(\psi), r_{\mathcal{L}, \psi}, \theta)$ を用いる。
仮定により開部分スキーム $U(\psi)$ は $X$ に等しい。従って
$r_{\mathcal{L}, \psi} : U(\psi) \to Y$ は $X$ 全体で定義される。
$f = r_{\mathcal{L}, \psi}$ とおく。(2) の写像は $\theta$ の成分である。
(3) は上で引用した補題の条件 (2) から従う。(1) は (3) と同補題の条件 (1) を
組み合わせて従う。(4) は構成の章の補題
\ref{constructions-lemma-invertible-map-into-proj} の最後の主張から従う。
実際、そこで述べられる写像 $\alpha$ は同型である。
\end{proof}

\begin{lemma}
\label{properties-lemma-map-into-proj-quasi-compact}
$X$ をスキームとする。$\mathcal{L}$ を可逆 $\mathcal{O}_X$-加群とする。
$S = \Gamma_*(X, \mathcal{L})$ とおく。
(a) $X$ の各点が開部分スキーム $X_s$ のいずれかに含まれ、ここで
$s \in S_{+}$ は斉次元であり、(b) $X$ は準コンパクトであると仮定する。このとき、上の補題
\ref{properties-lemma-map-into-proj} のスキームの標準射
$f : X \longrightarrow \text{Proj}(S)$ は準コンパクトで、像は稠密である。
\end{lemma}

\begin{proof}
$f$ が準コンパクトであることを示すには、$f^{-1}(D_{+}(s))$ が
任意の斉次元 $s \in S_{+}$ に対して準コンパクトであることを示せば十分である。
$X = \bigcup_{i = 1, \ldots, n} X_i$ をアフィン開集合の有限合併として書く。
補題 \ref{properties-lemma-affine-cap-s-open} により、各共通部分 $X_s \cap X_i$ はアフィンである。
従って $X_s = \bigcup_{i = 1, \ldots, n} X_s \cap X_i$ は準コンパクトである。
$f$ の像が稠密でないと仮定して矛盾を導く。開集合 $D_+(s)$ は、
斉次元 $s \in S_+$ にわたって $\text{Proj}(S)$ の位相の基底をなすので、
そのような $s$ であって、$D_+(s) \not = \emptyset$ かつ
$f(X) \cap D_+(s) = \emptyset$ を満たすものを選べる。
補題 \ref{properties-lemma-map-into-proj} により、これは
$X_s = \emptyset$ を意味する。補題 \ref{properties-lemma-invert-s-sections} により、これは
ある冪 $s^n$ が $\mathcal{L}^{\otimes n\deg(s)}$ の零切断であることを意味する。
さらにこれは $D_+(s) = \emptyset$ を意味し、求める矛盾である。
\end{proof}

\begin{lemma}
\label{properties-lemma-ample-immersion-into-proj}
$X$ をスキームとする。$\mathcal{L}$ を可逆 $\mathcal{O}_X$-加群とする。
$S = \Gamma_*(X, \mathcal{L})$ とおく。
$\mathcal{L}$ が豊富であると仮定する。このとき、補題
\ref{properties-lemma-map-into-proj} のスキームの標準射
$f : X \longrightarrow \text{Proj}(S)$ は稠密な像を持つ開埋め込みである。
\end{lemma}

\begin{proof}
補題 \ref{properties-lemma-affine-s-opens-cover-quasi-separated} により、$X$ は準分離的である。
斉次元 $s_1, \ldots, s_n \in S_{+}$ を、$X_{s_i}$ がアフィンで、かつ
$X = \bigcup X_{s_i}$ となるように有限個選ぶ。$s_i$ の次数を $d_i$ とする。
$D_{+}(s_i)$ の $f$ による逆像は $X_{s_i}$ である。補題
\ref{properties-lemma-map-into-proj} を参照せよ。補題 \ref{properties-lemma-invert-s-sections} により環写像
$$
(S^{(d_i)})_{(s_i)} = \Gamma(D_{+}(s_i), \mathcal{O}_{\text{Proj}(S)})
\longrightarrow
\Gamma(X_{s_i}, \mathcal{O}_X)
$$
は同型である。従って $f$ は同型 $X_{s_i} \to D_{+}(s_i)$ を誘導する。
ゆえに $f$ は $X$ から開部分スキーム
$\bigcup_{i = 1, \ldots, n} D_{+}(s_i)$ への同型であり、この開部分スキームは
$\text{Proj}(S)$ に含まれる。
像は補題 \ref{properties-lemma-map-into-proj-quasi-compact} により稠密である。
\end{proof}

\begin{lemma}
\label{properties-lemma-open-in-proj-ample}
$X$ をスキームとする。
$S$ を次数付き環とする。$X$ が準コンパクトであり、開埋め込み
$$
j : X \longrightarrow Y = \text{Proj}(S).
$$
が存在すると仮定する。このとき $j^*\mathcal{O}_Y(d)$ は、ある
$d > 0$ に対して可逆な豊富層である。
\end{lemma}

\begin{proof}
これは構成の章の補題 \ref{constructions-lemma-ample-on-proj} である。
\end{proof}

\begin{proposition}
\label{properties-proposition-characterize-ample}
$X$ を準コンパクトスキームとする。
$\mathcal{L}$ を $X$ 上の可逆層とする。
$S = \Gamma_*(X, \mathcal{L})$ とおく。
次の条件は同値である。
\begin{enumerate}
\item
\label{properties-item-ample}
$\mathcal{L}$ は豊富である。
\item
\label{properties-item-immersion}
開集合 $X_s$（ここで $s \in S_{+}$ は斉次元）が $X$ を被覆し、付随する射
$X \to \text{Proj}(S)$ は開埋め込みである。
\item
\label{properties-item-s-basis}
開集合 $X_s$（ここで $s \in S_{+}$ は斉次元）は $X$ の位相の基底をなす。
\item
\label{properties-item-s-affine-basis}
開集合 $X_s$（ここで $s \in S_{+}$ は斉次元）のうち、アフィンであるものは
$X$ の位相の基底をなす。
\item
\label{properties-item-qc-gg}
任意の準連接層 $\mathcal{F}$ が $X$ 上に与えられたとき、$n \geq 1$ による標準写像
$$
\Gamma(X, \mathcal{F} \otimes_{\mathcal{O}_X} \mathcal{L}^{\otimes n})
\otimes_{\mathbf{Z}} \mathcal{L}^{\otimes -n}
\longrightarrow
\mathcal{F}
$$
の像の和は $\mathcal{F}$ に等しい。
\item
\label{properties-item-qc-i-gg}
(\ref{properties-item-qc-gg}) と同じ性質が、$\mathcal{F}$ をすべての準連接イデアル層に
わたって動かしたときに成り立つ。
\item
\label{properties-item-c-gg}
$X$ は準分離的であり、任意の有限型準連接層 $\mathcal{F}$ が $X$ 上に与えられたとき、
整数 $n_0$ が存在し、
$\mathcal{F} \otimes_{\mathcal{O}_X} \mathcal{L}^{\otimes n}$ は、すべての
$n \geq n_0$ に対して大域生成である。
\item
\label{properties-item-c-q}
$X$ は準分離的であり、任意の有限型準連接層 $\mathcal{F}$ が $X$ 上に与えられたとき、
整数 $n > 0$、$k \geq 0$ が存在し、$\mathcal{F}$ は
$k$ 個の $\mathcal{L}^{\otimes - n}$ のコピーの直和の商である。
\item
\label{properties-item-c-i-q}
(\ref{properties-item-c-q}) と同じ性質が、$\mathcal{F}$ を $X$ 上のすべての有限型イデアル層に
わたって動かしたときに成り立つ。
\end{enumerate}
\end{proposition}

\begin{proof}
補題 \ref{properties-lemma-ample-immersion-into-proj} は
(\ref{properties-item-ample}) $\Rightarrow$ (\ref{properties-item-immersion}) である。
補題 \ref{properties-lemma-ample-power-ample} と \ref{properties-lemma-open-in-proj-ample} は含意
(\ref{properties-item-ample}) $\Leftarrow$ (\ref{properties-item-immersion}) を与える。
含意 (\ref{properties-item-immersion}) $\Rightarrow$
(\ref{properties-item-s-affine-basis}) $\Rightarrow$ (\ref{properties-item-s-basis}) は、
構成の章の節 \ref{constructions-section-proj} から明らかである。
補題 \ref{properties-lemma-affine-s-opens} は
(\ref{properties-item-s-basis}) $\Rightarrow$ (\ref{properties-item-ample}) である。
従って最初の 4 条件はすべて同値である。

\medskip\noindent
同値な条件 (1) -- (4) を仮定する。特に $X$ は分離的であることに注意する
（分離スキーム $\text{Proj}(S)$ の開部分スキームだからである）。
$\mathcal{F}$ を $X$ 上の準連接層とする。
斉次元 $s \in S_{+}$ を、$X_s$ がアフィンとなるように選ぶ。
任意の切断 $m \in \Gamma(X_s, \mathcal{F})$ は、上の
(\ref{properties-item-qc-gg}) に表示された写像の一つの像に属すると主張する。
このアフィン開集合 $X_s$ は $X$ を被覆するので、これは
(\ref{properties-item-qc-gg}) を導く。実際、補題 \ref{properties-lemma-invert-s-sections} により、
$m$ を $m' \otimes s^{-n}$ の像として書ける。ここで、ある
$n \geq 1$ と
$m' \in \Gamma(X, \mathcal{F} \otimes \mathcal{L}^{\otimes n})$ を取っている。
これで主張が証明された。

\medskip\noindent
(\ref{properties-item-qc-gg}) $\Rightarrow$ (\ref{properties-item-qc-i-gg}) は明らかである。
(\ref{properties-item-qc-i-gg}) を仮定して、$\mathcal{L}$ が豊富であることを示そう。
$x \in X$ を選ぶ。アフィン開集合 $U \subset X$ を取り、これは $x$ を含むとする。
$Z = X \setminus U$ とおく。スキームの章の節
\ref{schemes-section-reduced} により、$Z$ を被約閉部分スキームとみなせる。
$\mathcal{I} \subset \mathcal{O}_X$ を閉部分スキーム $Z$ に対応する準連接イデアル層とする。
仮定 (\ref{properties-item-qc-i-gg}) により、$n \geq 1$ と切断
$s \in \Gamma(X, \mathcal{I} \otimes \mathcal{L}^{\otimes n})$ が存在し、
$s$ は $x$ において消えない（より正確には
$s \not \in \mathfrak m_x \mathcal{I}_x \otimes \mathcal{L}_x^{\otimes n}$）。
$s$ を $\mathcal{L}^{\otimes n}$ の切断とみなせる。
この切断は明らかに $Z$ に沿って消えるので $X_s \subset U$ である。
従って補題 \ref{properties-lemma-affine-cap-s-open} により $X_s$ はアフィンである。
これで $\mathcal{L}$ が豊富であることが証明された。
ここまでで (1) -- (6) が同値であることを証明した。

\medskip\noindent
同値な条件 (1) -- (6) を仮定する。以下では、大域生成される二つの加群層の
テンソル積も大域生成されるという事実を断りなく用いる
（加群の章の補題 \ref{modules-lemma-tensor-product-globally-generated} を参照）。
(1) により、元 $s_i \in S_{d_i}$（ここで $d_i \geq 1$）を、
$X = \bigcup_{i = 1, \ldots, n} X_{s_i}$ となるように取れる。
$d = d_1\ldots d_n$ とおく。このとき $\mathcal{L}^{\otimes d}$ は
$$
s_1^{d/d_1}, \ldots, s_n^{d/d_n}.
$$
により大域生成される。従って $\mathcal{L}^{\otimes j}$ が大域生成ならば、
$\mathcal{L}^{\otimes j + dn}$ もすべての $n \geq 0$ に対して大域生成である。
$j \in \{0, \ldots, d - 1\}$ を固定する。任意の点 $x \in X$ に対して、
$n \geq 1$ と大域切断 $s$ が存在し、これは $\mathcal{L}^{j + dn}$ の切断で、
$x$ において消えない。これは (\ref{properties-item-qc-gg}) を
$\mathcal{F} = \mathcal{L}^{\otimes j}$ と豊富な可逆層
$\mathcal{L}^{\otimes d}$ に適用して従う。$X$ は準コンパクトなので、
整数 $n_i$ と大域切断 $s_i$ の有限リストを選べる。ここで後者は
$\mathcal{L}^{\otimes j + dn_i}$ の切断であり、これらは $X$ のどの点でも
同時には消えない。$\mathcal{L}^{\otimes d}$ は大域生成なので、これは
$\mathcal{L}^{\otimes j + dn}$ が大域生成であることを意味する。ここで
$n = \max\{n_i\}$ である。これを法 $d$ の各合同類について証明したので、
$n_0 = n_0(\mathcal{L})$ が存在し、$\mathcal{L}^{\otimes n}$ はすべての
$n \geq n_0$ に対して大域生成である。この時点で、$\mathcal{F}$ が
大域生成ならば、$\mathcal{F} \otimes \mathcal{L}^{\otimes n}$ もすべての
$n \geq n_0$ に対して
大域生成であることが分かる。

\medskip\noindent
引き続き同値な条件 (1) -- (6) を仮定する。
$\mathcal{F}$ を有限型準連接 $\mathcal{O}_X$-加群層とする。
$\mathcal{F}_n \subset \mathcal{F}$ を (\ref{properties-item-qc-gg}) の標準写像の像とする。
構成により $\mathcal{F}_n \otimes \mathcal{L}^{\otimes n}$ は大域生成である。
(\ref{properties-item-qc-gg}) により、$\mathcal{F}$ は部分層 $\mathcal{F}_n$、
$n \geq 1$ の和である。加群の章の補題
\ref{modules-lemma-finite-type-quasi-compact-colimit} により、
$\mathcal{F} = \sum_{n = 1, \ldots, N} \mathcal{F}_n$ であるような
$N \geq 1$ が存在する。従って
$\mathcal{F} \otimes \mathcal{L}^{\otimes n}$ は
$n \geq N + n_0(\mathcal{L})$ ならば大域生成である。ここで
$n_0(\mathcal{L})$ は上のものである。よって (1) -- (6) は
(\ref{properties-item-c-gg}) を含意する。

\medskip\noindent
(\ref{properties-item-c-gg}) を仮定する。$\mathcal{F}$ を有限型準連接
$\mathcal{O}_X$-加群層とする。(\ref{properties-item-c-gg}) により、整数 $n \geq 1$ が存在し、
標準写像
$$
\Gamma(X, \mathcal{F} \otimes_{\mathcal{O}_X} \mathcal{L}^{\otimes n})
\otimes_{\mathbf{Z}} \mathcal{L}^{\otimes -n}
\longrightarrow
\mathcal{F}
$$
は全射である。$I$ を
$\Gamma(X, \mathcal{F} \otimes_{\mathcal{O}_X} \mathcal{L}^{\otimes n})$ の
有限部分集合全体の集合とし、包含関係で半順序を入れる。このとき $I$ は有向半順序集合である。
$i = \{s_1, \ldots, s_{r(i)}\}$ に対して $\mathcal{F}_i \subset \mathcal{F}$ を写像
$$
\bigoplus\nolimits_{j = 1, \ldots, r(i)} \mathcal{L}^{\otimes -n}
\longrightarrow
\mathcal{F}
$$
の像とする。この写像は $s_j$ による乗法を第 $j$ 因子上で行う。
上の全射性から $\mathcal{F} = \colim_{i \in I} \mathcal{F}_i$ が従う。
従って加群の章の補題 \ref{modules-lemma-finite-type-quasi-compact-colimit} を適用でき、
$\mathcal{F} = \mathcal{F}_i$ と結論でき、これはある $i$ に対して成り立つ。
これで (\ref{properties-item-c-q}) が証明された。言い換えると、
(\ref{properties-item-c-gg}) $\Rightarrow$ (\ref{properties-item-c-q}) である。

\medskip\noindent
含意 (\ref{properties-item-c-q}) $\Rightarrow$ (\ref{properties-item-c-i-q}) は自明である。

\medskip\noindent
最後に (\ref{properties-item-c-i-q}) を仮定する。
$\mathcal{I} \subset \mathcal{O}_X$ を準連接イデアル層とする。
補題 \ref{properties-lemma-quasi-coherent-colimit-finite-type} により
（ここで $X$ が準分離的であるという条件を用いる）、
$\mathcal{I} = \colim_\alpha I_\alpha$ と書け、各 $I_\alpha$ は有限型準連接である。
仮定により各 $I_\alpha$ は
$\mathcal{L}$ の負のテンソル冪の商なので、$\mathcal{I}$ についても同じ結論を得る
（ただし冪の有限性や有界性はもちろん失われる）。従って
(\ref{properties-item-c-i-q}) は (\ref{properties-item-qc-i-gg}) を含意する。
これで命題の証明は終わる。
\end{proof}

\begin{lemma}
\label{properties-lemma-ample-on-locally-closed}
$X$ をスキームとする。$\mathcal{L}$ を豊富な可逆 $\mathcal{O}_X$-加群とする。
$i : X' \to X$ をスキームの射とする。次の条件の少なくとも一つが成り立つと仮定する。
\begin{enumerate}
\item $i$ は準コンパクトな埋め込みである。
\item $X'$ は準コンパクトであり、$i$ は埋め込みである。
\item $i$ は準コンパクトであり、$X'$ と $i(X')$ の間の同相写像を誘導する。
\item $X'$ は準コンパクトであり、$i$ は $X'$ と $i(X')$ の間の同相写像を誘導する。
\end{enumerate}
このとき $i^*\mathcal{L}$ は $X'$ 上で豊富である。
\end{lemma}

\begin{proof}
場合 (1) と (3) では、スキーム $X'$ は準コンパクトである。実際、定義
\ref{properties-definition-ample} により $X$ が準コンパクトである。従って (2) と (4) を証明すれば十分である。
(2) は (4) の特別な場合なので、(4) を証明すれば十分である。

\medskip\noindent
条件 (4) を仮定する。$s \in \Gamma(X, \mathcal{L}^{\otimes d})$ に対して、
$s' = i^*s$ を $s$ の $X'$ への引き戻しとする。
$s'$ は $(i^*\mathcal{L})^{\otimes d}$ の切断であることに注意する。
命題 \ref{properties-proposition-characterize-ample} により、開集合 $X_s$ は、
$s \in \Gamma(X, \mathcal{L}^{\otimes d})$ にわたって
$X$ の位相の基底をなす。加群の章の注意
\ref{modules-remark-pullback-s-open} により $X'_{s'} = i^{-1}(X_s)$ であり、
$X' \to i(X')$ は同相写像なので、開集合 $X'_{s'}$ は $X'$ の位相の基底をなす。
従って命題 \ref{properties-proposition-characterize-ample} により $i^*\mathcal{L}$ は豊富である。
\end{proof}

\begin{lemma}
\label{properties-lemma-ample-on-product}
$S$ を準分離スキームとする。$X$、$Y$ を $S$ 上のスキームとする。
$\mathcal{L}$ を豊富な可逆 $\mathcal{O}_X$-加群とし、
$\mathcal{N}$ を豊富な可逆 $\mathcal{O}_Y$-加群とする。このとき
$\mathcal{M} = \text{pr}_1^*\mathcal{L}
\otimes_{\mathcal{O}_{X \times_S Y}} \text{pr}_2^*\mathcal{N}$ は
$X \times_S Y$ 上の豊富な可逆層である。
\end{lemma}

\begin{proof}
射 $i : X \times_S Y \to X \times Y$ は準コンパクトな埋め込みである。
スキームの章の補題 \ref{schemes-lemma-fibre-product-after-map} を参照せよ。
一方、$\mathcal{M}$ は $i$ による、$X \times Y$ 上の対応する可逆加群の引き戻しである。
補題 \ref{properties-lemma-ample-on-locally-closed} により、$X \times Y$ の場合を証明すれば十分である。
$\mathcal{M}$ について、$X \times Y$ 上で定義
\ref{properties-definition-ample} の (1) と (2) を確認する。

\medskip\noindent
$X$ と $Y$ は準コンパクトなので $X \times Y$ も準コンパクトである。
$z \in X \times Y$ を点とする。その射影を $x \in X$ および $y \in Y$ とする。
$n > 0$ と $s \in \Gamma(X, \mathcal{L}^{\otimes n})$ を、
$X_s$ が $x$ のアフィン開近傍となるように選ぶ。
$m > 0$ と $t \in \Gamma(Y, \mathcal{N}^{\otimes m})$ を、
$Y_t$ が $y$ のアフィン開近傍となるように選ぶ。
このとき $r = \text{pr}_1^*s \otimes \text{pr}_2^*t$ は $\mathcal{M}$ の切断であり、
$(X \times Y)_r = X_s \times Y_t$ である。これは $z$ のアフィン開近傍であり、
証明は完了する。
\end{proof}

\section{アフィンスキームと準アフィンスキーム}
\label{properties-section-affine-quasi-affine}

\begin{lemma}
\label{properties-lemma-quasi-affine-O-ample}
$X$ をスキームとする。
このとき、$X$ が準アフィンであることと
$\mathcal{O}_X$ が豊富であることは同値である。
\end{lemma}

\begin{proof}
$X$ が準アフィンであるとする。$A = \Gamma(X, \mathcal{O}_X)$ とおく。
補題 \ref{properties-lemma-quasi-affine} の開埋め込み
$$
j : X \longrightarrow \Spec(A)
$$
を考える。
$\Spec(A) = \text{Proj}(A[T])$ であることに注意せよ。構成の章の例
\ref{constructions-example-trivial-proj} を参照せよ。
従って補題 \ref{properties-lemma-open-in-proj-ample} を適用し、
$\mathcal{O}_X$ が豊富であることが得られる。

\medskip\noindent
$\mathcal{O}_X$ が豊富であるとする。
次数付き環として
$\Gamma_*(X, \mathcal{O}_X) \cong A[T]$
であることに注意せよ。従って、上で見たように
$\Spec(A) = \text{Proj}(A[T])$ が任意の環 $A$ に対して成り立つことを考慮すれば、
結果は補題 \ref{properties-lemma-ample-immersion-into-proj} および
\ref{properties-lemma-quasi-affine} から従う。
\end{proof}

\begin{lemma}
\label{properties-lemma-quasi-affine-locally-closed}
$X$ を準アフィンスキームとする。任意の準コンパクトな埋め込み
$i : X' \to X$ に対して、スキーム $X'$ は準アフィンである。
\end{lemma}

\begin{proof}
これは豊富な可逆層に関する素材を用いずに直接証明できる。
読者にはこれをナプキンの上で証明してみることを勧める。
$X$ は準アフィンなので、補題 \ref{properties-lemma-quasi-affine-O-ample} により
$\mathcal{O}_X$ は豊富である。
すると補題 \ref{properties-lemma-ample-on-locally-closed} により
$\mathcal{O}_{X'}$ は豊富である。従って補題
\ref{properties-lemma-quasi-affine-O-ample} により $X'$ は準アフィンである。
\end{proof}

\begin{lemma}
\label{properties-lemma-characterize-affine}
$X$ をスキームとする。有限個の元
$f_1, \ldots, f_n \in \Gamma(X, \mathcal{O}_X)$ が存在して、
\begin{enumerate}
\item 各 $X_{f_i}$ が $X$ のアフィン開集合であり、
\item $f_1, \ldots, f_n$ が $\Gamma(X, \mathcal{O}_X)$ において生成するイデアルが
単位イデアルに等しい
\end{enumerate}
と仮定する。このとき $X$ はアフィンである。
\end{lemma}

\begin{proof}
補題の通りの $f_1, \ldots, f_n$ があるとする。
$1 = \sum g_i f_i$ と、ある $g_j \in \Gamma(X, \mathcal{O}_X)$ を用いて書けるので、
$X = \bigcup X_{f_i}$ であることは明らかである。
（すべての $f_i$ が一つの点で消えることはできない。）
各 $X_{f_i}$ はアフィンであるため準コンパクトなので、
$X$ は準コンパクトである。従って上の補題
\ref{properties-lemma-quasi-affine-O-ample} により $X$ は準アフィンである。
開埋め込み
$$
j : X \to \Spec(\Gamma(X, \mathcal{O}_X)),
$$
を考える。補題 \ref{properties-lemma-quasi-affine} を参照せよ。
右辺の標準開集合 $D(f_i)$ の逆像は左辺の $X_{f_i}$ に等しく、
射 $j$ は同型
$X_{f_i} \cong D(f_i)$ を誘導する。補題
\ref{properties-lemma-invert-f-affine} を参照せよ。
$f_i$ たちは単位イデアルを生成するから
$\Spec(\Gamma(X, \mathcal{O}_X))
= \bigcup_{i = 1, \ldots, n} D(f_i)$ である。従って $j$ は同型である。
\end{proof}

\section{準連接層と豊富な可逆層}
\label{properties-section-ample-quasi-coherent}

\noindent
この節の主題は次の通りである。豊富な可逆層があれば、
すべての準連接層は次数付き加群から得られる。

\begin{situation}
\label{properties-situation-ample}
$X$ をスキームとする。
$\mathcal{L}$ を $X$ 上の豊富な可逆層とする。
$S = \Gamma_*(X, \mathcal{L})$ を次数付き環とみなす。
$Y = \text{Proj}(S)$ とおく。
$f : X \to Y$ を補題 \ref{properties-lemma-map-into-proj} の標準射とする。
これには $\mathbf{Z}$-次数付き $\mathcal{O}_X$-代数の射
$\bigoplus f^*\mathcal{O}_Y(n) \to \bigoplus \mathcal{L}^{\otimes n}$ が付随する。
\end{situation}

\noindent
次の補題は実際にはその次の補題の特別な場合であるが、
先にその成立を指摘するのがよいと思われる。

\begin{lemma}
\label{properties-lemma-ample-gcd-is-one}
状況 \ref{properties-situation-ample} において、
標準射 $f : X \to Y$ は $X$ を開部分スキーム
$W = W_1 \subset Y$ の中に写す。ここで $\mathcal{O}_Y(1)$ は可逆であり、
すべての乗法射
$\mathcal{O}_Y(n) \otimes_{\mathcal{O}_Y} \mathcal{O}_Y(m) \to
\mathcal{O}_Y(n + m)$
は同型である（構成の章の補題
\ref{constructions-lemma-where-invertible} を参照せよ）。
さらに、射 $f^*\mathcal{O}_Y(n) \to \mathcal{L}^{\otimes n}$ はすべて同型である。
\end{lemma}

\begin{proof}
命題 \ref{properties-proposition-characterize-ample} により、ある整数
$n_0$ が存在して、$\mathcal{L}^{\otimes n}$ はすべての
$n \geq n_0$ に対して大域生成である。点 $x \in X$ をとる。
上記により、$a \in S_{n_0}$ および $b \in S_{n_0 + 1}$ を見つけられ、
$a$ と $b$ は $x$ で消えない。
従って
$f(x) \in D_{+}(a) \cap D_{+}(b) = D_{+}(ab)$ である。
構成の章の補題 \ref{constructions-lemma-where-invertible} により、
望む通り $f(x) \in W_1$ である。射 $f$ の構成で用いた
構成の章の補題 \ref{constructions-lemma-invertible-map-into-proj} により、射
$f^*\mathcal{O}_Y(n_0) \to \mathcal{L}^{\otimes n_0}$ および
$f^*\mathcal{O}_Y(n_0 + 1) \to \mathcal{L}^{\otimes n_0 + 1}$
は $x$ の近傍で同型である。代数構造との両立性および
$f$ が $W$ の中に写すという事実により、すべての射
$f^*\mathcal{O}_Y(n) \to \mathcal{L}^{\otimes n}$ は $x$ の近傍で同型であると結論される。
従って証明できた。
\end{proof}

\noindent
局所環付き空間 $X$、可逆層 $\mathcal{L}$、および
$\mathcal{O}_X$-加群 $\mathcal{F}$ が与えられると、次数付き
$\Gamma_*(X, \mathcal{L})$-加群
$$
\Gamma_*(X, \mathcal{L}, \mathcal{F}) =
\bigoplus\nolimits_{n \in \mathbf{Z}}
\Gamma(X, \mathcal{F} \otimes_{\mathcal{O}_X} \mathcal{L}^{\otimes n}).
$$
があることを、加群の章の定義 \ref{modules-definition-gamma-star} から思い出そう。
次の補題は、状況 \ref{properties-situation-ample} において、この次数付き加群から
準連接 $\mathcal{O}_X$-加群 $\mathcal{F}$ を復元できることを述べる。
$X = \mathbf{P}^n_R$ という特別な場合にこの補題を証明する、
構成の章の補題 \ref{constructions-lemma-quasi-coherent-projective-space} も参照せよ。

\begin{lemma}
\label{properties-lemma-ample-quasi-coherent}
状況 \ref{properties-situation-ample} において、
$\mathcal{F}$ を $X$ 上の準連接層とする。
$M = \Gamma_*(X, \mathcal{L}, \mathcal{F})$ を次数付き $S$-加群とみなす。
$\mathcal{F}$ に関して関手的な同型
$$
f^*\widetilde{M} \longrightarrow \mathcal{F}
$$
が存在し、
$M_0 \to \Gamma(\text{Proj}(S), \widetilde{M}) \to \Gamma(X, \mathcal{F})$
は恒等写像である。
\end{lemma}

\begin{proof}
$s \in S_{+}$ を、$X_s$ が $X$ のアフィン開集合となる斉次元とする。
$\widetilde{M}|_{D_{+}(s)}$ が $S_{(s)}$-加群 $M_{(s)}$ に対応することを思い出そう。
構成の章の補題 \ref{constructions-lemma-proj-sheaves} を参照せよ。
$f^{-1}(D_{+}(s)) = X_s$ であることを思い出そう。
$X$ は豊富な可逆層を持つので、節 \ref{properties-section-ample} により
準コンパクトかつ準分離的である。
補題 \ref{properties-lemma-invert-s-sections} により標準同型
$M_{(s)} = \Gamma_*(X, \mathcal{L}, \mathcal{F})_{(s)} \to
\Gamma(X_s, \mathcal{F})$ がある。
$\mathcal{F}$ は準連接なので、これは標準同型
$$
f^*\widetilde{M}|_{X_s} \to \mathcal{F}|_{X_s}
$$
を与える。$\mathcal{L}$ は $X$ 上豊富なので、$X$ は $X_s$ の形のアフィン開集合によって覆われる。
従って、表示した射が重なり上で貼り合わされることを証明すれば十分である。
その証明は省略する。
\end{proof}

\begin{remark}
\label{properties-remark-neurotic}
補題 \ref{properties-lemma-ample-quasi-coherent} の仮定と記号の下で考える。
その補題の表示した射を $\theta_\mathcal{F}$ と記す。
補題 \ref{properties-lemma-ample-gcd-is-one} の同型
$f^*\mathcal{O}_Y(n) \to \mathcal{L}^{\otimes n}$ は、まさに
$\theta_{\mathcal{L}^{\otimes n}}$ であることに注意せよ。
乗法射
$$
\widetilde{M} \otimes_{\mathcal{O}_Y} \mathcal{O}_Y(n)
\longrightarrow
\widetilde{M(n)}
$$
を考える。構成の章の等式
(\ref{constructions-equation-multiply-more-generally}) を参照せよ。
これを $X$ へ引き戻し、
$$
\xymatrix{
f^*\widetilde{M} \otimes_{\mathcal{O}_X} f^*\mathcal{O}_Y(n)
\ar[r]
\ar[d]_{\theta_\mathcal{F} \otimes \theta_{\mathcal{L}^{\otimes n}}}
&
f^*\widetilde{M(n)}
\ar[d]^{\theta_{\mathcal{F} \otimes \mathcal{L}^{\otimes n}}}
\\
\mathcal{F} \otimes \mathcal{L}^{\otimes n} \ar[r]^{\text{id}} &
\mathcal{F} \otimes \mathcal{L}^{\otimes n}
}
$$
を考える。ここで明らかな同一視
$M(n) = \Gamma_*(X, \mathcal{L}, \mathcal{F} \otimes \mathcal{L}^{\otimes n})$ を用いた。
この図式は可換である。証明は省略する。
\end{remark}

\noindent
次の補題は補題 \ref{properties-lemma-ample-quasi-coherent} から導けるはずであり
（または逆にそこから前者を導けるはずである）が、単に証明を繰り返す方が簡単に思われる。

\begin{lemma}
\label{properties-lemma-proj-quasi-coherent}
$S$ を次数付き環とし、$X = \text{Proj}(S)$ が準コンパクトであるとする。
$\mathcal{F}$ を準連接 $\mathcal{O}_X$-加群とする。
$M = \bigoplus_{n \in \mathbf{Z}} \Gamma(X, \mathcal{F}(n))$ を次数付き
$S$-加群とみなす。構成の章の節
\ref{constructions-section-invertible-on-proj} を参照せよ。
構成の章の補題
\ref{constructions-lemma-comparison-proj-quasi-coherent} の射
$$
\widetilde{M} \longrightarrow \mathcal{F}
$$
は同型である。$X$ が標準開集合 $D_+(f)$ で覆われ、
そこで $f$ の次数が $1$ ならば、誘導される射
$M_n \to \Gamma(X, \mathcal{F}(n))$ は恒等写像である。
\end{lemma}

\begin{proof}
$X$ は準コンパクトなので、正の次数の斉次元
$f_1, \ldots, f_n \in S$ であって
$X = D_+(f_1) \cup \ldots \cup D_+(f_n)$ となるものを見つけられる。
$d$ を $f_1, \ldots, f_n$ の次数の最小公倍数とする。
$f_i$ をある冪で置き換えることで、各 $f_i$ の次数が $d$ であると仮定してよい。
すると $\mathcal{L} = \mathcal{O}_X(d)$ は可逆であり、乗法射
$\mathcal{O}_X(ad) \otimes \mathcal{O}_X(bd) \to \mathcal{O}_X((a + b)d)$
は同型であり、各 $f_i$ は大域切断 $s_i$ を定める。
これは $\mathcal{L}$ の切断であり、$X_{s_i} = D_+(f_i)$ を満たす。構成の章の補題
\ref{constructions-lemma-where-invertible} および
\ref{constructions-lemma-principal-open} を参照せよ。
従って $\Gamma(X, \mathcal{F}(ad)) =
\Gamma(X, \mathcal{F} \otimes \mathcal{L}^{\otimes a})$ である。
$\widetilde{M}|_{D_{+}(f_i)}$ が $S_{(f_i)}$-加群 $M_{(f_i)}$ に対応することを思い出そう。
構成の章の補題 \ref{constructions-lemma-proj-sheaves} を参照せよ。
$f_i$ の次数は $d$ なので、$M_{(f_i)}$ の同型類は、
$M$ の斉次直和成分のうち次数が $d$ で割り切れるもののみに依存する。より正確には、
$M_{(f_i)}$ の同型類は次数付き $\Gamma_*(X, \mathcal{L})$-加群
$\Gamma_*(X, \mathcal{L}, \mathcal{F})$ と、$s_i$、すなわち $f_i$ の
$\Gamma_*(X, \mathcal{L})$ における像のみに依存する。
仮定によりスキーム $X$ は準コンパクトであり、構成の章の補題
\ref{constructions-lemma-proj-separated} により分離的である。
補題 \ref{properties-lemma-invert-s-sections} により標準同型
$$
M_{(f_i)} = \Gamma_*(X, \mathcal{L}, \mathcal{F})_{(s_i)} \to
\Gamma(X_{s_i}, \mathcal{F}).
$$
がある。そこで、構成の章の補題
\ref{constructions-lemma-comparison-proj-quasi-coherent}
の射の構成は、その射が $D_+(f_i)$ 上で同型であることを示す。
$X$ はこれらの開集合で覆われるので、射は同型である。
最後の主張の証明は省略する。
\end{proof}

\section{適切なアフィン開集合を見つけること}
\label{properties-section-finding-affine-opens}

\noindent
この節では、より一般的な状況からそうでない状況までにおける、
アフィン開集合の存在に関するいくつかの結果をまとめる。

\begin{lemma}
\label{properties-lemma-maximal-points-affine}
$X$ を準分離スキームとする。
$Z_1, \ldots, Z_n$ を $X$ の互いに異なる既約成分とする。
位相空間の章の節 \ref{topology-section-irreducible-components} を参照せよ。
$\eta_i \in Z_i$ をそれらの一般点とする。
スキームの章の補題 \ref{schemes-lemma-scheme-sober} を参照せよ。
アフィン開近傍 $\eta_i \in U_i$ であって、$U_i \cap U_j = \emptyset$ がすべての $i \not = j$ に対して成り立つものが存在する。
特に、$U = U_1 \cup \ldots \cup U_n$ はすべての点
$\eta_1, \ldots, \eta_n$ を含むアフィン開集合である。
\end{lemma}

\begin{proof}
$V_i$ を、$\eta_i$ を含み、閉集合
$Z_1 \cup \ldots \hat Z_i \ldots \cup Z_n$ と交わらない任意のアフィン開集合とする。
$X$ は準分離的なので、各 $i$ に対して和集合
$W_i = \bigcup_{j, j \not = i} V_i \cap V_j$ は $V_i$ の準コンパクト開集合であり、$\eta_i$ を含まない。
代数学の章の補題 \ref{algebra-lemma-standard-open-containing-maximal-point} により、$U_i \subset V_i$ なる開近傍を、$\eta_i$ を含み $W_i$ と交わらないように見つけられる。
最後に、$U$ は環 $R_1 \times \ldots \times R_n$ のスペクトルなのでアフィンである。
ここで $R_i = \mathcal{O}_X(U_i)$ である。スキームの章の補題
\ref{schemes-lemma-disjoint-union-affines} を参照せよ。
\end{proof}

\begin{remark}
\label{properties-remark-maximal-points-affine}
$X$ が準分離的でないなら、上の補題 \ref{properties-lemma-maximal-points-affine} は偽である。
次が例である。
$R = \mathbf{Q}[x, y_1, y_2, \ldots]/((x-i)y_i)$ とおく。
極小素イデアル $\mathfrak p = (y_1, y_2, \ldots)$ を考える。
$R$ の $\Spec(R)$ の二つのコピーを（準コンパクトでない）開集合
$\Spec(R) \setminus V(\mathfrak p)$ に沿って貼り合わせ、スキーム $X$ を得る。
貼り合わせについてはスキームの章の例
\ref{schemes-example-affine-space-zero-doubled} を参照せよ。
このとき、$X$ の $\mathfrak p$ に対応する二つの極大点は、
共通のアフィン開集合に含まれない。その理由は、
$\Spec(R)$ の、$\mathfrak p$ を含む任意の開集合が、無限個の「直線」
$x = i$, $y_j = 0$, $j \not = i$、すなわちパラメータ $y_i$ を持つものを含むことである。
詳細は省略する。
\end{remark}

\noindent
上の例にもかかわらず、既約閉部分集合の「ほとんどの」有限集合に対しては、
少なくとも $X$ が準コンパクトなら、上の補題 \ref{properties-lemma-maximal-points-affine} を適用できる。
これは $X$ が分離的な稠密開集合を含むからである。

\begin{lemma}
\label{properties-lemma-quasi-compact-dense-open-separated}
$X$ を準コンパクトスキームとする。
分離的な稠密開集合 $V \subset X$ が存在する。
\end{lemma}

\begin{proof}
$X = \bigcup_{i = 1, \ldots, n} U_i$ が $n$ 個のアフィン開部分スキームの和であるとする。
$n$ に関する帰納法で補題を証明する。$n = 1$ の場合は明らかである。
帰納法の仮定によって存在する分離的な稠密開部分スキーム
$V' \subset \bigcup_{i = 1, \ldots, n - 1} U_i$ をとる。
$$
V = V' \amalg (U_n \setminus \overline{V'}).
$$
を考える。$V$ が $X$ の分離的な稠密開部分スキームであることは明らかである。
\end{proof}

\noindent
$X$ が準コンパクトかつ準分離的であっても、
分離的かつ準コンパクトな稠密開集合は存在しないことが分かる。
例の章の補題
\ref{examples-lemma-no-dense-separated-quasi-compact-open-in-qcqs} を参照せよ。
次は上の補題 \ref{properties-lemma-maximal-points-affine} のわずかな精密化である。

\begin{lemma}
\label{properties-lemma-point-and-maximal-points-affine}
$X$ を準分離スキームとする。$Z_1, \ldots, Z_n$ を
$X$ の互いに異なる既約成分とする。$\eta_i \in Z_i$ をそれらの一般点とする。
$x \in X$ を任意にとる。アフィン開集合 $U \subset X$ が存在し、$x$ とすべての $\eta_i$ を含む。
\end{lemma}

\begin{proof}
$x \in Z_1 \cap \ldots \cap Z_r$ かつ
$x \not \in Z_{r + 1}, \ldots, Z_n$ であるとする。するとアフィン開集合
$W \subset X$ を、$x \in W$ かつ $W \cap Z_i = \emptyset$ が $i = r + 1, \ldots, n$ に対して成り立つように選べる。
明らかに、$\eta_i \in W$ は $i = 1, \ldots, r$ に対して成り立つ。
補題 \ref{properties-lemma-maximal-points-affine} により、互いに交わらないアフィン開集合 $U_i \subset X$ を、$\eta_i \in U_i$ が $i = r + 1, \ldots, n$ に対して成り立つように選べる。
$X$ は準分離的なので開集合 $W \cap U_i$ は準コンパクトであり、$\eta_i$ を含まない。これは $i = r + 1, \ldots, n$ に対して成り立つ。従って代数学の章の補題
\ref{algebra-lemma-standard-open-containing-maximal-point} により、
$U_i$ を、$W \cap U_i = \emptyset$ となるように $i = r + 1, \ldots, n$ に対して縮小できる。
このとき和集合
$U = W \cup \bigcup_{i = r + 1, \ldots, n} U_i$ は交わりを持たず、従って
（スキームの章の補題 \ref{schemes-lemma-disjoint-union-affines} により）
求めるアフィン開集合である。
\end{proof}

\begin{lemma}
\label{properties-lemma-ample-finite-set-in-affine}
$X$ をスキームとする。次のいずれかを仮定する。
\begin{enumerate}
\item スキーム $X$ は準アフィンである。
\item スキーム $X$ はアフィンスキームの局所閉部分スキームと同型である。
\item $X$ 上に豊富な可逆層が存在する。
\item スキーム $X$ は、次数付き環の
$\text{Proj}(S)$ の局所閉部分スキームと同型である。ここで $S$ はその次数付き環である。
\end{enumerate}
このとき、任意の有限部分集合 $E \subset X$ に対して、
$U \subset X$ なるアフィン開集合であって $E \subset U$ を満たすものが存在する。
\end{lemma}

\begin{proof}
性質の章の定義 \ref{properties-definition-quasi-affine} により、
準アフィンスキームはアフィンスキームの準コンパクト開部分スキームである。
任意のアフィンスキーム $\Spec(R)$ は $\text{Proj}(R[X])$ と同型である。
ここで $R[X]$ は $\deg(X) = 1$ とおくことで次数付けられている。
命題 \ref{properties-proposition-characterize-ample} により、$X$ が豊富な可逆層を持つなら、
$X$ はある次数付き環の $\text{Proj}(S)$ の開部分スキームと同型である。ここで $S$ はその次数付き環である。
従って、場合 (4) で補題を証明すれば十分である。
（読者には場合 (2) を直接証明してみることを勧める。）

\medskip\noindent
従って、$X \subset \text{Proj}(S)$ が局所閉部分スキームであり、
$S$ がある次数付き環であると仮定する。$T = \overline{X} \setminus X$ とおく。
標準開集合 $D_{+}(f)$ が $\text{Proj}(S)$ の位相の基底をなすことを思い出そう。
$E$ は有限なので、有限個の斉次元 $f_i \in S_{+}$ を
$$
E \subset
D_{+}(f_1) \cup \ldots \cup D_{+}(f_n) \subset
\text{Proj}(S) \setminus T
$$
となるように選べる。$E = \{\mathfrak p_1, \ldots, \mathfrak p_m\}$ を $\text{Proj}(S)$ の部分集合として考える。
イデアル $I = (f_1, \ldots, f_n) \subset S$ を考える。
イデアル $I \not \subset \mathfrak p_j$ がすべての $j = 1, \ldots, m$ に対して成り立つので、
代数学の章の補題 \ref{algebra-lemma-graded-silly} から、斉次元
$f \in I$ であって、$f \not \in \mathfrak p_j$ がすべての $j = 1, \ldots, m$ に対して成り立つものが存在することが分かる。
すると $E \subset D_{+}(f) \subset
D_{+}(f_1) \cup \ldots \cup D_{+}(f_n)$ である。$D_{+}(f)$ は
$T$ と交わらないので、$X \cap D_{+}(f)$ はアフィンスキーム
$D_{+}(f)$ の閉部分スキームであり、従って望む通り $X$ のアフィン開集合である。
\end{proof}

\begin{lemma}
\label{properties-lemma-ample-finite-set-in-principal-affine}
$X$ をスキームとする。$\mathcal{L}$ を $X$ 上の豊富な可逆層とする。

$$
E \subset W \subset X
$$
とし、$E$ は有限、$W$ は $X$ の開集合とする。このとき、ある $n > 0$ と
$s \in \Gamma(X, \mathcal{L}^{\otimes n})$ が存在して、
$X_s$ はアフィンであり、$E \subset X_s \subset W$ となる。
\end{lemma}

\begin{proof}
読者は、補題 \ref{properties-lemma-ample-finite-set-in-affine} の証明を修正してこの補題を証明できるが、ここでは代わりにその補題から導く。
補題 \ref{properties-lemma-ample-finite-set-in-affine} により、$U \subset W$ となるアフィン開集合を選べ、$E \subset U$ とできる。
次の次数付き環を考える。
$S = \Gamma_*(X, \mathcal{L}) =
\bigoplus_{n \geq 0} \Gamma(X, \mathcal{L}^{\otimes n})$.

$x \in E$ ごとに、切断の次数付きイデアル $\mathfrak p_x \subset S$ をとる。これは $x$ で消える切断からなる。$\mathfrak p_x$ が素イデアルであることは明らかであり、また、$\mathcal{L}$ のある冪が大域生成されるので、$S_{+} \not \subset \mathfrak p_x$ であることも明らかである。
次数付きイデアル $I \subset S$ を、$X \setminus U$ のすべての点で消える切断からなるものとしてとる。$X_s$ が位相の基底をなすことから、$I \not \subset \mathfrak p_x$ がすべての $x \in E$ に対して成り立つ。
代数学の章の補題 \ref{algebra-lemma-graded-silly} の（次数付き）素イデアル回避により、斉次元の $s \in I$ であって、$s \not \in \mathfrak p_x$ がすべての $x \in E$ に対して成り立つものを見つけられる。
そのとき $E \subset X_s \subset U$ であり、$X_s$ は補題 \ref{properties-lemma-affine-cap-s-open} によりアフィンである。
\end{proof}

\begin{lemma}
\label{properties-lemma-quasi-affine-invertible-nonvanishing-section}
$X$ を準アフィンスキームとする。$\mathcal{L}$ を可逆 $\mathcal{O}_X$-加群とする。
$E \subset W \subset X$ とし、$E$ は有限、$W$ は開集合とする。このとき、ある $s \in \Gamma(X, \mathcal{L})$ が存在して、
$X_s$ はアフィンであり、$E \subset X_s \subset W$ となる。
\end{lemma}

\begin{proof}
この補題の証明は、代数学の章の補題 \ref{algebra-lemma-silly} の証明と多くの点で共通している。
$E = \{x_1, \ldots, x_n\}$ と書く。$E = W = \emptyset$ ならば $s = 0$ でよい。$W \not = \emptyset$ ならば、必要なら点を一つ加えることにより $E \not = \emptyset$ としてよい。したがって $n \geq 1$ と仮定できる。
$n$ に関する帰納法で補題を証明する。

\medskip\noindent
基底の場合: $n = 1$ 。$W$ を、$x_1$ を含む $W$ 内のアフィン開近傍で置き換え、$W$ がアフィンであると仮定できる。
補題 \ref{properties-lemma-quasi-affine-O-ample} と命題 \ref{properties-proposition-characterize-ample} を組み合わせると、任意の準連接 $\mathcal{O}_X$-加群は大域生成される。
したがって、大域切断 $s$ であって $\mathcal{L}$ のものであり、$x_1$ で消えないものが存在する。
一方、被約誘導閉部分スキーム $Z \subset X$ を $X \setminus W$ 上にとる。準連接理想層 $\mathcal{I}$ に $Z$ のものとして大域生成を適用すると、$f$ を $\mathcal{I}$ の大域切断として $x_1$ で消えないものが見つかる。そこで $s' = fs$ とおく。これは $\mathcal{L}$ の大域切断であり、$x_1$ で消えず、$X_{s'} \subset W$ を満たす。したがって、$X_{s'}$ は補題 \ref{properties-lemma-affine-cap-s-open} によりアフィンである。

\medskip\noindent
$n > 1$ の帰納段階。もし特殊化 $x_i \leadsto x_j$ が $i \not = j$ であるとき存在すれば、
$\{x_1, \ldots, x_n\} \setminus \{x_i\}$ について補題を証明すれば十分であり、帰納法によって完了する。従って $x_i$ たちの間に特殊化はないと仮定できる。
補題 \ref{properties-lemma-ample-finite-set-in-affine} または補題 \ref{properties-lemma-ample-finite-set-in-principal-affine} のいずれかにより、$W$ がアフィンであると仮定できる。
帰納法により、大域切断 $s$ であって $\mathcal{L}$ のものであり、$X_s \subset W$ がアフィンで、$x_1, \ldots, x_{n - 1}$ を含むものを見つけられる。$x_n \in X_s$ ならば完了である。
$s$ が $x_n$ で零であると仮定する。$n = 1$ の場合により、大域切断 $s'$ を見つけられ、これは $\mathcal{L}$ のものであり、
$\{x_n\} \subset X_{s'} \subset
W \setminus \overline{\{x_1, \ldots, x_{n - 1}\}}$ を満たす。
ここで、$x_n$ が $x_1, \ldots, x_{n - 1}$ の特殊化ではないことを用いた。
そのとき $s + s'$ は、$\mathcal{L}$ の大域切断であり、$x_1, \ldots, x_n$ で消えず、$X_{s + s'} \subset W$ となる。従って先ほどと同様に結論できる。
\end{proof}

\begin{lemma}
\label{properties-lemma-ring-affine-open-injective-local-ring}
$X$ をスキーム、$x \in X$ を点とする。次の各場合に、$U \subset X$ なるアフィン開近傍であって $x$ を含み、標準写像
$\mathcal{O}_X(U) \to \mathcal{O}_{X, x}$ が単射であるものが存在する。
\begin{enumerate}
\item $X$ が整である。
\item $X$ が局所的に Noetherian である。
\item $X$ が被約であり、既約成分の個数が有限である。
\end{enumerate}
\end{lemma}

\begin{proof}
代数学への翻訳の後、これは代数学の章の補題 \ref{algebra-lemma-subring-of-local-ring} から従う。
\end{proof}

\begin{lemma}
\label{properties-lemma-standard-open-two-affines}
$U$, $V$ をアフィンスキームとし、$W \to U$ および $W \to V$ を開埋め込みとする。任意の $w \in W$ に対して、アフィン開近傍 $W' \subset W$ が存在し、$w$ を含み、$W'$ が $U$ と $V$ の両方の標準開集合に写る。
\end{lemma}

\begin{proof}
スキーム $X$ をとる。これは $U$ と $V$ を $W$ に沿って貼り合わせて得られる。スキームの章の補題 \ref{schemes-lemma-standard-open-two-affines} を適用する。
\end{proof}
